%============================================================
%  Bd. 18 - Rationale Zahlen
%============================================================

\documentclass[main.tex]{subfiles}

\ifSubfilesClassLoaded{
  \BandSetupStandalone{B18}
}{
}

\title{Bd. 18 - Rationale Zahlen}
\author{Martin Kunze}
\date{}
\setcounter{file}{18}

\hypersetup{
  pdftitle={Bd. 18 - Rationale Zahlen},
  pdfauthor={Martin Kunze}
}

% Lokale Fallbacks spiegeln die zentrale Zahlennotation und halten den Band
% auch bei einer isolierten Weiterverarbeitung lesbar.
\providecommand{\PositiveIntegers}{\ensuremath{\mathbb Z_{>0}}}
\providecommand{\RatPairCarrier}{\ensuremath{D_{\mathbb Q}}}
\providecommand{\RatPairRel}{\ensuremath{\mathrel{\sim_{\mathbb Q}}}}
\providecommand{\RatClass}[2]{\ensuremath{\left[#1,#2\right]_{\mathbb Q}}}
\providecommand{\RatEmbed}{\ensuremath{\iota_{\mathbb Q}}}
\providecommand{\RatOrderGraph}{\ensuremath{R^{\leq}_{\mathbb Q}}}
\providecommand{\RatRecipRep}[2]{%
  \ensuremath{\operatorname{rec}_{\mathbb Q}\!\left(#1,#2\right)}%
}
\providecommand{\RationalOrderedField}[1]{%
  \ensuremath{\mathsf{QOF}\!\left(#1\right)}%
}

\begin{document}

\title{Bd. 18 - Rationale Zahlen}
\author{Martin Kunze}
\date{}
\hypersetup{pageanchor=false}
\maketitle
\hypersetup{pageanchor=true}
\setcounter{tocdepth}{3}
\tableofcontents
\setcounter{file}{18}
\renewcommand{\FormulaBandID}{18}

\chapter{Einleitung}

Die rationalen Zahlen entstehen aus ganzzahligen Zählern und positiven
ganzzahligen Nennern. Ein Paar \((a,b)\) wird als formaler Bruch
\(a/b\) gelesen. Verschiedene Bruchpaare beschreiben dieselbe Zahl; daher
werden genau diejenigen Paare identifiziert, deren Kreuzprodukte
übereinstimmen. Die Äquivalenz- und Quotiententheorie aus Band~11 sowie die
ganzen Zahlen aus Band~17 liefern die dafür benötigten Grundlagen.

Wie bei der Konstruktion der ganzen Zahlen wird jede Operation erst nach dem
Nachweis ihrer Repräsentantenunabhängigkeit eingeführt. Positive Nenner
vermeiden ein zusätzliches Vorzeichen im Nenner und liefern zugleich eine
einfache Quotientenordnung. Abschließend werden die gewonnenen Gesetze als
archimedisch geordnete Körperstruktur zusammengefasst. Diese Zusammenfassung
ist eine metasprachliche Bündelung der zuvor einzeln bewiesenen Aussagen; sie
setzt keinen bereits eingeführten abstrakten Körperbegriff voraus.

\chapter{Konstruktion der rationalen Zahlen}

\section{Bruchpaare}

\FormulaDefDeltaK[Träger der Bruchpaare]{%
  \RatPairCarrier\coloneqq\mathbb Z\times\PositiveIntegers%
}{RationalPairCarrier}{%
  \DeltaRow{Mengen}{\RatPairCarrier\dsep\PositiveIntegers}%
  \DeltaRow{Neue Symbole}{\RatPairCarrier}%
}

\paragraph{Metasprachliche Motivation der Bruchpaarrelation.}
In einer bereits vorhandenen Bruchrechnung würden die formalen Brüche
\(a/b\) und \(c/d\) genau dann dieselbe Zahl bezeichnen, wenn nach dem
Erweitern auf den gemeinsamen Nenner \(bd\) ihre Zähler übereinstimmen;
dies führt auf die Kreuzproduktgleichung \(a\cdot d=c\cdot b\). Diese
Überlegung dient hier nur der Motivation. Eine Division ganzer Zahlen mit
rationalem Ergebnis steht an dieser Stelle noch nicht zur Verfügung und
darf daher nicht zur Definition herangezogen werden. Formal wird allein die
folgende Relation auf Bruchpaaren definiert; erst die anschließenden Beweise
von Reflexivität, Symmetrie und Transitivität rechtfertigen den Übergang zu
ihren Äquivalenzklassen.

\FormulaDefDeltaK[Äquivalenz von Bruchpaaren]{%
  a,c\in\mathbb Z\dsep b,d\in\PositiveIntegers\vdash
  (a,b)\RatPairRel(c,d)
  \coloneqq a\cdot d=c\cdot b%
}{RationalPairRelDef}{%
  \DeltaRow{Mengen}{a\dsep b\dsep c\dsep d}%
  \DeltaRow{Relationsoperatoren}{\RatPairRel}%
  \DeltaRow{Trägermenge}{\RatPairCarrier}%
  \DeltaRow{Neue Symbole}{\RatPairRel}%
}

\FormulaThmDeltaK[Charakterisierung der Bruchpaaräquivalenz]{%
  a,c\in\mathbb Z\dsep b,d\in\PositiveIntegers\vdash
  \bigl((a,b)\RatPairRel(c,d)
  \leftrightarrow a\cdot d=c\cdot b\bigr)%
}{RationalPairRelChar}{%
  \DeltaRow{Mengen}{a\dsep b\dsep c\dsep d}%
  \DeltaRow{Relationsoperatoren}{\RatPairRel}%
}
\begin{tabproof}
  \proofstep{1}{a,c\in\mathbb Z}{\rA}
  \proofstep{2}{b,d\in\PositiveIntegers}{\rA}
  \proofstep{1,2}{%
    (a,b)\RatPairRel(c,d)\leftrightarrow a\cdot d=c\cdot b}{%
    \FormulaRefAuto{RationalPairRelDef}[def]{1,2}}
\end{tabproof}

\FormulaThmDeltaK[Reflexivität auf Bruchpaaren]{%
  a\in\mathbb Z\dsep b\in\PositiveIntegers
  \vdash(a,b)\RatPairRel(a,b)%
}{RationalPairRelReflexive}{%
  \DeltaRow{Mengen}{a\dsep b}%
  \DeltaRow{Relationsoperatoren}{\RatPairRel}%
}
\begin{tabproof}
  \proofstep{1}{a\in\mathbb Z}{\rA}
  \proofstep{2}{b\in\PositiveIntegers}{\rA}
  \proofstep{}{a\cdot b=a\cdot b}{\rII}
  \proofstep{1,2}{(a,b)\RatPairRel(a,b)}{%
    \FormulaRefAuto{RationalPairRelChar}[thm]{1,2,3}}
\end{tabproof}

\FormulaThmDeltaK[Symmetrie auf Bruchpaaren]{%
  a,c\in\mathbb Z\dsep b,d\in\PositiveIntegers\dsep
  (a,b)\RatPairRel(c,d)
  \vdash(c,d)\RatPairRel(a,b)%
}{RationalPairRelSymmetric}{%
  \DeltaRow{Mengen}{a\dsep b\dsep c\dsep d}%
  \DeltaRow{Relationsoperatoren}{\RatPairRel}%
}
\begin{tabproof}
  \proofstep{1}{a,c\in\mathbb Z}{\rA}
  \proofstep{2}{b,d\in\PositiveIntegers}{\rA}
  \proofstep{3}{(a,b)\RatPairRel(c,d)}{\rA}
  \proofstep{1,2,3}{a\cdot d=c\cdot b}{%
    \FormulaRefAuto{RationalPairRelChar}[thm]{1,2,3}}
  \proofstep{1,2,3}{c\cdot b=a\cdot d}{%
    \FormulaRefAuto{a=b\vdash b=a}{4}}
  \proofstep{1,2,3}{(c,d)\RatPairRel(a,b)}{%
    \FormulaRefAuto{RationalPairRelChar}[thm]{1,2,5}}
\end{tabproof}

\FormulaThmDeltaK[Transitivität auf Bruchpaaren]{%
  \begin{aligned}[t]
    &a,c,e\in\mathbb Z\dsep b,d,f\in\PositiveIntegers\dsep
      (a,b)\RatPairRel(c,d)\dsep{}\\
    &(c,d)\RatPairRel(e,f)
      \vdash(a,b)\RatPairRel(e,f)
  \end{aligned}%
}{RationalPairRelTransitive}{%
  \DeltaRow{Mengen}{a\dsep b\dsep c\dsep d\dsep e\dsep f}%
  \DeltaRow{Relationsoperatoren}{\RatPairRel}%
}

\begin{tabproofwide}
  \proofstepwidestarbreakkeep[1]{a,c,e\in\mathbb Z}{\rA}
  \proofstepwidestarbreakkeep[2]{b,d,f\in\PositiveIntegers}{\rA}
  \proofstepwidestarbreakkeep[3]{(a,b)\RatPairRel(c,d)}{\rA}
  \proofstepwidestarbreakkeep[4]{(c,d)\RatPairRel(e,f)}{\rA}
  \proofstepwidestarbreakkeep[1,2,3]{a\cdot d=c\cdot b}{%
    \FormulaRefAuto{RationalPairRelChar}[thm]{1,2,3}}
  \proofstepwidestarbreakkeep[1,2,4]{c\cdot f=e\cdot d}{%
    \FormulaRefAuto{RationalPairRelChar}[thm]{1,2,4}}
  \proofstepwide[1,2,3]{(a\cdot f)\cdot d}{=}{%
    (a\cdot d)\cdot f}{%
    \FormulaRefAuto{IntMulAssociative}[thm],
    \FormulaRefAuto{IntMulCommutative}[thm]}
  \proofstepwide[1,2,3]{(a\cdot d)\cdot f}{=}{%
    (c\cdot b)\cdot f}{\rIE{5}}
  \proofstepwide[1,2,3,4]{(c\cdot b)\cdot f}{=}{%
    (e\cdot b)\cdot d}{%
    \FormulaRefAuto{IntMulAssociative}[thm],
    \FormulaRefAuto{IntMulCommutative}[thm],6}
  \proofstepwidestarbreakkeep[1,2,3,4]{(a\cdot f)\cdot d=(e\cdot b)\cdot d}{%
    \rChain{7,8,9}}
  \proofstepwidestar[1,2,3,4]{a \cdot f \in \mathbb Z}{%
      \FormulaRefAuto{IntMulClosure}[thm]{1}}
  \proofstepwidestar[1,2,3,4]{e \cdot b \in \mathbb Z}{%
      \FormulaRefAuto{IntMulClosure}[thm]{1}}
  \proofstepwidestarbreakkeep[1,2,3,4]{a\cdot f=e\cdot b}{%
    \FormulaRefAuto{IntegerPositiveFactorCancellation}[thm]{%
      11,%
      12,\rAEa{\rAEb{2}},10}}
  \proofstepwidestarbreakkeep[1,2,3,4]{(a,b)\RatPairRel(e,f)}{%
    \FormulaRefAuto{RationalPairRelChar}[thm]{1,2,13}}
\end{tabproofwide}

\FormulaThmDeltaK[Die Bruchpaarrelation ist eine Äquivalenzrelation]{%
  \EqRel{\RatPairCarrier,\RatPairRel}%
}{RationalPairEquivalence}{%
  \DeltaRow{Mengen}{\RatPairCarrier}%
  \DeltaRow{Relationsoperatoren}{\RatPairRel}%
}

\begin{tabproofsplitwide}
\proofpartwide{Reflexivität}
  \proofstepwidestar[1]{x\in\RatPairCarrier}{\rA}
  \proofstepwidestar[]{\RatPairCarrier=\mathbb Z\times\PositiveIntegers}{\FormulaRefAuto{RationalPairCarrier}[def]}
  \proofstepwidestar[1]{x\in\mathbb Z\times\PositiveIntegers}{\rIE{2,1}}
  \proofstepwidestar[1]{\exists a\in\mathbb Z\,\exists b\in\PositiveIntegers\,x=(a,b)}{\FormulaRefAuto{x\in A\times B\eqvdash\exists a\in A\,\exists b\in B\,x=(a,b)}{3}}
  \proofstepwidestar[5]{a\in\mathbb Z\land(b\in\PositiveIntegers\land x=(a,b))}{\rA}
  \proofstepwidestar[5]{(a,b)\RatPairRel(a,b)}{\FormulaRefAuto{RationalPairRelReflexive}[thm]{\rAEa{5},\rAEa{\rAEb{5}}}}
  \proofstepwidestar[5]{x\RatPairRel x}{\rIE{\rAEb{\rAEb{5}},6}}
  \proofstepwidestar[1]{x\RatPairRel x}{\rEE{4,5,7}}
  \proofstepwidestar[]{\forall x\in\RatPairCarrier\,(x\RatPairRel x)}{\rUI{\rRI{1,8}}}
\closeproofpartwide
\proofpartwide{Symmetrie}
  \setcounter{proofstepnr}{9}
  \proofstepwidestar[10]{(x\in\RatPairCarrier\land y\in\RatPairCarrier)\land x\RatPairRel y}{\rA}
  \proofstepwidestar[10]{x\in\mathbb Z\times\PositiveIntegers\land y\in\mathbb Z\times\PositiveIntegers}{\rIE{2,\rAEa{10}}}
  \proofstepwidestar[10]{\exists a\in\mathbb Z\,\exists b\in\PositiveIntegers\,x=(a,b)}{\FormulaRefAuto{x\in A\times B\eqvdash\exists a\in A\,\exists b\in B\,x=(a,b)}{\rAEa{11}}}
  \proofstepwidestar[10]{\exists c\in\mathbb Z\,\exists d\in\PositiveIntegers\,y=(c,d)}{\FormulaRefAuto{x\in A\times B\eqvdash\exists a\in A\,\exists b\in B\,x=(a,b)}{\rAEb{11}}}
  \proofstepwidestar[14]{a\in\mathbb Z\land(b\in\PositiveIntegers\land x=(a,b))}{\rA}
  \proofstepwidestar[15]{c\in\mathbb Z\land(d\in\PositiveIntegers\land y=(c,d))}{\rA}
  \proofstepwidestar[10,14,15]{(a,b)\RatPairRel(c,d)}{\rIE{\rAEb{\rAEb{14}},\rAEb{\rAEb{15}},\rAEb{10}}}
  \proofstepwidestar[10,14,15]{(c,d)\RatPairRel(a,b)}{\FormulaRefAuto{RationalPairRelSymmetric}[thm]{\rAI{\rAEa{14},\rAEa{15}},\rAI{\rAEa{\rAEb{14}},\rAEa{\rAEb{15}}},16}}
  \proofstepwidestar[10,14,15]{y\RatPairRel x}{\rIE{\rAEb{\rAEb{14}},\rAEb{\rAEb{15}},17}}
  \proofstepwidestar[10,14]{y\RatPairRel x}{\rEE{13,15,18}}
  \proofstepwidestar[10]{y\RatPairRel x}{\rEE{12,14,19}}
  \proofstepwidestar[]{\forall x,y\in\RatPairCarrier\,(x\RatPairRel y\rightarrow y\RatPairRel x)}{\rUI{\rRI{10,20}}}
\closeproofpartwide
\proofpartwide{Transitivität}
  \setcounter{proofstepnr}{21}
  \proofstepwidestar[22]{((x\in\RatPairCarrier\land y\in\RatPairCarrier)\land z\in\RatPairCarrier)\land(x\RatPairRel y\land y\RatPairRel z)}{\rA}
  \proofstepwidestar[22]{(x\in\mathbb Z\times\PositiveIntegers\land y\in\mathbb Z\times\PositiveIntegers)\land z\in\mathbb Z\times\PositiveIntegers}{\rIE{2,\rAEa{22}}}
  \proofstepwidestar[22]{\exists a\in\mathbb Z\,\exists b\in\PositiveIntegers\,x=(a,b)}{\FormulaRefAuto{x\in A\times B\eqvdash\exists a\in A\,\exists b\in B\,x=(a,b)}{\rAEa{\rAEa{23}}}}
  \proofstepwidestar[22]{\exists c\in\mathbb Z\,\exists d\in\PositiveIntegers\,y=(c,d)}{\FormulaRefAuto{x\in A\times B\eqvdash\exists a\in A\,\exists b\in B\,x=(a,b)}{\rAEb{\rAEa{23}}}}
  \proofstepwidestar[22]{\exists e\in\mathbb Z\,\exists f\in\PositiveIntegers\,z=(e,f)}{\FormulaRefAuto{x\in A\times B\eqvdash\exists a\in A\,\exists b\in B\,x=(a,b)}{\rAEb{23}}}
  \proofstepwidestar[27]{a\in\mathbb Z\land(b\in\PositiveIntegers\land x=(a,b))}{\rA}
  \proofstepwidestar[28]{c\in\mathbb Z\land(d\in\PositiveIntegers\land y=(c,d))}{\rA}
  \proofstepwidestar[29]{e\in\mathbb Z\land(f\in\PositiveIntegers\land z=(e,f))}{\rA}
  \proofstepwidestar[22,27,28]{(a,b)\RatPairRel(c,d)}{\rIE{\rAEb{\rAEb{27}},\rAEb{\rAEb{28}},\rAEa{\rAEb{22}}}}
  \proofstepwidestar[22,28,29]{(c,d)\RatPairRel(e,f)}{\rIE{\rAEb{\rAEb{28}},\rAEb{\rAEb{29}},\rAEb{\rAEb{22}}}}
  \proofstepwidestar[22,27,28,29]{(a,b)\RatPairRel(e,f)}{\FormulaRefAuto{RationalPairRelTransitive}[thm]{\rAI{\rAEa{27},\rAI{\rAEa{28},\rAEa{29}}},\rAI{\rAEa{\rAEb{27}},\rAI{\rAEa{\rAEb{28}},\rAEa{\rAEb{29}}}},30,31}}
  \proofstepwidestar[22,27,28,29]{x\RatPairRel z}{\rIE{\rAEb{\rAEb{27}},\rAEb{\rAEb{29}},32}}
  \proofstepwidestar[22,27,28]{x\RatPairRel z}{\rEE{26,29,33}}
  \proofstepwidestar[22,27]{x\RatPairRel z}{\rEE{25,28,34}}
  \proofstepwidestar[22]{x\RatPairRel z}{\rEE{24,27,35}}
  \proofstepwidestar[]{\forall x,y,z\in\RatPairCarrier\,(x\RatPairRel y\land y\RatPairRel z\rightarrow x\RatPairRel z)}{\rUI{\rRI{22,36}}}
\closeproofpartwide
\proofpartwide{Äquivalenzrelation}
  \setcounter{proofstepnr}{37}
  \proofstepwidestar[]{\EqRel{\RatPairCarrier,\RatPairRel}}{\FormulaRefAuto{EqRelDef}[def]{9,21,37}}
\closeproofpartwide
\end{tabproofsplitwide}

\section{Quotient und Bruchklassen}

\FormulaDefDeltaK[Menge der rationalen Zahlen]{%
  \mathbb Q\coloneqq
  \QuotientSet{\RatPairCarrier}{\RatPairRel}%
}{RationalsAsQuotient}{%
  \DeltaRow{Mengen}{\RatPairCarrier\dsep\mathbb Q}%
  \DeltaRow{Relationsoperatoren}{\RatPairRel}%
  \DeltaRow{Neue Symbole}{\mathbb Q}%
}

\FormulaDefDeltaK[Klasse eines Bruchpaares]{%
  a\in\mathbb Z\dsep b\in\PositiveIntegers\vdash
  \RatClass{a}{b}\coloneqq
  \EqClass{(a,b)}{\RatPairRel}{\RatPairCarrier}%
}{RationalClassDef}{%
  \DeltaRow{Mengen}{a\dsep b\dsep\RatPairCarrier}%
  \DeltaRow{Relationsoperatoren}{\RatPairRel}%
  \DeltaRow{Neue Symbole}{\RatClass{a}{b}}%
}

\FormulaThmDeltaK[Bruchklassen sind rationale Zahlen]{%
  a\in\mathbb Z\dsep b\in\PositiveIntegers
  \vdash\RatClass{a}{b}\in\mathbb Q%
}{RationalClassInQ}{%
  \DeltaRow{Mengen}{a\dsep b}%
}

\begin{tabproof}
  \proofstep{1}{a\in\mathbb Z}{\rA}
  \proofstep{2}{b\in\PositiveIntegers}{\rA}
  \proofstep{1,2}{(a,b)\in\RatPairCarrier}{%
    \FormulaRefAuto{RationalPairCarrier}[def]{1,2}}
  \proofstep{}{\EqRel{\RatPairCarrier,\RatPairRel}}{%
    \FormulaRefAuto{RationalPairEquivalence}[thm]}
  \proofstep{1,2}{\EqClass{(a,b)}{\RatPairRel}{\RatPairCarrier}
    \in\QuotientSet{\RatPairCarrier}{\RatPairRel}}{%
    \FormulaRefAuto{EqClassInQuotient}[thm]{3,4}}
  \proofstep{1,2}{\mathbb Q = \QuotientSet \RatPairCarrier \RatPairRel}{%
      \FormulaRefAuto{RationalsAsQuotient}[def]}
  \proofstep{1,2}{\RatClass a b = \EqClass { ( a , b ) } \RatPairRel \RatPairCarrier}{%
      \FormulaRefAuto{RationalClassDef}[def]{1,2}}
  \proofstep{1,2}{\RatClass{a}{b}\in\mathbb Q}{%
    \rIE{5,7,%
      6}}
\end{tabproof}

\FormulaThmDeltaK[Gleichheit von Bruchklassen]{%
  a,c\in\mathbb Z\dsep b,d\in\PositiveIntegers\vdash
  \bigl(\RatClass{a}{b}=\RatClass{c}{d}
  \leftrightarrow a\cdot d=c\cdot b\bigr)%
}{RationalClassEquality}{%
  \DeltaRow{Mengen}{a\dsep b\dsep c\dsep d}%
  \DeltaRow{Relationsoperatoren}{\RatPairRel}%
}

\begin{tabproofwide}
  \proofstepwidestar[1]{a,c\in\mathbb Z}{\rA}
  \proofstepwidestar[2]{b,d\in\PositiveIntegers}{\rA}
  \proofstepwidestar[]{\EqRel{\RatPairCarrier,\RatPairRel}}{\FormulaRefAuto{RationalPairEquivalence}[thm]}
  \proofstepwidestar[1,2]{(a,b),(c,d)\in\RatPairCarrier}{\FormulaRefAuto{RationalPairCarrier}[def]{1,2}}
  \proofstepwidestar[1,2]{\EqClass{(a,b)}{\RatPairRel}{\RatPairCarrier}=\EqClass{(c,d)}{\RatPairRel}{\RatPairCarrier}\leftrightarrow(a,b)\RatPairRel(c,d)}{\FormulaRefAuto{EqClassEqualIffRel}[thm]{\rAEa{4},\rAEb{4},3}}
  \proofstepwidestar[1,2]{\RatClass{a}{b}=\EqClass{(a,b)}{\RatPairRel}{\RatPairCarrier}}{\FormulaRefAuto{RationalClassDef}[def]{\rAEa{1},\rAEa{2}}}
  \proofstepwidestar[1,2]{\RatClass{c}{d}=\EqClass{(c,d)}{\RatPairRel}{\RatPairCarrier}}{\FormulaRefAuto{RationalClassDef}[def]{\rAEb{1},\rAEb{2}}}
  \proofstepwidestar[1,2]{\RatClass{a}{b}=\RatClass{c}{d}\leftrightarrow(a,b)\RatPairRel(c,d)}{\rIE{6,7,5}}
  \proofstepwidestar[1,2]{(a,b)\RatPairRel(c,d)\leftrightarrow a\cdot d=c\cdot b}{\FormulaRefAuto{RationalPairRelChar}[thm]{1,2}}
  \proofstepwidestar[1,2]{\RatClass{a}{b}=\RatClass{c}{d}\leftrightarrow a\cdot d=c\cdot b}{\rChain{8,9}}
\end{tabproofwide}

\FormulaThmDeltaK[Jede rationale Zahl besitzt ein Bruchpaar]{%
  q\in\mathbb Q\vdash
  \exists a\in\mathbb Z\,\exists b\in\PositiveIntegers\,
  q=\RatClass{a}{b}%
}{RationalPairRepresentative}{%
  \DeltaRow{Mengen}{q\dsep a\dsep b\dsep\RatPairCarrier}%
}

\begin{tabproofwide}
  \proofstepwidestarbreakkeep[1]{q\in\mathbb Q}{\rA}
  \proofstepwidestar[1]{\mathbb Q = \QuotientSet \RatPairCarrier \RatPairRel}{%
      \FormulaRefAuto{RationalsAsQuotient}[def]}
  \proofstepwidestarbreakkeep[1]{q\in\QuotientSet{\RatPairCarrier}{\RatPairRel}}{%
    \rIE{1,2}}
  \proofstepwidestarbreakkeep[]{\EqRel{\RatPairCarrier,\RatPairRel}}{%
    \FormulaRefAuto{RationalPairEquivalence}[thm]}
  \proofstepwidestarbreakkeep[1]{\exists x\in\RatPairCarrier\,
    q=\EqClass{x}{\RatPairRel}{\RatPairCarrier}}{%
    \FormulaRefAuto{QuotientRepresentative}[thm]{3,4}}
  \proofstepwidestarbreakkeep[6]{x\in\RatPairCarrier\land
    q=\EqClass{x}{\RatPairRel}{\RatPairCarrier}}{\rA}
  \proofstepwidestar[6]{\RatPairCarrier = \mathbb Z \times \PositiveIntegers}{%
      \FormulaRefAuto{RationalPairCarrier}[def]}
  \proofstepwidestarbreakkeep[6]{\exists a\in\mathbb Z\,
    \exists b\in\PositiveIntegers\,x=(a,b)}{%
    \FormulaRefAuto{x\in A\times B\eqvdash
      \exists a\in A\,\exists b\in B\,x=(a,b)}{%
      \rIE{\rAEa{6},7}}}
  \proofstepwidestarbreakkeep[9]{a\in\mathbb Z\land b\in\PositiveIntegers
    \land x=(a,b)}{\rA}
  \proofstepwidestar[6,9]{\RatClass a b = \EqClass { ( a , b ) } \RatPairRel \RatPairCarrier}{%
      \FormulaRefAuto{RationalClassDef}[def]{\rAEa{9},\rAEa{\rAEb{9}}}}
  \proofstepwidestarbreakkeep[6,9]{q=\RatClass{a}{b}}{%
    \rIE{\rAEb{6},\rAEb{\rAEb{9}},%
      10}}
  \proofstepwidestarbreakkeep[6,9]{\exists a\in\mathbb Z\,
    \exists b\in\PositiveIntegers\,q=\RatClass{a}{b}}{%
    \rEI{\rAI{\rAEa{9},\rEI{\rAI{\rAEa{\rAEb{9}},11}}}}}
  \proofstepwidestarbreakkeep[6]{\exists a\in\mathbb Z\,
    \exists b\in\PositiveIntegers\,q=\RatClass{a}{b}}{\rEE{8,9,12}}
  \proofstepwidestarbreakkeep[1]{\exists a\in\mathbb Z\,
    \exists b\in\PositiveIntegers\,q=\RatClass{a}{b}}{\rEE{5,6,13}}
\end{tabproofwide}

\FormulaThmDeltaK[Typisierung der Repräsentantenabbildung]{%
  \forall a\forall b\,\bigl((a\in\mathbb Z\land b\in\PositiveIntegers)\rightarrow\RatClass{a}{b}\in\mathbb Q\bigr)%
}{RationalPairClassTyping}{%
  \DeltaRow{Mengen}{a\dsep b}%
}
\begin{tabproofwide}
  \proofstepwidestar[1]{a\in\mathbb Z\land b\in\PositiveIntegers}{\rA}
  \proofstepwidestar[1]{\RatClass{a}{b}\in\mathbb Q}{%
    \FormulaRefAuto{RationalClassInQ}[thm]{\rAEa{1},\rAEb{1}}}
  \proofstepwidestar[]{\forall a\forall b\,\bigl((a\in\mathbb Z\land b\in\PositiveIntegers)\rightarrow\RatClass{a}{b}\in\mathbb Q\bigr)}{\rUI{\rUI{\rRI{1,2}}}}
\end{tabproofwide}

\FormulaThmDeltaK[Surjektive Repräsentantendarstellung]{%
  \forall z\in\mathbb Q\,\exists a\exists b\,\bigl((a\in\mathbb Z\land b\in\PositiveIntegers)\land z=\RatClass{a}{b}\bigr)%
}{RationalPairRepresentationSurjective}{%
  \DeltaRow{Mengen}{z\dsep a\dsep b}%
}
\begin{tabproofwide}
  \proofstepwidestar[1]{z\in\mathbb Q}{\rA}
  \proofstepwidestar[1]{\exists a\in\mathbb Z\,\exists b\in\PositiveIntegers\,z=\RatClass{a}{b}}{%
    \FormulaRefAuto{RationalPairRepresentative}[thm]{1}}
  \proofstepwidestar[1]{\exists a\exists b\,\bigl((a\in\mathbb Z\land b\in\PositiveIntegers)\land z=\RatClass{a}{b}\bigr)}{%
    \rRE{\rLREa{\FormulaRefAuto{BoundedPairExistsGuard}[thm]},2}}
  \proofstepwidestar[]{\forall z\in\mathbb Q\,\exists a\exists b\,\bigl((a\in\mathbb Z\land b\in\PositiveIntegers)\land z=\RatClass{a}{b}\bigr)}{\rUI{\rRI{1,3}}}
\end{tabproofwide}

\Needspace{16\baselineskip}
\section{Einbettung der ganzen Zahlen}

\FormulaDefDeltaK[Ganzzahlige Einbettung]{%
  \begin{aligned}[t]
    &\RatEmbed\colon\mathbb Z\to\mathbb Q,\\[-2pt]
    &\forall z\in\mathbb Z\,
      \bigl(\RatEmbed(z)\coloneqq\RatClass{z}{\IntOne}\bigr)
  \end{aligned}%
}{IntegerRationalEmbedding}{%
  \DeltaRow{Mengen}{z}%
  \DeltaRow{Funktionensymbol}{\RatEmbed}%
  \DeltaRow{Neue Symbole}{\RatEmbed}%
}
\begin{tabproof}
  \proofstep{1}{z\in\mathbb Z}{\rA}
  \proofstep{}{\IntOne\in\PositiveIntegers}{%
    \FormulaRefAuto{IntegerOnePositive}[thm]}
  \proofstep{1}{\RatClass{z}{\IntOne}\in\mathbb Q}{%
    \FormulaRefAuto{RationalClassInQ}[thm]{1,2}}
  \proofstep{}{%
    \forall z\in\mathbb Z\,
      \bigl(\RatClass{z}{\IntOne}\in\mathbb Q\bigr)}{%
    \rUI{\rRI{1,3}}}
\end{tabproof}

\FormulaThmDeltaK[Typisierung der ganzzahligen Einbettung]{%
  \RatEmbed\colon\mathbb Z\to\mathbb Q%
}{IntegerRationalEmbeddingFunction}{%
  \DeltaRow{Funktion}{\RatEmbed}%
}
\begin{tabproof}
  \proofstep{}{\RatEmbed\colon\mathbb Z\to\mathbb Q}{%
    \FormulaRefAuto{IntegerRationalEmbedding}[def]}
\end{tabproof}

\FormulaThmDeltaK[Injektivität der ganzzahligen Einbettung]{%
  \RatEmbed\colon\mathbb Z\inj\mathbb Q%
}{IntegerRationalEmbeddingInjective}{%
  \DeltaRow{Mengen}{z\dsep w}%
  \DeltaRow{Funktion}{\RatEmbed}%
}

\begin{tabproofsplitwide}
  \proofpartwideindR[Gleichheitsreflexion]{%
    z,w\in\mathbb Z\dsep\RatEmbed(z)=\RatEmbed(w)\vdash z=w%
  }{%
    z,w\in\mathbb Z\dsep\RatEmbed(z)=\RatEmbed(w)\vdash z=w
  }[IntegerRationalEmbeddingReflectsEquality]
    \proofstepwidestarbreakkeep[1]{z,w\in\mathbb Z}{\rA}
    \proofstepwidestarbreakkeep[2]{\RatEmbed(z)=\RatEmbed(w)}{\rA}
    \proofstepwidestarbreakkeep[1]{\RatEmbed(z)=\RatClass{z}{\IntOne}}{%
      \FormulaRefAuto{IntegerRationalEmbedding}[def]{\rAEa{1}}}
    \proofstepwidestarbreakkeep[1]{\RatEmbed(w)=\RatClass{w}{\IntOne}}{%
      \FormulaRefAuto{IntegerRationalEmbedding}[def]{\rAEb{1}}}
    \proofstepwidestar[1,2]{\IntOne \in \PositiveIntegers}{%
      \FormulaRefAuto{IntegerOnePositive}[thm]}
    \proofstepwidestarbreakkeep[1,2]{z\cdot\IntOne=w\cdot\IntOne}{%
      \FormulaRefAuto{RationalClassEquality}[thm]{%
        1,5,\rIE{2,3,4}}}
    \proofstepwidestar[1]{(z\cdot\IntOne=z)\land(w\cdot\IntOne=w)}{%
      \FormulaRefAuto{IntMulOne}[thm]{1}}
    \proofstepwidestarbreakkeep[1,2]{z=w}{%
      \rIE{6,7}}
  \closeproofpartwideindR

  \proofpartwide{Schluss}
    \proofstepwidestarbreakkeep[]{\RatEmbed\colon\mathbb Z\to\mathbb Q}{%
      \FormulaRefAuto{IntegerRationalEmbeddingFunction}[thm]}
    \proofstepwidestarbreakkeep[2]{z,w\in\mathbb Z\land
      \RatEmbed(z)=\RatEmbed(w)}{\rA}
    \proofstepwidestarbreakkeep[2]{z=w}{%
      \FormulaRefAuto{IntegerRationalEmbeddingReflectsEquality}{2}}
    \proofstepwidestarbreakkeep[]{\RatEmbed\colon\mathbb Z\inj\mathbb Q}{%
      \FormulaRefAuto{Injektive Funktion}[def]{1,\rUI{\rRI{2,3}}}}
  \closeproofpartwide
\end{tabproofsplitwide}

\subsection{Kanonische Teilmengenkopien}

Die Bilder der ganzzahligen Einbettung und ihrer Einschränkung auf die
natürliche Kopie bilden innerhalb von \(\mathbb Q\) eine wörtliche
Inklusionskette.

\FormulaDefDeltaK[Kanonische ganzzahlige Kopie in den rationalen Zahlen]{%
  \IntCopyInRat\coloneqq\RatEmbed[\mathbb Z]%
}{IntegerCopyInRationals}{%
  \DeltaRow{Mengen}{\IntCopyInRat}%
  \DeltaRow{Funktion}{\RatEmbed}%
  \DeltaRow{Neue Symbole}{\IntCopyInRat}%
}

\FormulaDefDeltaK[Kanonische natürliche Kopie in den rationalen Zahlen]{%
  \NatCopyInRat\coloneqq\RatEmbed[\NatCopyInInt]%
}{NaturalCopyInRationals}{%
  \DeltaRow{Mengen}{\NatCopyInInt\dsep\NatCopyInRat}%
  \DeltaRow{Funktion}{\RatEmbed}%
  \DeltaRow{Neue Symbole}{\NatCopyInRat}%
}

\FormulaThmDeltaK[Geschachtelte natürliche und ganzzahlige Kopien in
den rationalen Zahlen]{%
  \bigl(\NatCopyInRat\subseteq\IntCopyInRat\land
  \IntCopyInRat\subseteq\mathbb Q\bigr)
  \land\NatCopyInRat\subseteq\mathbb Q%
}{EmbeddedNaturalIntegerRationalTower}{%
  \DeltaRow{Mengen}{\NatCopyInInt\dsep\NatCopyInRat\dsep\IntCopyInRat}%
  \DeltaRow{Funktion}{\RatEmbed}%
}
\begin{tabproofsplitwide}
\proofpartwide{Natürliche Kopie in der ganzzahligen Kopie}
\proofstepwidestar[]{\RatEmbed\colon\mathbb Z\to\mathbb Q}{%
    \FormulaRefAuto{IntegerRationalEmbeddingFunction}[thm]}
  \proofstepwidestar[]{\NatCopyInInt\subseteq\mathbb Z}{%
    \FormulaRefAuto{NaturalCopySubsetIntegers}[thm]}
  \proofstepwidestar[]{\mathbb Z\subseteq\mathbb Z}{%
    \FormulaRefAuto{A\subseteq A}}
  \proofstepwidestar[]{%
    \RatEmbed[\NatCopyInInt]\subseteq\RatEmbed[\mathbb Z]}{%
    \FormulaRefAuto{F\colon M\to N\dsep A\subseteq B\dsep
      B\subseteq M\vdash F[A]\subseteq F[B]}[thm]{1,2,3}}
  \proofstepwidestar[]{\NatCopyInRat=\RatEmbed[\NatCopyInInt]}{%
    \FormulaRefAuto{NaturalCopyInRationals}[def]}
  \proofstepwidestar[]{\IntCopyInRat=\RatEmbed[\mathbb Z]}{%
    \FormulaRefAuto{IntegerCopyInRationals}[def]}
  \proofstepwidestar[]{\NatCopyInRat\subseteq\IntCopyInRat}{\rIE{5,6,4}}

\closeproofpartwide

\proofpartwide{Ganzzahlige Kopie im rationalen Träger}
  \setcounter{proofstepnr}{7}% Fortlaufende Schritte dieses Beweises.
\proofstepwidestar[]{\RatEmbed[\mathbb Z]\subseteq\mathbb Q}{%
    \FormulaRefAuto{F\colon A\to B\vdash F[A]\subseteq B}[thm]{1}}
  \proofstepwidestar[]{\IntCopyInRat\subseteq\mathbb Q}{\rIE{6,8}}

\closeproofpartwide

\proofpartwide{Natürliche Kopie im rationalen Träger}
  \setcounter{proofstepnr}{9}% Fortlaufende Schritte dieses Beweises.
\proofstepwidestar[]{\NatCopyInRat\subseteq\mathbb Q}{%
    \FormulaRefAuto{A\subseteq B\dsep B\subseteq C\vdash A\subseteq C}{7,9}}

\closeproofpartwide

\proofpartwide{Zusammenführung}
  \setcounter{proofstepnr}{10}% Fortlaufende Schritte dieses Beweises.
\proofstepwidestar[]{%
    \bigl(\NatCopyInRat\subseteq\IntCopyInRat\land
    \IntCopyInRat\subseteq\mathbb Q\bigr)
    \land\NatCopyInRat\subseteq\mathbb Q}{\rAI{\rAI{7,9},10}}

\closeproofpartwide
\end{tabproofsplitwide}

\FormulaThmDeltaK[Bijektive Identifikation der natürlichen Kopie]{%
  \RatEmbed\restriction_{\NatCopyInInt}\colon
  \NatCopyInInt\bij\NatCopyInRat%
}{NaturalCopyRationalBijection}{%
  \DeltaRow{Mengen}{\NatCopyInInt\dsep\NatCopyInRat}%
  \DeltaRow{Funktion}{\RatEmbed}%
}
\begin{tabproofwide}
  \proofstepwidestar[]{\RatEmbed\colon\mathbb Z\inj\mathbb Q}{%
    \FormulaRefAuto{IntegerRationalEmbeddingInjective}[thm]}
  \proofstepwidestar[]{\NatCopyInInt\subseteq\mathbb Z}{%
    \FormulaRefAuto{NaturalCopySubsetIntegers}[thm]}
  \proofstepwidestar[]{%
    \RatEmbed\restriction_{\NatCopyInInt}\colon
    \NatCopyInInt\inj\mathbb Q}{%
    \FormulaRefAuto{F\colon A\inj B\dsep C\subseteq A
      \vdash F\restriction_C\colon C\inj B}[thm]{1,2}}
  \proofstepwidestar[]{%
    \RatEmbed\restriction_{\NatCopyInInt}\colon\NatCopyInInt\bij
    (\RatEmbed\restriction_{\NatCopyInInt})[\NatCopyInInt]}{%
    \FormulaRefAuto{F\colon A\inj B\vdash F\colon A\bij F[A]}[thm]{3}}
  \proofstepwidestar[]{\RatEmbed \colon \mathbb Z \to \mathbb Q}{%
      \FormulaRefAuto{Injektive Funktion}[def]{1}}
  \proofstepwidestar[]{%
    (\RatEmbed\restriction_{\NatCopyInInt})[\NatCopyInInt]
    =\RatEmbed[\NatCopyInInt]}{%
    \FormulaRefAuto{C\subseteq A\dsep F\colon A\to B
      \vdash (F\restriction_C)[C]=F[C]}[thm]{%
        2,5}}
  \proofstepwidestar[]{\NatCopyInRat=\RatEmbed[\NatCopyInInt]}{%
    \FormulaRefAuto{NaturalCopyInRationals}[def]}
  \proofstepwidestar[]{%
    \RatEmbed\restriction_{\NatCopyInInt}\colon
    \NatCopyInInt\bij\NatCopyInRat}{\rIE{6,7,4}}
\end{tabproofwide}

\FormulaThmDeltaK[Bijektive Identifikation der ganzzahligen Kopie]{%
  \RatEmbed\colon\mathbb Z\bij\IntCopyInRat%
}{IntegerCopyRationalBijection}{%
  \DeltaRow{Mengen}{\IntCopyInRat}%
  \DeltaRow{Funktion}{\RatEmbed}%
}
\begin{tabproof}
  \proofstep{}{\RatEmbed\colon\mathbb Z\inj\mathbb Q}{%
    \FormulaRefAuto{IntegerRationalEmbeddingInjective}[thm]}
  \proofstep{}{\RatEmbed\colon\mathbb Z\bij\RatEmbed[\mathbb Z]}{%
    \FormulaRefAuto{F\colon A\inj B\vdash F\colon A\bij F[A]}[thm]{1}}
  \proofstep{}{\IntCopyInRat=\RatEmbed[\mathbb Z]}{%
    \FormulaRefAuto{IntegerCopyInRationals}[def]}
  \proofstep{}{\RatEmbed\colon\mathbb Z\bij\IntCopyInRat}{\rIE{3,2}}
\end{tabproof}

Im rationalen Kontext verwenden wir dieselben sichtbaren Zahlzeichen wie in
\(\mathbb N\) und \(\mathbb Z\).  Semantisch bezeichnet
\(\RatNumeral{n}\) für \(n\in\mathbb N\) das Bild
\(\RatOfNat{n}\).  Die wörtlichen Inklusionen werden durch
\(\NatCopyInRat\subseteq\IntCopyInRat\subseteq\mathbb Q\) ausgedrückt;
das Weglassen der Kopienindizes bleibt eine durch die Einbettungen
gerechtfertigte Identifikation.  Wo der umgebende Träger den Typ nicht
eindeutig bestimmt, bleiben \(\RatOfInt{z}\) beziehungsweise
\(\RatOfNat{n}\) ausdrücklich sichtbar.  Insbesondere stehen
\(\RatZero,\RatOne,\RatNumeral{2},
\RatNumeral{3},\RatNumeral{4},\ldots\) ohne Trägerindex.

\FormulaDefDeltaK[Unindizierte rationale Zahlzeichen]{%
  n\in\mathbb N\vdash
  \RatNumeral{n}\coloneqq\RatOfNat{n}%
}{RationalNumeralConvention}{%
  \DeltaRow{Mengen}{n}%
  \DeltaRow{Funktionensymbole}{\IntEmbed\dsep\RatEmbed}%
  \DeltaRow{Neue Schreibweise}{\RatNumeral{n}}%
}

\FormulaDefDeltaK[Null der rationalen Zahlen]{%
  \RatZero\coloneqq\RatOfInt{\IntZero}%
}{RationalZeroDef}{%
  \DeltaRow{Neue Symbole}{\RatZero}%
  \DeltaRow{Funktionensymbol}{\RatEmbed}%
}

\FormulaDefDeltaK[Eins der rationalen Zahlen]{%
  \RatOne\coloneqq\RatOfInt{\IntOne}%
}{RationalOneDef}{%
  \DeltaRow{Neue Symbole}{\RatOne}%
  \DeltaRow{Funktionensymbol}{\RatEmbed}%
}

\FormulaThmDeltaKR[Null- und Einsklasse]{%
  \begin{aligned}[t]
    &\text{(i)}\quad \RatZero=\RatClass{\IntZero}{\IntOne},\\[-2pt]
    &\text{(ii)}\quad \RatOne=\RatClass{\IntOne}{\IntOne}.
  \end{aligned}%
}{%
  \RatZero=\RatClass{\IntZero}{\IntOne}
  \dsep
  \RatOne=\RatClass{\IntOne}{\IntOne}%
}{RationalZeroOneClasses}{%
  \DeltaRow{Mengen}{\mathbb Q}%
}
\begin{tabproofsplit}
\proofpart{Nullklasse}
\proofstep{}{\RatZero=\RatEmbed(\IntZero)}{%
    \FormulaRefAuto{RationalZeroDef}[def]}
  \proofstep{}{\IntZero \in \mathbb Z}{%
      \FormulaRefAuto{IntZeroCarrier}[thm]}
  \proofstep{}{\RatEmbed(\IntZero)
    =\RatClass{\IntZero}{\IntOne}}{%
    \FormulaRefAuto{IntegerRationalEmbedding}[def]{%
      2}}
  \proofstep{}{\RatZero=\RatClass{\IntZero}{\IntOne}}{%
    \rChain{1,3}}

\closeproofpart

\proofpart{Einsklasse}
  \setcounter{proofstepnr}{4}% Fortlaufende Schritte dieses Beweises.
\proofstep{}{\RatOne=\RatEmbed(\IntOne)}{%
    \FormulaRefAuto{RationalOneDef}[def]}
  \proofstep{}{\IntOne \in \mathbb Z}{%
      \FormulaRefAuto{IntOneCarrier}[thm]}
  \proofstep{}{\RatEmbed(\IntOne)
    =\RatClass{\IntOne}{\IntOne}}{%
    \FormulaRefAuto{IntegerRationalEmbedding}[def]{%
      6}}
  \proofstep{}{\RatOne=\RatClass{\IntOne}{\IntOne}}{%
    \rChain{5,7}}

\closeproofpart

\proofpart{Zusammenführung}
  \setcounter{proofstepnr}{8}% Fortlaufende Schritte dieses Beweises.
\proofstep{}{\RatZero=\RatClass{\IntZero}{\IntOne}
    \land\RatOne=\RatClass{\IntOne}{\IntOne}}{%
    \rAI{4,8}}

\closeproofpart
\end{tabproofsplit}

\FormulaThmDeltaKR[Null und Eins sind rationale Zahlen]{%
  \begin{aligned}[t]
    &\text{(i)}\quad \RatZero\in\mathbb Q,\\[-2pt]
    &\text{(ii)}\quad \RatOne\in\mathbb Q.
  \end{aligned}%
}{%
  \RatZero\in\mathbb Q\dsep\RatOne\in\mathbb Q%
}{RationalZeroOneCarrier}{%
  \DeltaRow{Mengen}{\mathbb Q}%
}

\begin{tabproofsplit}
\proofpart{Null im rationalen Träger}
\proofstep{}{\RatEmbed\colon\mathbb Z\to\mathbb Q}{%
    \FormulaRefAuto{IntegerRationalEmbeddingFunction}[thm]}
  \proofstep{}{\IntZero\in\mathbb Z}{%
    \FormulaRefAuto{IntZeroCarrier}[thm]}
  \proofstep{}{\IntOne\in\mathbb Z}{%
    \FormulaRefAuto{IntOneCarrier}[thm]}
  \proofstep{}{\RatEmbed(\IntZero)\in\mathbb Q}{%
    \FormulaRefAuto{F\colon A\to B\dsep x\in A\vdash F(x)\in B}{1,2}}
  \proofstep{}{\RatEmbed(\IntOne)\in\mathbb Q}{%
    \FormulaRefAuto{F\colon A\to B\dsep x\in A\vdash F(x)\in B}{1,3}}
  \proofstep{}{\RatZero = \RatOfInt \IntZero}{%
      \FormulaRefAuto{RationalZeroDef}[def]}
  \proofstep{}{\RatZero\in\mathbb Q}{%
    \rIE{4,6}}

\closeproofpart

\proofpart{Eins im rationalen Träger}
  \setcounter{proofstepnr}{7}% Fortlaufende Schritte dieses Beweises.
\proofstep{}{\RatOne = \RatOfInt \IntOne}{%
      \FormulaRefAuto{RationalOneDef}[def]}
\proofstep{}{\RatOne\in\mathbb Q}{%
    \rIE{5,8}}

\closeproofpart

\proofpart{Zusammenführung}
  \setcounter{proofstepnr}{9}% Fortlaufende Schritte dieses Beweises.
\proofstep{}{\RatZero\in\mathbb Q\land\RatOne\in\mathbb Q}{%
    \rAI{7,9}}

\closeproofpart
\end{tabproofsplit}

\chapter{Arithmetik der rationalen Zahlen}

\section{Negation}

\FormulaThmDeltaK[Repräsentantenunabhängigkeit der rationalen Negation]{%
  a,c\in\mathbb Z\dsep b,d\in\PositiveIntegers\dsep
  \RatClass{a}{b}=\RatClass{c}{d}
  \vdash\RatClass{-a}{b}=\RatClass{-c}{d}%
}{RationalNegationCompatible}{%
  \DeltaRow{Mengen}{a\dsep b\dsep c\dsep d}%
}

\begin{tabproofwide}
  \proofstepwidestarbreakkeep[1]{a,c\in\mathbb Z}{\rA}
  \proofstepwidestarbreakkeep[2]{b,d\in\PositiveIntegers}{\rA}
  \proofstepwidestarbreakkeep[3]{\RatClass{a}{b}=\RatClass{c}{d}}{\rA}
  \proofstepwidestarbreakkeep[1,2,3]{a\cdot d=c\cdot b}{%
    \FormulaRefAuto{RationalClassEquality}[thm]{1,2,3}}
  \proofstepwide[1,2,3]{(-a)\cdot d}{=}{-(a\cdot d)}{%
    \FormulaRefAuto{IntMulLeftDistributive}[thm],
    \FormulaRefAuto{IntAddInverse}[thm]}
  \proofstepwide[1,2,3]{-(a\cdot d)}{=}{-(c\cdot b)}{\rIE{4}}
  \proofstepwide[1,2,3]{-(c\cdot b)}{=}{(-c)\cdot b}{%
    \FormulaRefAuto{IntMulLeftDistributive}[thm],
    \FormulaRefAuto{IntAddInverse}[thm]}
  \proofstepwidestarbreakkeep[1,2,3]{(-a)\cdot d=(-c)\cdot b}{\rChain{5,6,7}}
  \proofstepwidestar[1]{-a,-c\in\mathbb Z}{%
      \FormulaRefAuto{IntNegClosure}[thm]{1}}
  \proofstepwidestarbreakkeep[1,2,3]{\RatClass{-a}{b}=\RatClass{-c}{d}}{%
    \FormulaRefAuto{RationalClassEquality}[thm]{%
      9,2,8}}
\end{tabproofwide}

\Needspace{14\baselineskip}
\FormulaThmDeltaK[Quotientenabstieg der rationalen Negation]{%
  \begin{aligned}[t]
    &q\in\mathbb Q\vdash
      \exists!u\in\mathbb Q\,
      \exists a\in\mathbb Z\,\exists b\in\PositiveIntegers\,\Bigl({}\\[-2pt]
    &\qquad q=\RatClass{a}{b}\ \land\\[-2pt]
    &\qquad u=\RatClass{-a}{b}\Bigr)
  \end{aligned}%
}{RationalNegationQuotientDescent}{%
  \DeltaRow{Mengen}{q\dsep u\dsep a\dsep b}%
  \DeltaRow{Weitere Mengen}{z\dsep c\dsep d}%
}
\begin{tabproofwide}
  \proofstepwidestar[1]{q\in\mathbb Q}{%
    \rA}
  \proofstepwidestar[]{\forall a\forall b\,\bigl((a\in\mathbb Z\land b\in\PositiveIntegers)\rightarrow\RatClass{a}{b}\in\mathbb Q\bigr)}{%
    \FormulaRefAuto{RationalPairClassTyping}[thm]}
  \proofstepwidestar[]{\begin{aligned}[t]
&\forall z\in\mathbb Q\,\exists a\exists b\,\\[-2pt]
&\bigl((a\in\mathbb Z\land b\in\PositiveIntegers)\land z=\RatClass{a}{b}\bigr)
\end{aligned}}{%
    \FormulaRefAuto{RationalPairRepresentationSurjective}[thm]}
  \proofstepwidestar[4]{a\in\mathbb Z\land b\in\PositiveIntegers}{%
    \rA}
  \proofstepwidestar[4]{-a\in\mathbb Z}{%
    \FormulaRefAuto{IntNegClosure}[thm]{\rAEa{4}}}
  \proofstepwidestar[4]{\RatClass{-a}{b}\in\mathbb Q}{%
    \FormulaRefAuto{RationalClassInQ}[thm]{5,\rAEb{4}}}
  \proofstepwidestar[]{\begin{aligned}[t]
&\forall a\forall b\,\bigl((a\in\mathbb Z\land b\in\PositiveIntegers)\rightarrow\\[-2pt]
&\qquad \RatClass{-a}{b}\in\mathbb Q\bigr)
\end{aligned}}{%
    \rUI{\rUI{\rRI{4,6}}}}
  \proofstepwidestar[8]{\begin{aligned}[t]
&(a\in\mathbb Z\land b\in\PositiveIntegers)\land(c\in\mathbb Z\land d\in\PositiveIntegers)\\[-2pt]&\land\RatClass{a}{b}=\RatClass{c}{d}
\end{aligned}}{%
    \rA}
  \proofstepwidestar[8]{\RatClass{-a}{b}=\RatClass{-c}{d}}{%
    \FormulaRefAuto{RationalNegationCompatible}[thm]{\rAI{\rAEa{\rAEa{8}},\rAEa{\rAEa{\rAEb{8}}}},\rAI{\rAEb{\rAEa{8}},\rAEb{\rAEa{\rAEb{8}}}},\rAEb{\rAEb{8}}}}
  \proofstepwidestar[]{\begin{aligned}[t]
&\forall a\forall b\forall c\forall d\,\Bigl(\\[-2pt]
&(a\in\mathbb Z\land b\in\PositiveIntegers)\land(c\in\mathbb Z\land d\in\PositiveIntegers)\\[-2pt]&\land\RatClass{a}{b}=\RatClass{c}{d}\rightarrow\\[-2pt]
&\qquad \RatClass{-a}{b}=\RatClass{-c}{d}\Bigr)
\end{aligned}}{%
    \rUI{\rUI{\rUI{\rUI{\rRI{8,9}}}}}}
  \proofstepwidestar[1]{\begin{aligned}[t]
&\exists!u\in\mathbb Q\,\exists a\exists b\,\\[-2pt]
&\bigl((a\in\mathbb Z\land b\in\PositiveIntegers)\land q=\RatClass{a}{b}\land u=\RatClass{-a}{b}\bigr)
\end{aligned}}{%
    \FormulaRefAuto{RepresentativePairValueUnique}[thm]{2,3,7,10,1}}
  \proofstepwidestar[]{\begin{aligned}[t]
&\forall u\,\Bigl(\exists a\exists b\,\\[-2pt]
&\bigl((a\in\mathbb Z\land b\in\PositiveIntegers)\land q=\RatClass{a}{b}\land u=\RatClass{-a}{b}\bigr)\\[-2pt]
&\quad\leftrightarrow{}\exists a\in\mathbb Z\,\exists b\in\PositiveIntegers\,\\[-2pt]
&\bigl(q=\RatClass{a}{b}\land u=\RatClass{-a}{b}\bigr)\Bigr)
\end{aligned}}{%
    \rUI{\FormulaRefAuto{P \leftrightarrow Q \vdash Q \leftrightarrow P}[thm]{\FormulaRefAuto{BoundedPairExistsGuard}[thm]}}}
  \proofstepwidestar[1]{\begin{aligned}[t]
&\exists!u\in\mathbb Q\,\exists a\in\mathbb Z\,\exists b\in\PositiveIntegers\,\\[-2pt]
&\bigl(q=\RatClass{a}{b}\land u=\RatClass{-a}{b}\bigr)
\end{aligned}}{%
    \FormulaRefAuto{UniqueExistenceUnderEquivalentConjuncts}[thm]{12,11}}
\end{tabproofwide}

\Needspace{12\baselineskip}
\FormulaDefDeltaK[Negation rationaler Zahlen]{%
  q\in\mathbb Q\vdash
  -q\coloneqq\iota u\Bigl(
    u\in\mathbb Q\land
    \exists a\in\mathbb Z\,\exists b\in\PositiveIntegers\,
    (q=\RatClass{a}{b}\land u=\RatClass{-a}{b})
  \Bigr)%
}{RationalNegationDef}{%
  \DeltaRow{Mengen}{q\dsep u\dsep a\dsep b}%
  \DeltaRow{Neue Symbole}{-q}%
}

\FormulaThmDeltaK[Negation auf Bruchklassen]{%
  a\in\mathbb Z\dsep b\in\PositiveIntegers\vdash
  -\RatClass{a}{b}=\RatClass{-a}{b}%
}{RationalNegationClassValue}{%
  \DeltaRow{Mengen}{a\dsep b}%
}
\begin{tabproofwide}
  \proofstepwidestarbreakkeep[1]{a\in\mathbb Z}{\rA}
  \proofstepwidestarbreakkeep[2]{b\in\PositiveIntegers}{\rA}
  \proofstepwidestarbreakkeep[1,2]{\RatClass{a}{b}\in\mathbb Q}{%
    \FormulaRefAuto{RationalClassInQ}[thm]{1,2}}
  \proofstepwidestar[1]{-a\in\mathbb Z}{%
      \FormulaRefAuto{IntNegClosure}[thm]{1}}
  \proofstepwidestarbreakkeep[1,2]{\RatClass{-a}{b}\in\mathbb Q}{%
    \FormulaRefAuto{RationalClassInQ}[thm]{%
      4,2}}
  \proofstepwidestarbreakkeep[1,2]{\begin{aligned}[t]
    &\exists!u\in\mathbb Q\,
      \exists c\in\mathbb Z\,\exists d\in\PositiveIntegers\,{}\\
    &\bigl(\RatClass{a}{b}=\RatClass{c}{d}
      \land u=\RatClass{-c}{d}\bigr)
    \end{aligned}}{%
    \FormulaRefAuto{RationalNegationQuotientDescent}[thm]{3}}
  \proofstepwidestarbreakkeep[1,2]{-\RatClass{a}{b}=\RatClass{-a}{b}}{%
    \FormulaRefAuto{RationalNegationDef}[def]{1,2,5,6}}
\end{tabproofwide}

\FormulaThmDeltaKR[Abgeschlossenheit und doppelte Negation]{%
  \begin{aligned}[t]
    &q\in\mathbb Q\vdash{}\\[-2pt]
    &\text{(i)}\quad -q\in\mathbb Q,\\[-2pt]
    &\text{(ii)}\quad -(-q)=q.
  \end{aligned}%
}{%
  q\in\mathbb Q\vdash
  -q\in\mathbb Q\dsep-(-q)=q%
}{RationalNegationLaws}{%
  \DeltaRow{Mengen}{q}%
}


\begin{tabproofsplitwide}
\proofpartwide{Abgeschlossenheit der Negation}
\proofstepwidestarbreakkeep[1]{q\in\mathbb Q}{\rA}
  \proofstepwidestarbreakkeep[1]{\exists a\in\mathbb Z\,
    \exists b\in\PositiveIntegers\,q=\RatClass{a}{b}}{%
    \FormulaRefAuto{RationalPairRepresentative}[thm]{1}}
  \proofstepwidestarbreakkeep[3]{a\in\mathbb Z\land b\in\PositiveIntegers\land
    q=\RatClass{a}{b}}{\rA}
  \proofstepwidestar[3]{- \RatClass a b = \RatClass { - a } b}{%
      \FormulaRefAuto{RationalNegationClassValue}[thm]{%
        \rAEa{3},\rAEa{\rAEb{3}}}}
  \proofstepwidestarbreakkeep[3]{-q=\RatClass{-a}{b}}{%
    \rIE{\rAEb{\rAEb{3}},%
      4}}
  \proofstepwidestar[3]{- a \in \mathbb Z}{%
      \FormulaRefAuto{IntNegClosure}[thm]{\rAEa{3}}}
  \proofstepwidestar[3]{\RatClass { ( - a ) } b \in \mathbb Q}{%
      \FormulaRefAuto{RationalClassInQ}[thm]{%
      6,\rAEa{\rAEb{3}}}}
  \proofstepwidestarbreakkeep[3]{-q\in\mathbb Q}{%
    \rIE{5,7}}

\closeproofpartwide

\proofpartwide{Doppelte Negation}
  \setcounter{proofstepnr}{8}% Fortlaufende Schritte dieses Beweises.
\proofstepwidestar[3]{- a \in \mathbb Z}{%
      \FormulaRefAuto{IntNegClosure}[thm]{\rAEa{3}}}
\proofstepwidestar[3]{- \RatClass { ( - a ) } b = \RatClass { - { ( - a ) } } b}{%
      \FormulaRefAuto{RationalNegationClassValue}[thm]{%
      9,\rAEa{\rAEb{3}}}}
\proofstepwidestarbreakkeep[3]{-(-q)=\RatClass{-(-a)}{b}}{%
    \rIE{5,10}}
  \proofstepwidestar[3]{- ( - a ) = a}{%
      \FormulaRefAuto{IntegerDoubleNegation}[thm]{\rAEa{3}}}
  \proofstepwidestarbreakkeep[3]{-(-q)=q}{%
    \rIE{11,12,%
      \rAEb{\rAEb{3}}}}

\closeproofpartwide

\proofpartwide{Zusammenführung}
  \setcounter{proofstepnr}{13}% Fortlaufende Schritte dieses Beweises.
\proofstepwidestarbreakkeep[3]{-q\in\mathbb Q\land-(-q)=q}{\rAI{8,13}}
  \proofstepwidestarbreakkeep[1]{-q\in\mathbb Q\land-(-q)=q}{\rEE{2,3,14}}

\closeproofpartwide
\end{tabproofsplitwide}

\section{Addition}

\FormulaThmDeltaK[Repräsentantenunabhängigkeit der rationalen Addition]{%
  \begin{aligned}[t]
    &a,a',c,c'\in\mathbb Z\dsep
      b,b',d,d'\in\PositiveIntegers\dsep
      \RatClass{a}{b}=\RatClass{a'}{b'}\dsep{}\\
    &\RatClass{c}{d}=\RatClass{c'}{d'}\vdash
      \RatClass{a\cdot d+c\cdot b}{b\cdot d}
      =\RatClass{a'\cdot d'+c'\cdot b'}{b'\cdot d'}
  \end{aligned}%
}{RationalAdditionCompatible}{%
  \DeltaRow{Mengen}{a\dsep a'\dsep b\dsep b'\dsep
    c\dsep c'\dsep d\dsep d'}%
}

\begin{tabproofwide}
  \proofstepwidestarbreakkeep[1]{a,a',c,c'\in\mathbb Z}{\rA}
  \proofstepwidestarbreakkeep[2]{b,b',d,d'\in\PositiveIntegers}{\rA}
  \proofstepwidestarbreakkeep[3]{\RatClass{a}{b}=\RatClass{a'}{b'}}{\rA}
  \proofstepwidestarbreakkeep[4]{\RatClass{c}{d}=\RatClass{c'}{d'}}{\rA}
  \proofstepwidestarbreakkeep[1,2,3]{a\cdot b'=a'\cdot b}{%
    \FormulaRefAuto{RationalClassEquality}[thm]{1,2,3}}
  \proofstepwidestarbreakkeep[1,2,4]{c\cdot d'=c'\cdot d}{%
    \FormulaRefAuto{RationalClassEquality}[thm]{1,2,4}}
  \proofstepwidestarbreakkeep[2]{b\cdot d\in\PositiveIntegers}{%
    \FormulaRefAuto{PositiveIntegerMulClosure}[thm]{2}}
  \proofstepwidestarbreakkeep[2]{b'\cdot d'\in\PositiveIntegers}{%
    \FormulaRefAuto{PositiveIntegerMulClosure}[thm]{2}}
  \proofstepwidestar[1]{a \cdot d \in \mathbb Z}{%
      \FormulaRefAuto{IntMulClosure}[thm]{1}}
  \proofstepwidestar[1]{c \cdot b \in \mathbb Z}{%
      \FormulaRefAuto{IntMulClosure}[thm]{1}}
  \proofstepwidestarbreakkeep[1]{a\cdot d+c\cdot b\in\mathbb Z}{%
    \FormulaRefAuto{IntAddClosure}[thm]{%
      9,%
      10}}
  \proofstepwidestar[1]{{ a ' } \cdot { d ' } \in \mathbb Z}{%
      \FormulaRefAuto{IntMulClosure}[thm]{1}}
  \proofstepwidestar[1]{{ c ' } \cdot { b ' } \in \mathbb Z}{%
      \FormulaRefAuto{IntMulClosure}[thm]{1}}
  \proofstepwidestarbreakkeep[1]{a'\cdot d'+c'\cdot b'\in\mathbb Z}{%
    \FormulaRefAuto{IntAddClosure}[thm]{%
      12,%
      13}}
  \proofstepwidestarbreakkeep[1,2,3,4]{%
    (a\cdot d+c\cdot b)\cdot(b'\cdot d')
    =(a'\cdot d'+c'\cdot b')\cdot(b\cdot d)}{%
    \rIE{5,6,\FormulaRefAuto{IntMulAssociative}[thm],%
      \FormulaRefAuto{IntMulCommutative}[thm],%
      \FormulaRefAuto{IntMulLeftDistributive}[thm],%
      \FormulaRefAuto{IntMulRightDistributive}[thm]}}
  \proofstepwidestarbreakkeep[1,2,3,4]{%
    \RatClass{a\cdot d+c\cdot b}{b\cdot d}
    =\RatClass{a'\cdot d'+c'\cdot b'}{b'\cdot d'}}{%
    \FormulaRefAuto{RationalClassEquality}[thm]{11,14,7,8,15}}
\end{tabproofwide}

\FormulaThmDeltaK[Quotientenabstieg der rationalen Addition]{%
  \begin{aligned}[t]
    &q,r\in\mathbb Q\vdash
      \exists!u\in\mathbb Q\,
      \exists a,c\in\mathbb Z\,\exists b,d\in\PositiveIntegers\,\Bigl({}\\[-2pt]
    &\qquad q=\RatClass{a}{b}\ \land\\[-2pt]
    &\qquad r=\RatClass{c}{d}\ \land\\[-2pt]
    &\qquad u=\RatClass{a\cdot d+c\cdot b}{b\cdot d}\Bigr)
  \end{aligned}%
}{RationalAdditionQuotientDescent}{%
  \DeltaRow{Mengen}{q\dsep r\dsep u\dsep v\dsep
    a\dsep a'\dsep b\dsep b'\dsep c\dsep c'\dsep d\dsep d'}%
  \DeltaRow{Weitere Mengen}{e\dsep f\dsep g\dsep h\dsep z}%
}

\begin{tabproofwide}
  \proofstepwidestar[1]{q,r\in\mathbb Q}{%
    \rA}
  \proofstepwidestar[]{\forall a\forall b\,\bigl((a\in\mathbb Z\land b\in\PositiveIntegers)\rightarrow\RatClass{a}{b}\in\mathbb Q\bigr)}{%
    \FormulaRefAuto{RationalPairClassTyping}[thm]}
  \proofstepwidestar[]{\begin{aligned}[t]
&\forall z\in\mathbb Q\,\exists a\exists b\,\\[-2pt]
&\bigl((a\in\mathbb Z\land b\in\PositiveIntegers)\land z=\RatClass{a}{b}\bigr)
\end{aligned}}{%
    \FormulaRefAuto{RationalPairRepresentationSurjective}[thm]}
  \proofstepwidestar[4]{(a\in\mathbb Z\land b\in\PositiveIntegers)\land(c\in\mathbb Z\land d\in\PositiveIntegers)}{%
    \rA}
  \proofstepwidestar[4]{b\in\mathbb Z}{%
    \rAEa{\FormulaRefAuto{PositiveIntegersDef}[def]{\rAEb{\rAEa{4}}}}}
  \proofstepwidestar[4]{d\in\mathbb Z}{%
    \rAEa{\FormulaRefAuto{PositiveIntegersDef}[def]{\rAEb{\rAEb{4}}}}}
  \proofstepwidestar[4]{a\cdot d\in\mathbb Z}{%
    \FormulaRefAuto{IntMulClosure}[thm]{\rAEa{\rAEa{4}},6}}
  \proofstepwidestar[4]{c\cdot b\in\mathbb Z}{%
    \FormulaRefAuto{IntMulClosure}[thm]{\rAEa{\rAEb{4}},5}}
  \proofstepwidestar[4]{a\cdot d+c\cdot b\in\mathbb Z}{%
    \FormulaRefAuto{IntAddClosure}[thm]{7,8}}
  \proofstepwidestar[4]{b\cdot d\in\PositiveIntegers}{%
    \FormulaRefAuto{PositiveIntegerMulClosure}[thm]{\rAEb{\rAEa{4}},\rAEb{\rAEb{4}}}}
  \proofstepwidestar[4]{\RatClass{a\cdot d+c\cdot b}{b\cdot d}\in\mathbb Q}{%
    \FormulaRefAuto{RationalClassInQ}[thm]{9,10}}
  \proofstepwidestar[]{\begin{aligned}[t]
&\forall a\forall b\forall c\forall d\\[-2pt]&\bigl(((a\in\mathbb Z\land b\in\PositiveIntegers)\land(c\in\mathbb Z\land d\in\PositiveIntegers))\rightarrow\\[-2pt]
&\qquad \RatClass{a\cdot d+c\cdot b}{b\cdot d}\in\mathbb Q\bigr)
\end{aligned}}{%
    \rUI{\rUI{\rUI{\rUI{\rRI{4,11}}}}}}
  \proofstepwidestar[13]{\begin{aligned}[t]
&(a\in\mathbb Z\land b\in\PositiveIntegers)\land(c\in\mathbb Z\land d\in\PositiveIntegers)\land{}\\[-2pt]
&(e\in\mathbb Z\land f\in\PositiveIntegers)\land(g\in\mathbb Z\land h\in\PositiveIntegers)\land{}\\[-2pt]
&\RatClass{a}{b}=\RatClass{e}{f}\land \RatClass{c}{d}=\RatClass{g}{h}
\end{aligned}}{%
    \rA}
  \proofstepwidestar[13]{\RatClass{a\cdot d+c\cdot b}{b\cdot d}=\RatClass{e\cdot h+g\cdot f}{f\cdot h}}{%
    \FormulaRefAuto{RationalAdditionCompatible}[thm]{\rAIStar{\rAEa{\rAEa{13}},\rAEa{\rAEa{\rAEb{\rAEb{13}}}},\rAEa{\rAEa{\rAEb{13}}},\rAEa{\rAEa{\rAEb{\rAEb{\rAEb{13}}}}}},\rAIStar{\rAEb{\rAEa{13}},\rAEb{\rAEa{\rAEb{\rAEb{13}}}},\rAEb{\rAEa{\rAEb{13}}},\rAEb{\rAEa{\rAEb{\rAEb{\rAEb{13}}}}}},\rAEa{\rAEb{\rAEb{\rAEb{\rAEb{13}}}}},\rAEb{\rAEb{\rAEb{\rAEb{\rAEb{13}}}}}}}
  \proofstepwidestar[]{\begin{aligned}[t]
&\forall a\forall b\forall c\forall d\forall e\forall f\forall g\forall h\,\Bigl(\\[-2pt]
&(a\in\mathbb Z\land b\in\PositiveIntegers)\land(c\in\mathbb Z\land d\in\PositiveIntegers)\land{}\\[-2pt]
&(e\in\mathbb Z\land f\in\PositiveIntegers)\land(g\in\mathbb Z\land h\in\PositiveIntegers)\land{}\\[-2pt]
&\RatClass{a}{b}=\RatClass{e}{f}\land \RatClass{c}{d}=\RatClass{g}{h}\rightarrow\\[-2pt]
&\qquad \RatClass{a\cdot d+c\cdot b}{b\cdot d}=\RatClass{e\cdot h+g\cdot f}{f\cdot h}\Bigr)
\end{aligned}}{%
    \rUI{\rUI{\rUI{\rUI{\rUI{\rUI{\rUI{\rUI{\rRI{13,14}}}}}}}}}}
  \proofstepwidestar[1]{\begin{aligned}[t]
&\exists!u\in\mathbb Q\,\exists a\exists b\exists c\exists d\,\\[-2pt]
&\bigl((a\in\mathbb Z\land b\in\PositiveIntegers)\land(c\in\mathbb Z\land d\in\PositiveIntegers)\\[-2pt]&\land q=\RatClass{a}{b}\land{}\\[-2pt]
&\qquad r=\RatClass{c}{d}\land u=\RatClass{a\cdot d+c\cdot b}{b\cdot d}\bigr)
\end{aligned}}{%
    \FormulaRefAuto{RepresentativePairBinaryValueUnique}[thm]{2,3,12,15,\rAEa{1},\rAEb{1}}}
  \proofstepwidestar[]{\begin{aligned}[t]
&\forall u\,\Bigl(\exists a\exists b\exists c\exists d\,\\[-2pt]
&\bigl((a\in\mathbb Z\land b\in\PositiveIntegers)\land(c\in\mathbb Z\land d\in\PositiveIntegers)\\[-2pt]&\land q=\RatClass{a}{b}\land{}\\[-2pt]
&\qquad r=\RatClass{c}{d}\land u=\RatClass{a\cdot d+c\cdot b}{b\cdot d}\bigr)\\[-2pt]
&\quad\leftrightarrow{}\exists a,c\in\mathbb Z\,\exists b,d\in\PositiveIntegers\,\\[-2pt]
&\bigl(q=\RatClass{a}{b}\land r=\RatClass{c}{d}\land{}\\[-2pt]
&\qquad u=\RatClass{a\cdot d+c\cdot b}{b\cdot d}\bigr)\Bigr)
\end{aligned}}{%
    \rUI{\FormulaRefAuto{P \leftrightarrow Q \vdash Q \leftrightarrow P}[thm]{\FormulaRefAuto{BoundedFourExistsGroupedGuard}[thm]}}}
  \proofstepwidestar[1]{\begin{aligned}[t]
&\exists!u\in\mathbb Q\,\exists a,c\in\mathbb Z\,\exists b,d\in\PositiveIntegers\,\\[-2pt]
&\bigl(q=\RatClass{a}{b}\land r=\RatClass{c}{d}\land{}\\[-2pt]
&\qquad u=\RatClass{a\cdot d+c\cdot b}{b\cdot d}\bigr)
\end{aligned}}{%
    \FormulaRefAuto{UniqueExistenceUnderEquivalentConjuncts}[thm]{17,16}}
\end{tabproofwide}

\FormulaDefDeltaK[Addition rationaler Zahlen]{%
  q,r\in\mathbb Q\vdash
  q+r\coloneqq\iota u\Bigl(
  \begin{aligned}[t]
    &u\in\mathbb Q\land
      \exists a,c\in\mathbb Z\,\exists b,d\in\PositiveIntegers\,\bigl({}\\
    &q=\RatClass{a}{b}\land r=\RatClass{c}{d}\land
      u=\RatClass{a\cdot d+c\cdot b}{b\cdot d}\bigr)
  \end{aligned}
  \Bigr)%
}{RationalAdditionDef}{%
  \DeltaRow{Mengen}{q\dsep r\dsep u\dsep a\dsep b\dsep c\dsep d}%
  \DeltaRow{Neue Symbole}{+}%
}

\FormulaThmDeltaK[Addition auf Bruchklassen]{%
  a,c\in\mathbb Z\dsep b,d\in\PositiveIntegers\vdash
  \RatClass{a}{b}+\RatClass{c}{d}
  =\RatClass{a\cdot d+c\cdot b}{b\cdot d}%
}{RationalAdditionClassValue}{%
  \DeltaRow{Mengen}{a\dsep b\dsep c\dsep d}%
}
\begin{tabproofwide}
  \proofstepwidestarbreakkeep[1]{a,c\in\mathbb Z}{\rA}
  \proofstepwidestarbreakkeep[2]{b,d\in\PositiveIntegers}{\rA}
  \proofstepwidestarbreakkeep[1,2]{\RatClass{a}{b},\RatClass{c}{d}\in\mathbb Q}{%
    \rAIRepeat{2}{\FormulaRefAuto{RationalClassInQ}[thm]}{1,2}}
  \proofstepwidestar[1,2]{a d\in\mathbb Z}{%
      \FormulaRefAuto{IntMulClosure}[thm]{1,2}}
  \proofstepwidestar[1,2]{c b\in\mathbb Z}{%
      \FormulaRefAuto{IntMulClosure}[thm]{1,2}}
  \proofstepwidestar[2]{b d\in\PositiveIntegers}{%
      \FormulaRefAuto{PositiveIntegerMulClosure}[thm]{2}}
  \proofstepwidestar[1,2]{a d+c b\in\mathbb Z}{%
      \FormulaRefAuto{IntAddClosure}[thm]{%
        4,%
        5}}
  \proofstepwidestarbreakkeep[1,2]{\RatClass{a d+c b}{b d}\in\mathbb Q}{%
    \FormulaRefAuto{RationalClassInQ}[thm]{%
      7,%
      6}}
  \proofstepwidestarbreakkeep[1,2]{\exists!u\in\mathbb Q\,
    \exists x,z\in\mathbb Z\,\exists y,t\in\PositiveIntegers\,
    \begin{aligned}[t]
      (&\RatClass{a}{b}=\RatClass{x}{y}\land
       \RatClass{c}{d}=\RatClass{z}{t}\land{}\\[-2pt]
       &u=\RatClass{x t+z y}{y t})
    \end{aligned}}{%
    \FormulaRefAuto{RationalAdditionQuotientDescent}[thm]{3}}
  \proofstepwidestarbreakkeep[1,2]{\RatClass{a}{b}+\RatClass{c}{d}
    =\RatClass{a d+c b}{b d}}{%
    \FormulaRefAuto{RationalAdditionDef}[def]{1,2,3,8,9}}
\end{tabproofwide}

\FormulaThmDeltaK[Abgeschlossenheit der rationalen Addition]{%
  q,r\in\mathbb Q\vdash q+r\in\mathbb Q%
}{RationalAddClosure}{%
  \DeltaRow{Mengen}{q\dsep r}%
}
\begin{tabproofwide}
  \proofstepwidestarbreakkeep[1]{q,r\in\mathbb Q}{\rA}
  \proofstepwidestarbreakkeep[1]{\exists a\in\mathbb Z\,
    \exists b\in\PositiveIntegers\,q=\RatClass{a}{b}}{%
    \FormulaRefAuto{RationalPairRepresentative}[thm]{\rAEa{1}}}
  \proofstepwidestarbreakkeep[1]{\exists c\in\mathbb Z\,
    \exists d\in\PositiveIntegers\,r=\RatClass{c}{d}}{%
    \FormulaRefAuto{RationalPairRepresentative}[thm]{\rAEb{1}}}
  \proofstepwidestarbreakkeep[4]{a\in\mathbb Z\land b\in\PositiveIntegers
    \land q=\RatClass{a}{b}}{\rA}
  \proofstepwidestarbreakkeep[5]{c\in\mathbb Z\land d\in\PositiveIntegers
    \land r=\RatClass{c}{d}}{\rA}
  \proofstepwidestar[4,5]{\RatClass{a}{b}+\RatClass{c}{d}=\RatClass{a d+c b}{b d}}{%
      \FormulaRefAuto{RationalAdditionClassValue}[thm]{4,5}}
  \proofstepwidestarbreakkeep[4,5]{q+r=\RatClass{a d+c b}{b d}}{%
    \rIE{\rAEb{\rAEb{4}},\rAEb{\rAEb{5}},%
      6}}
  \proofstepwidestarbreakkeep[4,5]{a d\in\mathbb Z}{%
    \FormulaRefAuto{IntMulClosure}[thm]{4,5}}
  \proofstepwidestarbreakkeep[4,5]{c b\in\mathbb Z}{%
    \FormulaRefAuto{IntMulClosure}[thm]{4,5}}
  \proofstepwidestarbreakkeep[4,5]{a d+c b\in\mathbb Z}{%
    \FormulaRefAuto{IntAddClosure}[thm]{8,9}}
  \proofstepwidestarbreakkeep[4,5]{b d\in\PositiveIntegers}{%
    \FormulaRefAuto{PositiveIntegerMulClosure}[thm]{4,5}}
  \proofstepwidestarbreakkeep[4,5]{\RatClass{a d+c b}{b d}\in\mathbb Q}{%
    \FormulaRefAuto{RationalClassInQ}[thm]{10,11}}
  \proofstepwidestarbreakkeep[4,5]{q+r\in\mathbb Q}{\rIE{7,12}}
  \proofstepwidestarbreakkeep[1]{q+r\in\mathbb Q}{%
    \rEE{2,4,\rEE{3,5,13}}}
\end{tabproofwide}

\section{Multiplikation}

\paragraph{Metasprachliche Motivation der Multiplikation.}
Die aus der gewohnten Bruchrechnung bekannte Vorschrift
\[
  \frac{a}{b}\cdot\frac{c}{d}=\frac{a\cdot c}{b\cdot d}
\]
legt nahe, zwei Bruchklassen durch Multiplikation ihrer Zähler und ihrer
Nenner zu verknüpfen. Auch diese Schreibweise ist zunächst nur eine
Motivation: Eine rationale Zahl ist eine Äquivalenzklasse und besitzt im
Allgemeinen viele Bruchpaarrepräsentanten. Die angegebene Vorschrift
definiert deshalb erst dann eine Operation auf \(\mathbb Q\), wenn ihr Wert
von der Wahl beider Repräsentanten unabhängig ist. Genau dies beweist der
folgende Satz; der anschließende Quotientenabstieg liefert daraus den
eindeutigen Wert auf den Bruchklassen, bevor die Multiplikation formal
definiert wird.

\FormulaThmDeltaK[Repräsentantenunabhängigkeit der rationalen Multiplikation]{%
  \begin{aligned}[t]
    &a,a',c,c'\in\mathbb Z\dsep
      b,b',d,d'\in\PositiveIntegers\dsep
      \RatClass{a}{b}=\RatClass{a'}{b'}\dsep{}\\
    &\RatClass{c}{d}=\RatClass{c'}{d'}\vdash
      \RatClass{a\cdot c}{b\cdot d}
      =\RatClass{a'\cdot c'}{b'\cdot d'}
  \end{aligned}%
}{RationalMultiplicationCompatible}{%
  \DeltaRow{Mengen}{a\dsep a'\dsep b\dsep b'\dsep
    c\dsep c'\dsep d\dsep d'}%
}

\begin{tabproofwide}
  \proofstepwidestarbreakkeep[1]{a,a',c,c'\in\mathbb Z}{\rA}
  \proofstepwidestarbreakkeep[2]{b,b',d,d'\in\PositiveIntegers}{\rA}
  \proofstepwidestarbreakkeep[3]{\RatClass{a}{b}=\RatClass{a'}{b'}}{\rA}
  \proofstepwidestarbreakkeep[4]{\RatClass{c}{d}=\RatClass{c'}{d'}}{\rA}
  \proofstepwidestarbreakkeep[1,2,3]{a\cdot b'=a'\cdot b}{%
    \FormulaRefAuto{RationalClassEquality}[thm]{1,2,3}}
  \proofstepwidestarbreakkeep[1,2,4]{c\cdot d'=c'\cdot d}{%
    \FormulaRefAuto{RationalClassEquality}[thm]{1,2,4}}
  \proofstepwidestarbreakkeep[1]{a\cdot c,a'\cdot c'\in\mathbb Z}{%
    \rAIRepeat{2}{\FormulaRefAuto{IntMulClosure}[thm]}{1}}
  \proofstepwidestarbreakkeep[2]{b\cdot d,b'\cdot d'\in\PositiveIntegers}{%
    \rAIRepeat{2}{\FormulaRefAuto{PositiveIntegerMulClosure}[thm]}{2}}
  \proofstepwidestarbreakkeep[1,2,3,4]{%
    (a\cdot c)\cdot(b'\cdot d')
    =(a'\cdot c')\cdot(b\cdot d)}{%
    \rIE{5,6,\FormulaRefAuto{IntMulAssociative}[thm],%
      \FormulaRefAuto{IntMulCommutative}[thm]}}
  \proofstepwidestarbreakkeep[1,2,3,4]{\RatClass{a\cdot c}{b\cdot d}
    =\RatClass{a'\cdot c'}{b'\cdot d'}}{%
    \FormulaRefAuto{RationalClassEquality}[thm]{7,8,9}}
\end{tabproofwide}

\FormulaThmDeltaK[Quotientenabstieg der rationalen Multiplikation]{%
  \begin{aligned}[t]
    &q,r\in\mathbb Q\vdash
      \exists!u\in\mathbb Q\,
      \exists a,c\in\mathbb Z\,\exists b,d\in\PositiveIntegers\,\Bigl({}\\[-2pt]
    &\qquad q=\RatClass{a}{b}\ \land\\[-2pt]
    &\qquad r=\RatClass{c}{d}\ \land\\[-2pt]
    &\qquad u=\RatClass{a\cdot c}{b\cdot d}\Bigr)
  \end{aligned}%
}{RationalMultiplicationQuotientDescent}{%
  \DeltaRow{Mengen}{q\dsep r\dsep u\dsep v\dsep
    a\dsep a'\dsep b\dsep b'\dsep c\dsep c'\dsep d\dsep d'}%
  \DeltaRow{Weitere Mengen}{e\dsep f\dsep g\dsep h\dsep z}%
}

\begin{tabproofwide}
  \proofstepwidestar[1]{q,r\in\mathbb Q}{%
    \rA}
  \proofstepwidestar[]{\forall a\forall b\,\bigl((a\in\mathbb Z\land b\in\PositiveIntegers)\rightarrow\RatClass{a}{b}\in\mathbb Q\bigr)}{%
    \FormulaRefAuto{RationalPairClassTyping}[thm]}
  \proofstepwidestar[]{\begin{aligned}[t]
&\forall z\in\mathbb Q\,\exists a\exists b\,\\[-2pt]
&\bigl((a\in\mathbb Z\land b\in\PositiveIntegers)\land z=\RatClass{a}{b}\bigr)
\end{aligned}}{%
    \FormulaRefAuto{RationalPairRepresentationSurjective}[thm]}
  \proofstepwidestar[4]{(a\in\mathbb Z\land b\in\PositiveIntegers)\land(c\in\mathbb Z\land d\in\PositiveIntegers)}{%
    \rA}
  \proofstepwidestar[4]{a\cdot c\in\mathbb Z}{%
    \FormulaRefAuto{IntMulClosure}[thm]{\rAEa{\rAEa{4}},\rAEa{\rAEb{4}}}}
  \proofstepwidestar[4]{b\cdot d\in\PositiveIntegers}{%
    \FormulaRefAuto{PositiveIntegerMulClosure}[thm]{\rAEb{\rAEa{4}},\rAEb{\rAEb{4}}}}
  \proofstepwidestar[4]{\RatClass{a\cdot c}{b\cdot d}\in\mathbb Q}{%
    \FormulaRefAuto{RationalClassInQ}[thm]{5,6}}
  \proofstepwidestar[]{\begin{aligned}[t]
&\forall a\forall b\forall c\forall d\\[-2pt]&\bigl(((a\in\mathbb Z\land b\in\PositiveIntegers)\land(c\in\mathbb Z\land d\in\PositiveIntegers))\rightarrow\\[-2pt]
&\qquad \RatClass{a\cdot c}{b\cdot d}\in\mathbb Q\bigr)
\end{aligned}}{%
    \rUI{\rUI{\rUI{\rUI{\rRI{4,7}}}}}}
  \proofstepwidestar[9]{\begin{aligned}[t]
&(a\in\mathbb Z\land b\in\PositiveIntegers)\land(c\in\mathbb Z\land d\in\PositiveIntegers)\land{}\\[-2pt]
&(e\in\mathbb Z\land f\in\PositiveIntegers)\land(g\in\mathbb Z\land h\in\PositiveIntegers)\land{}\\[-2pt]
&\RatClass{a}{b}=\RatClass{e}{f}\land \RatClass{c}{d}=\RatClass{g}{h}
\end{aligned}}{%
    \rA}
  \proofstepwidestar[9]{\RatClass{a\cdot c}{b\cdot d}=\RatClass{e\cdot g}{f\cdot h}}{%
    \FormulaRefAuto{RationalMultiplicationCompatible}[thm]{\rAIStar{\rAEa{\rAEa{9}},\rAEa{\rAEa{\rAEb{\rAEb{9}}}},\rAEa{\rAEa{\rAEb{9}}},\rAEa{\rAEa{\rAEb{\rAEb{\rAEb{9}}}}}},\rAIStar{\rAEb{\rAEa{9}},\rAEb{\rAEa{\rAEb{\rAEb{9}}}},\rAEb{\rAEa{\rAEb{9}}},\rAEb{\rAEa{\rAEb{\rAEb{\rAEb{9}}}}}},\rAEa{\rAEb{\rAEb{\rAEb{\rAEb{9}}}}},\rAEb{\rAEb{\rAEb{\rAEb{\rAEb{9}}}}}}}
  \proofstepwidestar[]{\begin{aligned}[t]
&\forall a\forall b\forall c\forall d\forall e\forall f\forall g\forall h\,\Bigl(\\[-2pt]
&(a\in\mathbb Z\land b\in\PositiveIntegers)\land(c\in\mathbb Z\land d\in\PositiveIntegers)\land{}\\[-2pt]
&(e\in\mathbb Z\land f\in\PositiveIntegers)\land(g\in\mathbb Z\land h\in\PositiveIntegers)\land{}\\[-2pt]
&\RatClass{a}{b}=\RatClass{e}{f}\land \RatClass{c}{d}=\RatClass{g}{h}\rightarrow\\[-2pt]
&\qquad \RatClass{a\cdot c}{b\cdot d}=\RatClass{e\cdot g}{f\cdot h}\Bigr)
\end{aligned}}{%
    \rUI{\rUI{\rUI{\rUI{\rUI{\rUI{\rUI{\rUI{\rRI{9,10}}}}}}}}}}
  \proofstepwidestar[1]{\begin{aligned}[t]
&\exists!u\in\mathbb Q\,\exists a\exists b\exists c\exists d\,\\[-2pt]
&\bigl((a\in\mathbb Z\land b\in\PositiveIntegers)\land(c\in\mathbb Z\land d\in\PositiveIntegers)\\[-2pt]&\land q=\RatClass{a}{b}\land{}\\[-2pt]
&\qquad r=\RatClass{c}{d}\land u=\RatClass{a\cdot c}{b\cdot d}\bigr)
\end{aligned}}{%
    \FormulaRefAuto{RepresentativePairBinaryValueUnique}[thm]{2,3,8,11,\rAEa{1},\rAEb{1}}}
  \proofstepwidestar[]{\begin{aligned}[t]
&\forall u\,\Bigl(\exists a\exists b\exists c\exists d\,\\[-2pt]
&\bigl((a\in\mathbb Z\land b\in\PositiveIntegers)\land(c\in\mathbb Z\land d\in\PositiveIntegers)\\[-2pt]&\land q=\RatClass{a}{b}\land{}\\[-2pt]
&\qquad r=\RatClass{c}{d}\land u=\RatClass{a\cdot c}{b\cdot d}\bigr)\\[-2pt]
&\quad\leftrightarrow{}\exists a,c\in\mathbb Z\,\exists b,d\in\PositiveIntegers\,\\[-2pt]
&\bigl(q=\RatClass{a}{b}\land r=\RatClass{c}{d}\land{}\\[-2pt]
&\qquad u=\RatClass{a\cdot c}{b\cdot d}\bigr)\Bigr)
\end{aligned}}{%
    \rUI{\FormulaRefAuto{P \leftrightarrow Q \vdash Q \leftrightarrow P}[thm]{\FormulaRefAuto{BoundedFourExistsGroupedGuard}[thm]}}}
  \proofstepwidestar[1]{\begin{aligned}[t]
&\exists!u\in\mathbb Q\,\exists a,c\in\mathbb Z\,\exists b,d\in\PositiveIntegers\,\\[-2pt]
&\bigl(q=\RatClass{a}{b}\land r=\RatClass{c}{d}\land{}\\[-2pt]
&\qquad u=\RatClass{a\cdot c}{b\cdot d}\bigr)
\end{aligned}}{%
    \FormulaRefAuto{UniqueExistenceUnderEquivalentConjuncts}[thm]{13,12}}
\end{tabproofwide}

\Needspace{15\baselineskip}
\FormulaDefDeltaK[Multiplikation rationaler Zahlen]{%
  q,r\in\mathbb Q\vdash
  q\cdot r\coloneqq\iota u\Bigl(
  \begin{aligned}[t]
    &u\in\mathbb Q\land
      \exists a,c\in\mathbb Z\,\exists b,d\in\PositiveIntegers\,\bigl({}\\
    &q=\RatClass{a}{b}\land r=\RatClass{c}{d}\land
      u=\RatClass{a\cdot c}{b\cdot d}\bigr)
  \end{aligned}
  \Bigr)%
}{RationalMultiplicationDef}{%
  \DeltaRow{Mengen}{q\dsep r\dsep u\dsep a\dsep b\dsep c\dsep d}%
  \DeltaRow{Neue Symbole}{\cdot}%
}

\Needspace{9\baselineskip}
\FormulaThmDeltaK[Multiplikation auf Bruchklassen]{%
  a,c\in\mathbb Z\dsep b,d\in\PositiveIntegers\vdash
  \RatClass{a}{b}\cdot\RatClass{c}{d}
  =\RatClass{a\cdot c}{b\cdot d}%
}{RationalMultiplicationClassValue}{%
  \DeltaRow{Mengen}{a\dsep b\dsep c\dsep d}%
}
\begin{tabproofwide}
  \proofstepwidestarbreakkeep[1]{a,c\in\mathbb Z}{\rA}
  \proofstepwidestarbreakkeep[2]{b,d\in\PositiveIntegers}{\rA}
  \proofstepwidestarbreakkeep[1,2]{\RatClass{a}{b},\RatClass{c}{d}\in\mathbb Q}{%
    \rAIRepeat{2}{\FormulaRefAuto{RationalClassInQ}[thm]}{1,2}}
  \proofstepwidestar[1]{a c\in\mathbb Z}{%
      \FormulaRefAuto{IntMulClosure}[thm]{1}}
  \proofstepwidestar[2]{b d\in\PositiveIntegers}{%
      \FormulaRefAuto{PositiveIntegerMulClosure}[thm]{2}}
  \proofstepwidestarbreakkeep[1,2]{\RatClass{a c}{b d}\in\mathbb Q}{%
    \FormulaRefAuto{RationalClassInQ}[thm]{%
      4,%
      5}}
  \proofstepwidestarbreakkeep[1,2]{\begin{aligned}[t]
    &\exists!u\in\mathbb Q\,
      \exists x,z\in\mathbb Z\,\exists y,t\in\PositiveIntegers\,{}\\
    &\bigl(\RatClass{a}{b}=\RatClass{x}{y}
      \land\RatClass{c}{d}=\RatClass{z}{t}\land{}\\
    &\hspace{8em}u=\RatClass{x z}{y t}\bigr)
    \end{aligned}}{%
    \FormulaRefAuto{RationalMultiplicationQuotientDescent}[thm]{3}}
  \proofstepwidestarbreakkeep[1,2]{\RatClass{a}{b}\cdot\RatClass{c}{d}
    =\RatClass{a c}{b d}}{%
    \FormulaRefAuto{RationalMultiplicationDef}[def]{1,2,3,6,7}}
\end{tabproofwide}

\FormulaThmDeltaK[Abgeschlossenheit der rationalen Multiplikation]{%
  q,r\in\mathbb Q\vdash q\cdot r\in\mathbb Q%
}{RationalMulClosure}{%
  \DeltaRow{Mengen}{q\dsep r}%
}
\begin{tabproofwide}
  \proofstepwidestarbreakkeep[1]{q,r\in\mathbb Q}{\rA}
  \proofstepwidestarbreakkeep[1]{\exists a\in\mathbb Z\,
    \exists b\in\PositiveIntegers\,q=\RatClass{a}{b}}{%
    \FormulaRefAuto{RationalPairRepresentative}[thm]{\rAEa{1}}}
  \proofstepwidestarbreakkeep[1]{\exists c\in\mathbb Z\,
    \exists d\in\PositiveIntegers\,r=\RatClass{c}{d}}{%
    \FormulaRefAuto{RationalPairRepresentative}[thm]{\rAEb{1}}}
  \proofstepwidestarbreakkeep[4]{a\in\mathbb Z\land b\in\PositiveIntegers
    \land q=\RatClass{a}{b}}{\rA}
  \proofstepwidestarbreakkeep[5]{c\in\mathbb Z\land d\in\PositiveIntegers
    \land r=\RatClass{c}{d}}{\rA}
  \proofstepwidestar[4,5]{\RatClass{a}{b}\cdot\RatClass{c}{d}=\RatClass{a c}{b d}}{%
      \FormulaRefAuto{RationalMultiplicationClassValue}[thm]{4,5}}
  \proofstepwidestarbreakkeep[4,5]{q\cdot r=\RatClass{a c}{b d}}{%
    \rIE{\rAEb{\rAEb{4}},\rAEb{\rAEb{5}},%
      6}}
  \proofstepwidestar[4,5]{a c\in\mathbb Z}{%
      \FormulaRefAuto{IntMulClosure}[thm]{4,5}}
  \proofstepwidestar[4,5]{b d\in\PositiveIntegers}{%
      \FormulaRefAuto{PositiveIntegerMulClosure}[thm]{4,5}}
  \proofstepwidestarbreakkeep[4,5]{\RatClass{a c}{b d}\in\mathbb Q}{%
    \FormulaRefAuto{RationalClassInQ}[thm]{%
      8,%
      9}}
  \proofstepwidestarbreakkeep[4,5]{q\cdot r\in\mathbb Q}{\rIE{7,10}}
  \proofstepwidestarbreakkeep[1]{q\cdot r\in\mathbb Q}{%
    \rEE{2,4,\rEE{3,5,11}}}
\end{tabproofwide}

\section{Grundlegende Rechengesetze}

\FormulaThmDeltaK[Assoziativität der rationalen Addition]{%
  p,q,r\in\mathbb Q\vdash(p+q)+r=p+(q+r)%
}{RationalAddAssociative}{%
  \DeltaRow{Mengen}{p\dsep q\dsep r}%
}
\begin{tabproofwide}
  \proofstepwidestarbreakkeep[1]{p,q,r\in\mathbb Q}{\rA}
  \proofstepwidestarbreakkeep[1]{\exists a\in\mathbb Z\,
    \exists b\in\PositiveIntegers\,p=\RatClass{a}{b}}{%
    \FormulaRefAuto{RationalPairRepresentative}[thm]{\rAEa{1}}}
  \proofstepwidestarbreakkeep[1]{\exists c\in\mathbb Z\,
    \exists d\in\PositiveIntegers\,q=\RatClass{c}{d}}{%
    \FormulaRefAuto{RationalPairRepresentative}[thm]{\rAEa{\rAEb{1}}}}
  \proofstepwidestarbreakkeep[1]{\exists e\in\mathbb Z\,
    \exists f\in\PositiveIntegers\,r=\RatClass{e}{f}}{%
    \FormulaRefAuto{RationalPairRepresentative}[thm]{\rAEb{\rAEb{1}}}}
  \proofstepwidestarbreakkeep[5]{a,c,e\in\mathbb Z\land
    b,d,f\in\PositiveIntegers\land p=\RatClass{a}{b}\land
    q=\RatClass{c}{d}\land r=\RatClass{e}{f}}{\rA}
  \proofstepwidestarbreakkeep[5]{\RatClass{a}{b}+\RatClass{c}{d}
    =\RatClass{a d+c b}{b d}}{%
    \FormulaRefAuto{RationalAdditionClassValue}[thm]{5}}
  \proofstepwidestar[5]{\RatClass{a d+c b}{b d}+\RatClass{e}{f}=\RatClass{(a d+c b)f+e(b d)}{(b d)f}}{%
      \FormulaRefAuto{RationalAdditionClassValue}[thm]{5,6}}
  \proofstepwidestarbreakkeep[5]{%
    (\RatClass{a}{b}+\RatClass{c}{d})+\RatClass{e}{f}
    =\RatClass{(a d+c b)f+e(b d)}{(b d)f}}{%
    \rIE{6,7}}
  \proofstepwidestarbreakkeep[5]{\RatClass{c}{d}+\RatClass{e}{f}
    =\RatClass{c f+e d}{d f}}{%
    \FormulaRefAuto{RationalAdditionClassValue}[thm]{5}}
  \proofstepwidestar[5]{\RatClass{a}{b}+\RatClass{c f+e d}{d f}=\RatClass{a(d f)+(c f+e d)b}{b(d f)}}{%
      \FormulaRefAuto{RationalAdditionClassValue}[thm]{5,9}}
  \proofstepwidestarbreakkeep[5]{%
    \RatClass{a}{b}+(\RatClass{c}{d}+\RatClass{e}{f})
    =\RatClass{a(d f)+(c f+e d)b}{b(d f)}}{%
    \rIE{9,10}}
  \proofstepwidestarbreakkeep[5]{%
    (a d+c b)f+e(b d)=a(d f)+(c f+e d)b}{%
    \FormulaRefAuto{IntAddAssociative}[thm]{5},%
    \FormulaRefAuto{IntAddCommutative}[thm]{5},%
    \FormulaRefAuto{IntMulAssociative}[thm]{5},%
    \FormulaRefAuto{IntMulCommutative}[thm]{5},%
    \FormulaRefAuto{IntMulLeftDistributive}[thm]{5},%
    \FormulaRefAuto{IntMulRightDistributive}[thm]{5}}
  \proofstepwidestarbreakkeep[5]{(b d)f=b(d f)}{%
    \FormulaRefAuto{IntMulAssociative}[thm]{5}}
  \proofstepwidestarbreakkeep[5]{%
    (\RatClass{a}{b}+\RatClass{c}{d})+\RatClass{e}{f}
    =\RatClass{a}{b}+(\RatClass{c}{d}+\RatClass{e}{f})}{%
    \rIE{8,11,12,13}}
  \proofstepwidestarbreakkeep[5]{(p+q)+r=p+(q+r)}{\rIE{5,14}}
  \proofstepwidestarbreakkeep[1]{(p+q)+r=p+(q+r)}{%
    \rEE{2,3,4,5,15}}
\end{tabproofwide}

\FormulaThmDeltaK[Kommutativität der rationalen Addition]{%
  q,r\in\mathbb Q\vdash q+r=r+q%
}{RationalAddCommutative}{%
  \DeltaRow{Mengen}{q\dsep r}%
}
\begin{tabproofwide}
  \proofstepwidestarbreakkeep[1]{q,r\in\mathbb Q}{\rA}
  \proofstepwidestarbreakkeep[1]{\exists a\in\mathbb Z\,
    \exists b\in\PositiveIntegers\,q=\RatClass{a}{b}}{%
    \FormulaRefAuto{RationalPairRepresentative}[thm]{\rAEa{1}}}
  \proofstepwidestarbreakkeep[1]{\exists c\in\mathbb Z\,
    \exists d\in\PositiveIntegers\,r=\RatClass{c}{d}}{%
    \FormulaRefAuto{RationalPairRepresentative}[thm]{\rAEb{1}}}
  \proofstepwidestarbreakkeep[4]{a,c\in\mathbb Z\land b,d\in\PositiveIntegers
    \land q=\RatClass{a}{b}\land r=\RatClass{c}{d}}{\rA}
  \proofstepwidestarbreakkeep[4]{\RatClass{a}{b}+\RatClass{c}{d}
    =\RatClass{c}{d}+\RatClass{a}{b}}{%
    \FormulaRefAuto{RationalAdditionClassValue}[thm]{4},%
    \FormulaRefAuto{IntAddCommutative}[thm],%
    \FormulaRefAuto{IntMulCommutative}[thm]}
  \proofstepwidestarbreakkeep[4]{q+r=r+q}{\rIE{4,5}}
  \proofstepwidestarbreakkeep[1]{q+r=r+q}{\rEE{2,3,4,6}}
\end{tabproofwide}

\FormulaThmDeltaK[Nullgesetz der rationalen Addition]{%
  q\in\mathbb Q\vdash q+\RatZero=q%
}{RationalAddZero}{%
  \DeltaRow{Mengen}{q}%
}
\begin{tabproofwide}
  \proofstepwidestarbreakkeep[1]{q\in\mathbb Q}{\rA}
  \proofstepwidestarbreakkeep[1]{\exists a\in\mathbb Z\,
    \exists b\in\PositiveIntegers\,q=\RatClass{a}{b}}{%
    \FormulaRefAuto{RationalPairRepresentative}[thm]{1}}
  \proofstepwidestarbreakkeep[3]{a\in\mathbb Z\land b\in\PositiveIntegers
    \land q=\RatClass{a}{b}}{\rA}
  \proofstepwidestarbreakkeep[3]{\RatClass{a}{b}+\RatZero
    =\RatClass{a}{b}}{%
    \FormulaRefAuto{RationalAdditionClassValue}[thm]{3},%
    \FormulaRefAuto{RationalZeroOneClasses}[thm],%
    \FormulaRefAuto{IntAddZero}[thm],%
    \FormulaRefAuto{IntMulOne}[thm],%
    \FormulaRefAuto{IntegerMulZeroFromAxioms}[thm]}
  \proofstepwidestarbreakkeep[3]{q+\RatZero=q}{\rIE{3,4}}
  \proofstepwidestarbreakkeep[1]{q+\RatZero=q}{\rEE{2,3,5}}
\end{tabproofwide}

\FormulaThmDeltaK[Additives Inverses rationaler Zahlen]{%
  q\in\mathbb Q\vdash q+(-q)=\RatZero%
}{RationalAddInverse}{%
  \DeltaRow{Mengen}{q}%
}
\begin{tabproofwide}
  \proofstepwidestarbreakkeep[1]{q\in\mathbb Q}{\rA}
  \proofstepwidestarbreakkeep[1]{\exists a\in\mathbb Z\,
    \exists b\in\PositiveIntegers\,q=\RatClass{a}{b}}{%
    \FormulaRefAuto{RationalPairRepresentative}[thm]{1}}
  \proofstepwidestarbreakkeep[3]{a\in\mathbb Z\land b\in\PositiveIntegers
    \land q=\RatClass{a}{b}}{\rA}
  \proofstepwidestar[3]{-\RatClass{a}{b}=\RatClass{-a}{b}}{%
      \FormulaRefAuto{RationalNegationClassValue}[thm]{3}}
  \proofstepwidestarbreakkeep[3]{-q=\RatClass{-a}{b}}{%
    \rIE{\rAEb{\rAEb{3}},%
      4}}
  \proofstepwidestarbreakkeep[3]{\RatClass{a}{b}+\RatClass{-a}{b}
    =\RatZero}{%
    \FormulaRefAuto{RationalAdditionClassValue}[thm]{3},%
    \FormulaRefAuto{RationalZeroOneClasses}[thm],%
    \FormulaRefAuto{IntAddInverse}[thm],%
    \FormulaRefAuto{IntMulOne}[thm],%
    \FormulaRefAuto{IntegerMulZeroFromAxioms}[thm]}
  \proofstepwidestarbreakkeep[3]{q+(-q)=\RatZero}{\rIE{3,5,6}}
  \proofstepwidestarbreakkeep[1]{q+(-q)=\RatZero}{\rEE{2,3,7}}
\end{tabproofwide}

\FormulaThmDeltaK[Assoziativität der rationalen Multiplikation]{%
  p,q,r\in\mathbb Q\vdash(p\cdot q)\cdot r=p\cdot(q\cdot r)%
}{RationalMulAssociative}{%
  \DeltaRow{Mengen}{p\dsep q\dsep r}%
}
\begin{tabproofwide}
  \proofstepwidestarbreakkeep[1]{p,q,r\in\mathbb Q}{\rA}
  \proofstepwidestarbreakkeep[1]{\exists a\in\mathbb Z\,
    \exists b\in\PositiveIntegers\,p=\RatClass{a}{b}}{%
    \FormulaRefAuto{RationalPairRepresentative}[thm]{\rAEa{1}}}
  \proofstepwidestarbreakkeep[1]{\exists c\in\mathbb Z\,
    \exists d\in\PositiveIntegers\,q=\RatClass{c}{d}}{%
    \FormulaRefAuto{RationalPairRepresentative}[thm]{\rAEa{\rAEb{1}}}}
  \proofstepwidestarbreakkeep[1]{\exists e\in\mathbb Z\,
    \exists f\in\PositiveIntegers\,r=\RatClass{e}{f}}{%
    \FormulaRefAuto{RationalPairRepresentative}[thm]{\rAEb{\rAEb{1}}}}
  \proofstepwidestarbreakkeep[5]{a,c,e\in\mathbb Z\land
    b,d,f\in\PositiveIntegers\land p=\RatClass{a}{b}\land
    q=\RatClass{c}{d}\land r=\RatClass{e}{f}}{\rA}
  \proofstepwidestarbreakkeep[5]{%
    (\RatClass{a}{b}\cdot\RatClass{c}{d})\cdot\RatClass{e}{f}
    =\RatClass{a}{b}\cdot(\RatClass{c}{d}\cdot\RatClass{e}{f})}{%
    \FormulaRefAuto{RationalMultiplicationClassValue}[thm]{5},%
    \FormulaRefAuto{IntMulAssociative}[thm]}
  \proofstepwidestarbreakkeep[5]{(p\cdot q)\cdot r=p\cdot(q\cdot r)}{\rIE{5,6}}
  \proofstepwidestarbreakkeep[1]{(p\cdot q)\cdot r=p\cdot(q\cdot r)}{%
    \rEE{2,3,4,5,7}}
\end{tabproofwide}

\FormulaThmDeltaK[Kommutativität der rationalen Multiplikation]{%
  q,r\in\mathbb Q\vdash q\cdot r=r\cdot q%
}{RationalMulCommutative}{%
  \DeltaRow{Mengen}{q\dsep r}%
}
\begin{tabproofwide}
  \proofstepwidestarbreakkeep[1]{q,r\in\mathbb Q}{\rA}
  \proofstepwidestarbreakkeep[1]{\exists a\in\mathbb Z\,
    \exists b\in\PositiveIntegers\,q=\RatClass{a}{b}}{%
    \FormulaRefAuto{RationalPairRepresentative}[thm]{\rAEa{1}}}
  \proofstepwidestarbreakkeep[1]{\exists c\in\mathbb Z\,
    \exists d\in\PositiveIntegers\,r=\RatClass{c}{d}}{%
    \FormulaRefAuto{RationalPairRepresentative}[thm]{\rAEb{1}}}
  \proofstepwidestarbreakkeep[4]{a,c\in\mathbb Z\land b,d\in\PositiveIntegers
    \land q=\RatClass{a}{b}\land r=\RatClass{c}{d}}{\rA}
  \proofstepwidestarbreakkeep[4]{\RatClass{a}{b}\cdot\RatClass{c}{d}
    =\RatClass{c}{d}\cdot\RatClass{a}{b}}{%
    \FormulaRefAuto{RationalMultiplicationClassValue}[thm]{4},%
    \FormulaRefAuto{IntMulCommutative}[thm]}
  \proofstepwidestarbreakkeep[4]{q\cdot r=r\cdot q}{\rIE{4,5}}
  \proofstepwidestarbreakkeep[1]{q\cdot r=r\cdot q}{\rEE{2,3,4,6}}
\end{tabproofwide}

\FormulaThmDeltaK[Einheitsgesetz der rationalen Multiplikation]{%
  q\in\mathbb Q\vdash q\cdot\RatOne=q%
}{RationalMulOne}{%
  \DeltaRow{Mengen}{q}%
}
\begin{tabproofwide}
  \proofstepwidestarbreakkeep[1]{q\in\mathbb Q}{\rA}
  \proofstepwidestarbreakkeep[1]{\exists a\in\mathbb Z\,
    \exists b\in\PositiveIntegers\,q=\RatClass{a}{b}}{%
    \FormulaRefAuto{RationalPairRepresentative}[thm]{1}}
  \proofstepwidestarbreakkeep[3]{a\in\mathbb Z\land b\in\PositiveIntegers
    \land q=\RatClass{a}{b}}{\rA}
  \proofstepwidestarbreakkeep[3]{\RatClass{a}{b}\cdot\RatOne
    =\RatClass{a}{b}}{%
    \FormulaRefAuto{RationalMultiplicationClassValue}[thm]{3},%
    \FormulaRefAuto{RationalZeroOneClasses}[thm],%
    \FormulaRefAuto{IntMulOne}[thm]}
  \proofstepwidestarbreakkeep[3]{q\cdot\RatOne=q}{\rIE{3,4}}
  \proofstepwidestarbreakkeep[1]{q\cdot\RatOne=q}{\rEE{2,3,5}}
\end{tabproofwide}

\FormulaThmDeltaK[Nullgesetz der rationalen Multiplikation]{%
  q\in\mathbb Q\vdash q\cdot\RatZero=\RatZero%
}{RationalMulZero}{%
  \DeltaRow{Mengen}{q}%
}
\begin{tabproofwide}
  \proofstepwidestarbreakkeep[1]{q\in\mathbb Q}{\rA}
  \proofstepwidestarbreakkeep[1]{\exists a\in\mathbb Z\,
    \exists b\in\PositiveIntegers\,q=\RatClass{a}{b}}{%
    \FormulaRefAuto{RationalPairRepresentative}[thm]{1}}
  \proofstepwidestarbreakkeep[3]{a\in\mathbb Z\land b\in\PositiveIntegers
    \land q=\RatClass{a}{b}}{\rA}
  \proofstepwidestarbreakkeep[3]{\RatClass{a}{b}\cdot\RatZero
    =\RatZero}{%
    \FormulaRefAuto{RationalMultiplicationClassValue}[thm]{3},%
    \FormulaRefAuto{RationalZeroOneClasses}[thm],%
    \FormulaRefAuto{IntegerMulZeroFromAxioms}[thm],%
    \FormulaRefAuto{IntMulOne}[thm]}
  \proofstepwidestarbreakkeep[3]{q\cdot\RatZero=\RatZero}{\rIE{3,4}}
  \proofstepwidestarbreakkeep[1]{q\cdot\RatZero=\RatZero}{\rEE{2,3,5}}
\end{tabproofwide}

\FormulaThmDeltaK[Distributivität rationaler Zahlen]{%
  p,q,r\in\mathbb Q\vdash
  p\cdot(q+r)=p\cdot q+p\cdot r%
}{RationalDistributive}{%
  \DeltaRow{Mengen}{p\dsep q\dsep r}%
}
\begin{tabproofwide}
  \proofstepwidestarbreakkeep[1]{p,q,r\in\mathbb Q}{\rA}
  \proofstepwidestarbreakkeep[1]{\exists a\in\mathbb Z\,
    \exists b\in\PositiveIntegers\,p=\RatClass{a}{b}}{%
    \FormulaRefAuto{RationalPairRepresentative}[thm]{\rAEa{1}}}
  \proofstepwidestarbreakkeep[1]{\exists c\in\mathbb Z\,
    \exists d\in\PositiveIntegers\,q=\RatClass{c}{d}}{%
    \FormulaRefAuto{RationalPairRepresentative}[thm]{\rAEa{\rAEb{1}}}}
  \proofstepwidestarbreakkeep[1]{\exists e\in\mathbb Z\,
    \exists f\in\PositiveIntegers\,r=\RatClass{e}{f}}{%
    \FormulaRefAuto{RationalPairRepresentative}[thm]{\rAEb{\rAEb{1}}}}
  \proofstepwidestarbreakkeep[5]{a,c,e\in\mathbb Z\land
    b,d,f\in\PositiveIntegers\land p=\RatClass{a}{b}\land
    q=\RatClass{c}{d}\land r=\RatClass{e}{f}}{\rA}
  \proofstepwidestarbreakkeep[5]{\RatClass{c}{d}+\RatClass{e}{f}
    =\RatClass{c f+e d}{d f}}{%
    \FormulaRefAuto{RationalAdditionClassValue}[thm]{5}}
  \proofstepwidestar[5]{\RatClass{a}{b}\cdot\RatClass{c f+e d}{d f}=\RatClass{a(c f+e d)}{b(d f)}}{%
      \FormulaRefAuto{RationalMultiplicationClassValue}[thm]{5,6}}
  \proofstepwidestarbreakkeep[5]{%
    \RatClass{a}{b}\cdot(\RatClass{c}{d}+\RatClass{e}{f})
    =\RatClass{a(c f+e d)}{b(d f)}}{%
    \rIE{6,7}}
  \proofstepwidestarbreakkeep[5]{\RatClass{a}{b}\cdot\RatClass{c}{d}
    =\RatClass{a c}{b d}}{%
    \FormulaRefAuto{RationalMultiplicationClassValue}[thm]{5}}
  \proofstepwidestarbreakkeep[5]{\RatClass{a}{b}\cdot\RatClass{e}{f}
    =\RatClass{a e}{b f}}{%
    \FormulaRefAuto{RationalMultiplicationClassValue}[thm]{5}}
  \proofstepwidestar[5]{\RatClass{a c}{b d}+\RatClass{a e}{b f}=\RatClass{(a c)(b f)+(a e)(b d)}{(b d)(b f)}}{%
      \FormulaRefAuto{RationalAdditionClassValue}[thm]{5,9,10}}
  \proofstepwidestarbreakkeep[5]{%
    \RatClass{a}{b}\cdot\RatClass{c}{d}
      +\RatClass{a}{b}\cdot\RatClass{e}{f}
    =\RatClass{(a c)(b f)+(a e)(b d)}{(b d)(b f)}}{%
    \rIE{9,10,11}}
  \proofstepwidestarbreakkeep[5]{%
    a(c f+e d),(a c)(b f)+(a e)(b d)\in\mathbb Z}{%
    \FormulaRefAuto{IntMulClosure}[thm]{5},%
    \FormulaRefAuto{IntAddClosure}[thm]{5}}
  \proofstepwidestarbreakkeep[5]{%
    b(d f),(b d)(b f)\in\PositiveIntegers}{%
    \FormulaRefAuto{PositiveIntegerMulClosure}[thm]{5}}
  \proofstepwidestarbreakkeep[5]{%
    a(c f+e d)\bigl((b d)(b f)\bigr)
    =\bigl((a c)(b f)+(a e)(b d)\bigr)b(d f)}{%
    \FormulaRefAuto{IntMulLeftDistributive}[thm]{5},%
    \FormulaRefAuto{IntMulRightDistributive}[thm]{5},%
    \FormulaRefAuto{IntMulAssociative}[thm]{5},%
    \FormulaRefAuto{IntMulCommutative}[thm]{5}}
  \proofstepwidestarbreakkeep[5]{%
    \RatClass{a(c f+e d)}{b(d f)}
    =\RatClass{(a c)(b f)+(a e)(b d)}{(b d)(b f)}}{%
    \FormulaRefAuto{RationalClassEquality}[thm]{13,14,15}}
  \proofstepwidestarbreakkeep[5]{%
    \RatClass{a}{b}\cdot(\RatClass{c}{d}+\RatClass{e}{f})
    =\RatClass{a}{b}\cdot\RatClass{c}{d}
      +\RatClass{a}{b}\cdot\RatClass{e}{f}}{\rIE{8,12,16}}
  \proofstepwidestarbreakkeep[5]{p\cdot(q+r)=p\cdot q+p\cdot r}{%
    \rIE{5,17}}
  \proofstepwidestarbreakkeep[1]{p\cdot(q+r)=p\cdot q+p\cdot r}{%
    \rEE{2,3,4,5,18}}
\end{tabproofwide}

\FormulaThmDeltaK[Additivität der ganzzahligen Einbettung]{%
  z,w\in\mathbb Z\vdash
  \RatEmbed(z+w)=\RatEmbed(z)+\RatEmbed(w)%
}{IntegerRationalEmbeddingAdditive}{%
  \DeltaRow{Mengen}{z\dsep w}%
  \DeltaRow{Funktion}{\RatEmbed}%
}
\begin{tabproofwide}
  \proofstepwidestarbreakkeep[1]{z,w\in\mathbb Z}{\rA}
  \proofstepwidestar[1]{z+w\in\mathbb Z}{%
      \FormulaRefAuto{IntegerAddClosureAxiom}[ax]{1}}
  \proofstepwidestarbreakkeep[1]{\RatEmbed(z+w)
    =\RatClass{z+w}{\IntOne}}{%
    \FormulaRefAuto{IntegerRationalEmbedding}[def]{%
      2}}
  \proofstepwidestar[]{\IntOne\in\mathbb Z}{%
      \FormulaRefAuto{IntegerOneCarrierAxiom}[ax]}
  \proofstepwidestar[]{\IntZero<\IntOne}{%
      \FormulaRefAuto{IntegerOnePositiveAxiom}[ax]}
  \proofstepwidestar[]{\IntOne\in\PositiveIntegers}{%
      \FormulaRefAuto{PositiveIntegersDef}[def]{\rAI{4,5}}}
  \proofstepwidestarbreakkeep[1]{\RatEmbed(z)+\RatEmbed(w)
    =\RatClass{z\cdot\IntOne+w\cdot\IntOne}
      {\IntOne\cdot\IntOne}}{%
    \FormulaRefAuto{IntegerRationalEmbedding}[def]{1},%
    \FormulaRefAuto{RationalAdditionClassValue}[thm]{%
      1,6}}
  \proofstepwidestarbreakkeep[1]{%
    \RatClass{z\cdot\IntOne+w\cdot\IntOne}
      {\IntOne\cdot\IntOne}
    =\RatClass{z+w}{\IntOne}}{%
    \FormulaRefAuto{IntegerMulOneAxiom}[ax]{1},%
    \FormulaRefAuto{IntegerAddCommutativeAxiom}[ax]{1}}
  \proofstepwidestarbreakkeep[1]{\RatEmbed(z+w)=\RatEmbed(z)+\RatEmbed(w)}{%
    \rIE{3,7,8}}
\end{tabproofwide}

\FormulaThmDeltaK[Multiplikativität der ganzzahligen Einbettung]{%
  z,w\in\mathbb Z\vdash
  \RatEmbed(z\cdot w)=\RatEmbed(z)\cdot\RatEmbed(w)%
}{IntegerRationalEmbeddingMultiplicative}{%
  \DeltaRow{Mengen}{z\dsep w}%
  \DeltaRow{Funktion}{\RatEmbed}%
}
\begin{tabproofwide}
  \proofstepwidestarbreakkeep[1]{z,w\in\mathbb Z}{\rA}
  \proofstepwidestar[1]{z\cdot w\in\mathbb Z}{%
      \FormulaRefAuto{IntegerMulClosureAxiom}[ax]{1}}
  \proofstepwidestarbreakkeep[1]{\RatEmbed(z\cdot w)
    =\RatClass{z\cdot w}{\IntOne}}{%
    \FormulaRefAuto{IntegerRationalEmbedding}[def]{%
      2}}
  \proofstepwidestar[]{\IntOne\in\mathbb Z}{%
      \FormulaRefAuto{IntegerOneCarrierAxiom}[ax]}
  \proofstepwidestar[]{\IntZero<\IntOne}{%
      \FormulaRefAuto{IntegerOnePositiveAxiom}[ax]}
  \proofstepwidestar[]{\IntOne\in\PositiveIntegers}{%
      \FormulaRefAuto{PositiveIntegersDef}[def]{\rAI{4,5}}}
  \proofstepwidestarbreakkeep[1]{\RatEmbed(z)\cdot\RatEmbed(w)
    =\RatClass{z\cdot w}{\IntOne\cdot\IntOne}}{%
    \FormulaRefAuto{IntegerRationalEmbedding}[def]{1},%
    \FormulaRefAuto{RationalMultiplicationClassValue}[thm]{%
      1,6}}
  \proofstepwidestar[1]{\IntOne \in \mathbb Z}{%
      \FormulaRefAuto{IntegerOneCarrierAxiom}[ax]}
  \proofstepwidestarbreakkeep[1]{\RatClass{z\cdot w}
      {\IntOne\cdot\IntOne}
    =\RatClass{z\cdot w}{\IntOne}}{%
    \FormulaRefAuto{IntegerMulOneAxiom}[ax]{%
      8}}
  \proofstepwidestarbreakkeep[1]{\RatEmbed(z\cdot w)
    =\RatEmbed(z)\cdot\RatEmbed(w)}{\rIE{3,7,9}}
\end{tabproofwide}

\section{Kehrwert und Division}

\FormulaThmDeltaK[Nullkriterium einer Bruchklasse]{%
  a\in\mathbb Z\dsep b\in\PositiveIntegers\vdash
  \bigl(\RatClass{a}{b}=\RatZero
  \leftrightarrow a=\IntZero\bigr)%
}{RationalClassZeroIffNumeratorZero}{%
  \DeltaRow{Mengen}{a\dsep b}%
}
\begin{tabproofsplitwide}
  \proofpartwideindR[Hinrichtung]{%
    a\in\mathbb Z\dsep b\in\PositiveIntegers\dsep
    \RatClass{a}{b}=\RatZero\vdash a=\IntZero
  }{%
    a\in\mathbb Z\dsep b\in\PositiveIntegers\dsep
    \RatClass{a}{b}=\RatZero\vdash a=\IntZero
  }[RationalClassZeroNumeratorForward]
    \proofstepwidestarbreakkeep[1]{a\in\mathbb Z}{\rA}
    \proofstepwidestarbreakkeep[2]{b\in\PositiveIntegers}{\rA}
    \proofstepwidestarbreakkeep[3]{\RatClass{a}{b}=\RatZero}{\rA}
    \proofstepwidestar[1,2,3]{\IntZero \in \mathbb Z}{%
      \FormulaRefAuto{IntegerZeroCarrierAxiom}[ax]}
    \proofstepwidestar[1,2,3]{\IntZero < \IntOne}{%
      \FormulaRefAuto{IntegerOnePositiveAxiom}[ax]}
    \proofstepwidestar[1,2,3]{\RatZero = \RatClass \IntZero \IntOne \land \RatOne = \RatClass \IntOne \IntOne}{%
      \FormulaRefAuto{RationalZeroOneClasses}[thm]}
    \proofstepwidestarbreakkeep[1,2,3]{a\cdot\IntOne
      =\IntZero\cdot b}{%
      \FormulaRefAuto{RationalClassEquality}[thm]{%
        1,4,2,%
        5,%
        \rIE{3,\rAEa{6}}}}
    \proofstepwidestar[1,2,3]{a \cdot \IntOne = a \land \IntOne \cdot a = a}{%
      \FormulaRefAuto{IntegerMulOneAxiom}[ax]{1}}
    \proofstepwidestar[1,2,3]{\IntZero \in \mathbb Z}{%
      \FormulaRefAuto{IntegerZeroCarrierAxiom}[ax]}
    \proofstepwidestar[2]{b\in\mathbb Z\land\IntZero<b}{%
      \FormulaRefAuto{PositiveIntegersDef}[def]{2}}
    \proofstepwidestar[1,2,3]{b\cdot\IntZero=\IntZero\land\IntZero\cdot b=\IntZero}{%
      \FormulaRefAuto{IntegerMulZeroFromAxioms}[thm]{%
          9,%
          \rAEa{10}}}
    \proofstepwidestarbreakkeep[1,2,3]{a=\IntZero}{%
      \rIE{7,8,%
        11}}
  \closeproofpartwideindR

  \proofpartwideindR[Rückrichtung]{%
    a\in\mathbb Z\dsep b\in\PositiveIntegers\dsep
    a=\IntZero\vdash\RatClass{a}{b}=\RatZero
  }{%
    a\in\mathbb Z\dsep b\in\PositiveIntegers\dsep
    a=\IntZero\vdash\RatClass{a}{b}=\RatZero
  }[RationalClassZeroNumeratorBackward]
    \proofstepwidestarbreakkeep[1]{a\in\mathbb Z}{\rA}
    \proofstepwidestarbreakkeep[2]{b\in\PositiveIntegers}{\rA}
    \proofstepwidestarbreakkeep[3]{a=\IntZero}{\rA}
    \proofstepwidestar[1,2,3]{a \cdot \IntOne = a \land \IntOne \cdot a = a}{%
      \FormulaRefAuto{IntegerMulOneAxiom}[ax]{1}}
    \proofstepwidestar[1,2,3]{\IntZero \in \mathbb Z}{%
      \FormulaRefAuto{IntegerZeroCarrierAxiom}[ax]}
    \proofstepwidestar[2]{b\in\mathbb Z\land\IntZero<b}{%
      \FormulaRefAuto{PositiveIntegersDef}[def]{2}}
    \proofstepwidestar[1,2,3]{b\cdot\IntZero=\IntZero\land\IntZero\cdot b=\IntZero}{%
      \FormulaRefAuto{IntegerMulZeroFromAxioms}[thm]{%
          5,%
          \rAEa{6}}}
    \proofstepwidestarbreakkeep[1,2,3]{a\cdot\IntOne
      =\IntZero\cdot b}{%
      \rIE{3,4,%
        7}}
    \proofstepwidestar[1,2,3]{\IntZero \in \mathbb Z}{%
      \FormulaRefAuto{IntegerZeroCarrierAxiom}[ax]}
    \proofstepwidestar[1,2,3]{\IntZero < \IntOne}{%
      \FormulaRefAuto{IntegerOnePositiveAxiom}[ax]}
    \proofstepwidestarbreakkeep[1,2,3]{\RatClass{a}{b}
      =\RatClass{\IntZero}{\IntOne}}{%
      \FormulaRefAuto{RationalClassEquality}[thm]{%
        1,9,2,%
        10,8}}
    \proofstepwidestar[1,2,3]{\RatZero = \RatClass \IntZero \IntOne \land \RatOne = \RatClass \IntOne \IntOne}{%
      \FormulaRefAuto{RationalZeroOneClasses}[thm]}
    \proofstepwidestarbreakkeep[1,2,3]{\RatClass{a}{b}=\RatZero}{%
      \rIE{11,\rAEa{12}}}
  \closeproofpartwideindR

  \proofpartwide{Schluss}
    \proofstepwidestarbreakkeep[1]{a\in\mathbb Z}{\rA}
    \proofstepwidestarbreakkeep[2]{b\in\PositiveIntegers}{\rA}
    \proofstepwidestarbreakkeep[1,2]{\RatClass{a}{b}=\RatZero
      \to a=\IntZero}{\FormulaRefAuto{RationalClassZeroNumeratorForward}{1,2}
}
    \proofstepwidestarbreakkeep[1,2]{a=\IntZero
      \to\RatClass{a}{b}=\RatZero}{\FormulaRefAuto{RationalClassZeroNumeratorBackward}{1,2}
}
    \proofstepwidestarbreakkeep[1,2]{\RatClass{a}{b}=\RatZero
      \leftrightarrow a=\IntZero}{\rLRI{3,4}}
  \closeproofpartwide
\end{tabproofsplitwide}

\FormulaDefDeltaK[Reziproker Repräsentantenwert]{%
  \begin{aligned}[t]
    &a\in\mathbb Z\dsep b\in\PositiveIntegers\dsep
      a\neq\IntZero\vdash{}\\
    &\RatRecipRep{a}{b}\coloneqq
      \begin{cases}
        \RatClass{b}{a},&\IntZero<a,\\
        \RatClass{-b}{-a},&a<\IntZero.
      \end{cases}
  \end{aligned}%
}{RationalReciprocalRepresentativeDef}{%
  \DeltaRow{Mengen}{a\dsep b}%
  \DeltaRow{Neue Symbole}{\RatRecipRep{a}{b}}%
}

Wegen der Totalität der Ganzzahlordnung gilt für
\(a\neq\IntZero\) genau einer der beiden Fälle. Im ersten Fall ist
\(a\) selbst ein positiver Nenner, im zweiten Fall ist \(-a\) positiv.

\FormulaThmDeltaK[Typisierung des reziproken Repräsentantenwertes]{%
  a\in\mathbb Z\dsep b\in\PositiveIntegers\dsep
  a\neq\IntZero
  \vdash\RatRecipRep{a}{b}\in\mathbb Q%
}{RationalReciprocalRepresentativeInQ}{%
  \DeltaRow{Mengen}{a\dsep b}%
}
\begin{tabproofwide}
  \proofstepwidestarbreakkeep[1]{a\in\mathbb Z}{\rA}
  \proofstepwidestarbreakkeep[2]{b\in\PositiveIntegers}{\rA}
  \proofstepwidestarbreakkeep[3]{a\neq\IntZero}{\rA}
  \proofstepwidestar[]{\IntZero\in\mathbb Z}{%
      \FormulaRefAuto{IntegerZeroCarrierAxiom}[ax]}
  \proofstepwidestarbreakkeep[1,3]{\IntZero<a\lor a<\IntZero}{%
    \FormulaRefAuto{IntegerOrderComparableAxiom}[ax]{%
      4,1},%
    \FormulaRefAuto{IntegerStrictOrderDef}[def]{1,3}}
  \proofstepwidestarbreakkeep[6]{\IntZero<a}{\rA}
  \proofstepwidestarbreakkeep[1,2,6]{\RatRecipRep{a}{b}=\RatClass{b}{a}}{%
    \FormulaRefAuto{RationalReciprocalRepresentativeDef}[def]{1,2,6}}
  \proofstepwidestar[2]{b\in\mathbb Z\land\IntZero<b}{%
      \FormulaRefAuto{PositiveIntegersDef}[def]{2}}
  \proofstepwidestar[1,6]{a\in\PositiveIntegers}{%
      \FormulaRefAuto{PositiveIntegersDef}[def]{1,6}}
  \proofstepwidestar[1,2,6]{\RatClass{b}{a}\in\mathbb Q}{%
      \FormulaRefAuto{RationalClassInQ}[thm]{%
      \rAEa{8},%
      9}}
  \proofstepwidestarbreakkeep[1,2,6]{\RatRecipRep{a}{b}\in\mathbb Q}{%
    \rIE{7,10}}
  \proofstepwidestarbreakkeep[12]{a<\IntZero}{\rA}
  \proofstepwidestarbreakkeep[1,12]{\IntZero<-a}{%
    \FormulaRefAuto{IntegerNegativeOrderReversal}[thm]{1,12}}
  \proofstepwidestarbreakkeep[1,2,12]{\RatRecipRep{a}{b}=\RatClass{-b}{-a}}{%
    \FormulaRefAuto{RationalReciprocalRepresentativeDef}[def]{1,2,12}}
  \proofstepwidestar[2]{b\in\mathbb Z\land\IntZero<b}{%
      \FormulaRefAuto{PositiveIntegersDef}[def]{2}}
  \proofstepwidestar[1,12]{-a\in\PositiveIntegers}{%
      \FormulaRefAuto{PositiveIntegersDef}[def]{1,13}}
  \proofstepwidestar[2]{-b\in\mathbb Z}{%
      \FormulaRefAuto{IntegerNegClosureAxiom}[ax]{%
        \rAEa{15}}}
  \proofstepwidestar[1,2,12]{\RatClass{-b}{-a}\in\mathbb Q}{%
      \FormulaRefAuto{RationalClassInQ}[thm]{%
      17,%
      16}}
  \proofstepwidestarbreakkeep[1,2,12]{\RatRecipRep{a}{b}\in\mathbb Q}{%
    \rIE{14,18}}
  \proofstepwidestarbreakkeep[1,2,3]{\RatRecipRep{a}{b}\in\mathbb Q}{%
    \rOE{5,6,11,12,19}}
\end{tabproofwide}

\FormulaThmDeltaK[Gleiche Kreuzprodukte haben gleiches Vorzeichen]{%
  \begin{aligned}[t]
    &a,c\in\mathbb Z\dsep b,d\in\PositiveIntegers\dsep
      a\neq\IntZero\dsep c\neq\IntZero\dsep a d=c b\vdash{}\\[-2pt]
    &\bigl((\IntZero<a\land\IntZero<c)
      \lor(a<\IntZero\land c<\IntZero)\bigr)
  \end{aligned}%
}{RationalCrossProductSameSign}{%
  \DeltaRow{Mengen}{a\dsep b\dsep c\dsep d}%
  \DeltaRow{Relationsoperatoren}{<}%
}
\begin{tabproofsplitwide}
  \proofpartwideindR[Positiver Zählerfall]{%
    \begin{aligned}[t]
      &a,c\in\mathbb Z\dsep b,d\in\PositiveIntegers\dsep
        a\neq\IntZero\dsep c\neq\IntZero\dsep{}\\[-2pt]
      &a d=c b\dsep\IntZero<a\vdash\IntZero<c
    \end{aligned}
  }{%
    a,c\in\mathbb Z\dsep b,d\in\PositiveIntegers\dsep
    a\neq\IntZero\dsep c\neq\IntZero\dsep
    a d=c b\dsep\IntZero<a\vdash\IntZero<c
  }[RationalCrossProductPositiveNumerators]
    \proofstepwidestarbreakkeep[1]{a,c\in\mathbb Z}{\rA}
    \proofstepwidestarbreakkeep[2]{b,d\in\PositiveIntegers}{\rA}
    \proofstepwidestarbreakkeep[3]{a\neq\IntZero}{\rA}
    \proofstepwidestarbreakkeep[4]{c\neq\IntZero}{\rA}
    \proofstepwidestarbreakkeep[5]{a d=c b}{\rA}
    \proofstepwidestarbreakkeep[6]{\IntZero<a}{\rA}
    \proofstepwidestarbreakkeep[2]{b\in\mathbb Z\land\IntZero<b}{%
      \FormulaRefAuto{PositiveIntegersDef}[def]{\rAEa{2}}}
    \proofstepwidestarbreakkeep[2]{d\in\mathbb Z\land\IntZero<d}{%
      \FormulaRefAuto{PositiveIntegersDef}[def]{\rAEb{2}}}
    \proofstepwidestar[1,2,6]{\IntZero \in \mathbb Z}{%
      \FormulaRefAuto{IntegerZeroCarrierAxiom}[ax]}
    \proofstepwidestarbreakkeep[1,2,6]{\IntZero\cdot d<a\cdot d}{%
      \FormulaRefAuto{IntegerPositiveFactorStrictOrderFromAxioms}[thm]{%
        9,\rAEa{1},8,6}}
    \proofstepwidestar[2]{d \cdot \IntZero = \IntZero \land \IntZero \cdot d = \IntZero}{%
      \FormulaRefAuto{IntegerMulZeroFromAxioms}[thm]{\rAEa{8}}}
    \proofstepwidestarbreakkeep[2]{\IntZero\cdot d=\IntZero}{%
      \rAEb{11}}
    \proofstepwidestarbreakkeep[1,2,6]{\IntZero<a d}{\rIE{10,12}}
    \proofstepwidestarbreakkeep[1,2,5,6]{\IntZero<c b}{\rIE{5,13}}
    \proofstepwidestar[2]{b \cdot \IntZero = \IntZero \land \IntZero \cdot b = \IntZero}{%
      \FormulaRefAuto{IntegerMulZeroFromAxioms}[thm]{\rAEa{7}}}
    \proofstepwidestarbreakkeep[2]{\IntZero\cdot b=\IntZero}{%
      \rAEb{15}}
    \proofstepwidestar[1,2,5,6]{\IntZero\leq c b\land\IntZero\neq c b}{%
      \FormulaRefAuto{IntegerStrictOrderDef}[def]{14}}
    \proofstepwidestarbreakkeep[1,2,5,6]{\IntZero\cdot b\leq c b}{%
      \rIE{16,\rAEa{17}}}
    \proofstepwidestar[1,2,5,6]{\IntZero \in \mathbb Z}{%
      \FormulaRefAuto{IntegerZeroCarrierAxiom}[ax]}
    \proofstepwidestar[1,2,5,6]{\IntZero\leq c\leftrightarrow\IntZero\cdot b\leq c\cdot b}{%
      \FormulaRefAuto{IntegerPositiveFactorOrderAxiom}[ax]{%
        19,\rAEb{1},7}}
    \proofstepwidestarbreakkeep[1,2,5,6]{\IntZero\leq c}{%
      \rLREb{20,18}}
    \proofstepwidestarbreakkeep[4]{\IntZero\neq c}{%
      \FormulaRefAuto{a\neq b\vdash b\neq a}[thm]{4}}
    \proofstepwidestarbreakkeep[1,2,4,5,6]{\IntZero<c}{%
      \FormulaRefAuto{IntegerStrictOrderDef}[def]{21,22}}
  \closeproofpartwideindR

  \proofpartwideindR[Negativer Zählerfall]{%
    \begin{aligned}[t]
      &a,c\in\mathbb Z\dsep b,d\in\PositiveIntegers\dsep
        a\neq\IntZero\dsep c\neq\IntZero\dsep{}\\[-2pt]
      &a d=c b\dsep a<\IntZero\vdash c<\IntZero
    \end{aligned}
  }{%
    a,c\in\mathbb Z\dsep b,d\in\PositiveIntegers\dsep
    a\neq\IntZero\dsep c\neq\IntZero\dsep
    a d=c b\dsep a<\IntZero\vdash c<\IntZero
  }[RationalCrossProductNegativeNumerators]
    \proofstepwidestarbreakkeep[1]{a,c\in\mathbb Z}{\rA}
    \proofstepwidestarbreakkeep[2]{b,d\in\PositiveIntegers}{\rA}
    \proofstepwidestarbreakkeep[3]{a\neq\IntZero}{\rA}
    \proofstepwidestarbreakkeep[4]{c\neq\IntZero}{\rA}
    \proofstepwidestarbreakkeep[5]{a d=c b}{\rA}
    \proofstepwidestarbreakkeep[6]{a<\IntZero}{\rA}
    \proofstepwidestarbreakkeep[2]{b\in\mathbb Z\land\IntZero<b}{%
      \FormulaRefAuto{PositiveIntegersDef}[def]{\rAEa{2}}}
    \proofstepwidestarbreakkeep[2]{d\in\mathbb Z\land\IntZero<d}{%
      \FormulaRefAuto{PositiveIntegersDef}[def]{\rAEb{2}}}
    \proofstepwidestar[1,2,6]{\IntZero \in \mathbb Z}{%
      \FormulaRefAuto{IntegerZeroCarrierAxiom}[ax]}
    \proofstepwidestarbreakkeep[1,2,6]{a\cdot d<\IntZero\cdot d}{%
      \FormulaRefAuto{IntegerPositiveFactorStrictOrderFromAxioms}[thm]{%
        \rAEa{1},9,8,6}}
    \proofstepwidestar[2]{d \cdot \IntZero = \IntZero \land \IntZero \cdot d = \IntZero}{%
      \FormulaRefAuto{IntegerMulZeroFromAxioms}[thm]{\rAEa{8}}}
    \proofstepwidestarbreakkeep[2]{\IntZero\cdot d=\IntZero}{%
      \rAEb{11}}
    \proofstepwidestarbreakkeep[1,2,6]{a d<\IntZero}{\rIE{10,12}}
    \proofstepwidestarbreakkeep[1,2,5,6]{c b<\IntZero}{\rIE{5,13}}
    \proofstepwidestar[2]{b \cdot \IntZero = \IntZero \land \IntZero \cdot b = \IntZero}{%
      \FormulaRefAuto{IntegerMulZeroFromAxioms}[thm]{\rAEa{7}}}
    \proofstepwidestarbreakkeep[2]{\IntZero\cdot b=\IntZero}{%
      \rAEb{15}}
    \proofstepwidestar[1,2,5,6]{c b\leq\IntZero\land c b\neq\IntZero}{%
      \FormulaRefAuto{IntegerStrictOrderDef}[def]{14}}
    \proofstepwidestarbreakkeep[1,2,5,6]{c b\leq\IntZero\cdot b}{%
      \rIE{16,\rAEa{17}}}
    \proofstepwidestar[1,2,5,6]{\IntZero \in \mathbb Z}{%
      \FormulaRefAuto{IntegerZeroCarrierAxiom}[ax]}
    \proofstepwidestar[1,2,5,6]{c\leq\IntZero\leftrightarrow c\cdot b\leq\IntZero\cdot b}{%
      \FormulaRefAuto{IntegerPositiveFactorOrderAxiom}[ax]{%
        \rAEb{1},19,7}}
    \proofstepwidestarbreakkeep[1,2,5,6]{c\leq\IntZero}{%
      \rLREb{20,18}}
    \proofstepwidestarbreakkeep[4]{c\neq\IntZero}{\rA}
    \proofstepwidestarbreakkeep[1,2,4,5,6]{c<\IntZero}{%
      \FormulaRefAuto{IntegerStrictOrderDef}[def]{21,22}}
  \closeproofpartwideindR

  \proofpartwide{Fallschluss}
    \proofstepwidestarbreakkeep[1]{a,c\in\mathbb Z}{\rA}
    \proofstepwidestarbreakkeep[2]{b,d\in\PositiveIntegers}{\rA}
    \proofstepwidestarbreakkeep[3]{a\neq\IntZero}{\rA}
    \proofstepwidestarbreakkeep[4]{c\neq\IntZero}{\rA}
    \proofstepwidestarbreakkeep[5]{a d=c b}{\rA}
    \proofstepwidestar[]{\IntZero\in\mathbb Z}{%
      \FormulaRefAuto{IntegerZeroCarrierAxiom}[ax]}
    \proofstepwidestarbreakkeep[1,3]{\IntZero<a\lor a<\IntZero}{%
      \FormulaRefAuto{IntegerOrderComparableAxiom}[ax]{%
        6,\rAEa{1}},%
      \FormulaRefAuto{IntegerStrictOrderDef}[def]{1,3}}
    \proofstepwidestarbreakkeep[8]{\IntZero<a}{\rA}
    \proofstepwidestarbreakkeep[1,2,3,4,5,8]{\IntZero<c}{%
      \FormulaRefAuto{RationalCrossProductPositiveNumerators}{1,2,3,4,5,8}}
    \proofstepwidestarbreakkeep[1,2,3,4,5,8]{%
      (\IntZero<a\land\IntZero<c)\lor(a<\IntZero\land c<\IntZero)}{%
      \rOIa{\rAI{8,9}}}
    \proofstepwidestarbreakkeep[11]{a<\IntZero}{\rA}
    \proofstepwidestarbreakkeep[1,2,3,4,5,11]{c<\IntZero}{%
      \FormulaRefAuto{RationalCrossProductNegativeNumerators}{1,2,3,4,5,11}}
    \proofstepwidestarbreakkeep[1,2,3,4,5,11]{%
      (\IntZero<a\land\IntZero<c)\lor(a<\IntZero\land c<\IntZero)}{%
      \rOIb{\rAI{11,12}}}
    \proofstepwidestarbreakkeep[1,2,3,4,5]{%
      (\IntZero<a\land\IntZero<c)\lor(a<\IntZero\land c<\IntZero)}{%
      \rOE{7,8,10,11,13}}
  \closeproofpartwide
\end{tabproofsplitwide}

\FormulaThmDeltaK[Repräsentantenunabhängigkeit des Kehrwertes]{%
  \begin{aligned}[t]
    &a,c\in\mathbb Z\dsep b,d\in\PositiveIntegers\dsep
      a\neq\IntZero\dsep c\neq\IntZero\dsep{}\\
    &\RatClass{a}{b}=\RatClass{c}{d}
      \vdash\RatRecipRep{a}{b}=\RatRecipRep{c}{d}
  \end{aligned}%
}{RationalReciprocalCompatible}{%
  \DeltaRow{Mengen}{a\dsep b\dsep c\dsep d}%
}
\begin{tabproofwide}
  \proofstepwidestarbreakkeep[1]{a,c\in\mathbb Z}{\rA}
  \proofstepwidestarbreakkeep[2]{b,d\in\PositiveIntegers}{\rA}
  \proofstepwidestarbreakkeep[3]{a\neq\IntZero}{\rA}
  \proofstepwidestarbreakkeep[4]{c\neq\IntZero}{\rA}
  \proofstepwidestarbreakkeep[5]{\RatClass{a}{b}=\RatClass{c}{d}}{\rA}
  \proofstepwidestarbreakkeep[1,2,5]{a\cdot d=c\cdot b}{%
    \FormulaRefAuto{RationalClassEquality}[thm]{1,2,5}}
  \proofstepwidestarbreakkeep[1,2,3,4,5]{%
    (\IntZero<a\land\IntZero<c)
    \lor(a<\IntZero\land c<\IntZero)}{%
    \FormulaRefAuto{RationalCrossProductSameSign}[thm]{1,2,3,4,6}}
  \proofstepwidestarbreakkeep[2]{b,d\in\mathbb Z}{%
    \rAIRepeat{2}{\rAEa{\FormulaRefAuto{PositiveIntegersDef}[def]}}{%
      \rAEa{2}\mid\rAEb{2}}}
  \proofstepwidestarbreakkeep[1,2,5]{b c=c b}{%
    \FormulaRefAuto{IntegerMulCommutativeAxiom}[ax]{1,8}}
  \proofstepwidestarbreakkeep[1,2,5]{c b=a d}{%
    \FormulaRefAuto{a=b\vdash b=a}[thm]{6}}
  \proofstepwidestarbreakkeep[1,2,5]{a d=d a}{%
    \FormulaRefAuto{IntegerMulCommutativeAxiom}[ax]{1,8}}
  \proofstepwidestarbreakkeep[1,2,5]{b c=d a}{\rChain{9,10,11}}
  \proofstepwidestarbreakkeep[13]{\IntZero<a\land\IntZero<c}{\rA}
  \proofstepwidestarbreakkeep[1,13]{a,c\in\PositiveIntegers}{%
    \rAIRepeat{2}{\FormulaRefAuto{PositiveIntegersDef}[def]}{%
      \rAEa{1},\rAEa{13}\mid\rAEb{1},\rAEb{13}}}
  \proofstepwidestarbreakkeep[1,2,13]{\RatRecipRep{a}{b}=\RatClass{b}{a}}{%
    \FormulaRefAuto{RationalReciprocalRepresentativeDef}[def]{1,2,\rAEa{13}}}
  \proofstepwidestarbreakkeep[1,2,13]{\RatRecipRep{c}{d}=\RatClass{d}{c}}{%
    \FormulaRefAuto{RationalReciprocalRepresentativeDef}[def]{1,2,\rAEb{13}}}
  \proofstepwidestarbreakkeep[1,2,5,13]{\RatClass{b}{a}=\RatClass{d}{c}}{%
    \FormulaRefAuto{RationalClassEquality}[thm]{8,14,12}}
  \proofstepwidestarbreakkeep[1,2,5,13]{\RatRecipRep{a}{b}
    =\RatRecipRep{c}{d}}{\rIE{15,16,17}}
  \proofstepwidestarbreakkeep[19]{a<\IntZero\land c<\IntZero}{\rA}
  \proofstepwidestarbreakkeep[1,2]{-b,-d\in\mathbb Z}{%
    \rAIRepeat{2}{\FormulaRefAuto{IntegerNegClosureAxiom}[ax]}{%
      \rAEa{8}\mid\rAEb{8}}}
  \proofstepwidestarbreakkeep[1,19]{\IntZero<-a}{%
    \FormulaRefAuto{IntegerNegativeOrderReversal}[thm]{%
      \rAEa{1},\rAEa{19}}}
  \proofstepwidestarbreakkeep[1,19]{\IntZero<-c}{%
    \FormulaRefAuto{IntegerNegativeOrderReversal}[thm]{%
      \rAEb{1},\rAEb{19}}}
  \proofstepwidestarbreakkeep[1,19]{-a,-c\in\PositiveIntegers}{%
    \rAIRepeat{2}{%
      \FormulaRefAuto{PositiveIntegersDef}[def]\circ
      \FormulaRefAuto{IntegerNegClosureAxiom}[ax]}{%
      \rAEa{1},21\mid\rAEb{1},22}}
  \proofstepwidestarbreakkeep[1,2,19]{\RatRecipRep{a}{b}=\RatClass{-b}{-a}}{%
    \FormulaRefAuto{RationalReciprocalRepresentativeDef}[def]{1,2,\rAEa{19}}}
  \proofstepwidestarbreakkeep[1,2,19]{\RatRecipRep{c}{d}=\RatClass{-d}{-c}}{%
    \FormulaRefAuto{RationalReciprocalRepresentativeDef}[def]{1,2,\rAEb{19}}}
  \proofstepwidestarbreakkeep[1,2,19]{(-b)(-c)=b c}{%
    \FormulaRefAuto{IntegerNegativeTimesNegative}[thm]{%
      \rAEa{8},\rAEb{1}}}
  \proofstepwidestarbreakkeep[1,2,19]{(-d)(-a)=d a}{%
    \FormulaRefAuto{IntegerNegativeTimesNegative}[thm]{%
      \rAEb{8},\rAEa{1}}}
  \proofstepwidestarbreakkeep[1,2,5,19]{(-b)(-c)=(-d)(-a)}{%
    \rIE{26,12,27}}
  \proofstepwidestarbreakkeep[1,2,5,19]{\RatClass{-b}{-a}
    =\RatClass{-d}{-c}}{%
    \FormulaRefAuto{RationalClassEquality}[thm]{20,23,28}}
  \proofstepwidestarbreakkeep[1,2,5,19]{\RatRecipRep{a}{b}
    =\RatRecipRep{c}{d}}{\rIE{24,25,29}}
  \proofstepwidestarbreakkeep[1,2,3,4,5]{\RatRecipRep{a}{b}
    =\RatRecipRep{c}{d}}{\rOE{7,13,18,19,30}}
\end{tabproofwide}

\Needspace{15\baselineskip}
\FormulaThmDeltaK[Quotientenabstieg des Kehrwertes]{%
  \begin{aligned}[t]
    &q\in\mathbb Q\dsep q\neq\RatZero\vdash{}\\
    &\exists!u\in\mathbb Q\,
      \exists a\in\mathbb Z\,\exists b\in\PositiveIntegers\,\Bigl({}\\[-2pt]
    &\qquad q=\RatClass{a}{b}\ \land\\[-2pt]
    &\qquad a\neq\IntZero\ \land\\[-2pt]
    &\qquad u=\RatRecipRep{a}{b}\Bigr)
  \end{aligned}%
}{RationalReciprocalQuotientDescent}{%
  \DeltaRow{Mengen}{q\dsep u\dsep v\dsep a\dsep a'\dsep b\dsep b'}%
  \DeltaRow{Weitere Mengen}{x\dsep t\dsep c\dsep d}%
}
\begin{tabproofwide}
  \proofstepwidestar[1]{q\in\mathbb Q}{\rA}
  \proofstepwidestar[2]{q\neq\RatZero}{\rA}
  \proofstepwidestar[1,2]{q\in\{t\in\mathbb Q\mid t\neq\RatZero\}}{%
    \FormulaRefAuto{x \in A\dsep P(x)\vdash x \in \{x \in A \mid P(x)\}}[thm]{1,2}}
  \proofstepwidestar[4]{a\in\mathbb Z\land b\in\PositiveIntegers\land a\neq\IntZero}{\rA}
  \proofstepwidestar[4]{\RatClass{a}{b}\in\mathbb Q}{%
    \FormulaRefAuto{RationalClassInQ}[thm]{\rAEa{4},\rAEa{\rAEb{4}}}}
  \proofstepwidestar[4]{\RatClass{a}{b}=\RatZero\leftrightarrow a=\IntZero}{%
    \FormulaRefAuto{RationalClassZeroIffNumeratorZero}[thm]{\rAEa{4},\rAEa{\rAEb{4}}}}
  \proofstepwidestar[4]{\RatClass{a}{b}\neq\RatZero}{%
    \FormulaRefAuto{P \leftrightarrow Q, \neg Q \vdash \neg P}[thm]{6,\rAEb{\rAEb{4}}}}
  \proofstepwidestar[4]{\RatClass{a}{b}\in\{t\in\mathbb Q\mid t\neq\RatZero\}}{%
    \FormulaRefAuto{x \in A\dsep P(x)\vdash x \in \{x \in A \mid P(x)\}}[thm]{5,7}}
  \proofstepwidestar[]{\begin{aligned}[t]
    &\forall a\forall b\,\Bigl((a\in\mathbb Z\land b\in\PositiveIntegers\land a\neq\IntZero)\\[-2pt]
    &\qquad\rightarrow\RatClass{a}{b}\in\{t\in\mathbb Q\mid t\neq\RatZero\}\Bigr)
  \end{aligned}}{\rUI{\rUI{\rRI{4,8}}}}
  \proofstepwidestar[10]{x\in\{t\in\mathbb Q\mid t\neq\RatZero\}}{\rA}
  \proofstepwidestar[10]{x\in\mathbb Q\land x\neq\RatZero}{%
    \FormulaRefAuto{x \in \{u \in A \mid P(u)\} \eqvdash x \in A \land P(x)}[thm]{10}}
  \proofstepwidestar[10]{\begin{aligned}[t]
    &\exists a\exists b\,\bigl((a\in\mathbb Z\land b\in\PositiveIntegers)\\[-2pt]
    &\qquad\land x=\RatClass{a}{b}\bigr)
  \end{aligned}}{%
    \rRE{\rUE{\FormulaRefAuto{RationalPairRepresentationSurjective}[thm]},\rAEa{11}}}
  \proofstepwidestar[13]{(a\in\mathbb Z\land b\in\PositiveIntegers)\land x=\RatClass{a}{b}}{\rA}
  \proofstepwidestar[13]{\RatClass{a}{b}=\RatZero\leftrightarrow a=\IntZero}{%
    \FormulaRefAuto{RationalClassZeroIffNumeratorZero}[thm]{\rAEa{\rAEa{13}},\rAEb{\rAEa{13}}}}
  \proofstepwidestar[10,13]{\RatClass{a}{b}\neq\RatZero}{\rIE{\rAEb{13},\rAEb{11}}}
  \proofstepwidestar[10,13]{a\neq\IntZero}{%
    \FormulaRefAuto{P \leftrightarrow Q, \neg P \vdash \neg Q}[thm]{14,15}}
  \proofstepwidestar[10,13]{\begin{aligned}[t]
    &\exists a\exists b\,\bigl((a\in\mathbb Z\land b\in\PositiveIntegers\land a\neq\IntZero)\\[-2pt]
    &\qquad\land x=\RatClass{a}{b}\bigr)
  \end{aligned}}{\rEIStar{\rAI{\rAIStar{\rAEa{\rAEa{13}},\rAEb{\rAEa{13}},16},\rAEb{13}}}}
  \proofstepwidestar[10]{\begin{aligned}[t]
    &\exists a\exists b\,\bigl((a\in\mathbb Z\land b\in\PositiveIntegers\land a\neq\IntZero)\\[-2pt]
    &\qquad\land x=\RatClass{a}{b}\bigr)
  \end{aligned}}{\rEEStar{12,13,17}}
  \proofstepwidestar[]{\begin{aligned}[t]
    &\forall x\in\{t\in\mathbb Q\mid t\neq\RatZero\}\,\exists a\exists b\,\\[-2pt]
    &\quad\bigl((a\in\mathbb Z\land b\in\PositiveIntegers\land a\neq\IntZero)\land x=\RatClass{a}{b}\bigr)
  \end{aligned}}{\rUI{\rRI{10,18}}}
  \proofstepwidestar[4]{\RatRecipRep{a}{b}\in\mathbb Q}{%
    \FormulaRefAuto{RationalReciprocalRepresentativeInQ}[thm]{\rAEa{4},\rAEa{\rAEb{4}},\rAEb{\rAEb{4}}}}
  \proofstepwidestar[]{\begin{aligned}[t]
    &\forall a\forall b\,\Bigl((a\in\mathbb Z\land b\in\PositiveIntegers\land a\neq\IntZero)\\[-2pt]
    &\qquad\rightarrow\RatRecipRep{a}{b}\in\mathbb Q\Bigr)
  \end{aligned}}{\rUI{\rUI{\rRI{4,20}}}}
  \proofstepwidestar[22]{\begin{aligned}[t]
    &(a\in\mathbb Z\land b\in\PositiveIntegers\land a\neq\IntZero)\land{}\\[-2pt]
    &(c\in\mathbb Z\land d\in\PositiveIntegers\land c\neq\IntZero)\land{}\\[-2pt]
    &\RatClass{a}{b}=\RatClass{c}{d}
  \end{aligned}}{\rA}
  \proofstepwidestar[22]{\RatRecipRep{a}{b}=\RatRecipRep{c}{d}}{%
    \FormulaRefAuto{RationalReciprocalCompatible}[thm]{\rAI{\rAEa{\rAEa{22}},\rAEa{\rAEa{\rAEb{22}}}},\rAI{\rAEa{\rAEb{\rAEa{22}}},\rAEa{\rAEb{\rAEa{\rAEb{22}}}}},\rAEb{\rAEb{\rAEa{22}}},\rAEb{\rAEb{\rAEa{\rAEb{22}}}},\rAEb{\rAEb{22}}}}
  \proofstepwidestar[]{\begin{aligned}[t]
    &\forall a\forall b\forall c\forall d\,\Bigl(\\[-2pt]
    &(a\in\mathbb Z\land b\in\PositiveIntegers\land a\neq\IntZero)\land{}\\[-2pt]
    &(c\in\mathbb Z\land d\in\PositiveIntegers\land c\neq\IntZero)\land{}\\[-2pt]
    &\RatClass{a}{b}=\RatClass{c}{d}\rightarrow\RatRecipRep{a}{b}=\RatRecipRep{c}{d}\Bigr)
  \end{aligned}}{\rUI{\rUI{\rUI{\rUI{\rRI{22,23}}}}}}
  \proofstepwidestar[1,2]{\begin{aligned}[t]
    &\exists!u\in\mathbb Q\,\exists a\exists b\,\\[-2pt]
    &\quad\bigl((a\in\mathbb Z\land b\in\PositiveIntegers\land a\neq\IntZero)\\[-2pt]
    &\qquad\land q=\RatClass{a}{b}\land u=\RatRecipRep{a}{b}\bigr)
  \end{aligned}}{\FormulaRefAuto{RepresentativePairValueUnique}[thm]{9,19,21,24,3}}
  \proofstepwidestar[]{\begin{aligned}[t]
&\forall u\,\Bigl(\exists a\exists b\,\\[-2pt]
    &\quad\bigl((a\in\mathbb Z\land b\in\PositiveIntegers\land a\neq\IntZero)\\[-2pt]
    &\qquad\land q=\RatClass{a}{b}\land u=\RatRecipRep{a}{b}\bigr)\\[-2pt]
&\quad\leftrightarrow{}\exists a\in\mathbb Z\,\exists b\in\PositiveIntegers\,\\[-2pt]
    &\quad\bigl(q=\RatClass{a}{b}\land a\neq\IntZero\land u=\RatRecipRep{a}{b}\bigr)\Bigr)
\end{aligned}}{%
    \rUI{\FormulaRefAuto{P \leftrightarrow Q \vdash Q \leftrightarrow P}[thm]{\FormulaRefAuto{BoundedPairExistsGuardWithCondition}[thm]}}}
  \proofstepwidestar[1,2]{\begin{aligned}[t]
    &\exists!u\in\mathbb Q\,\exists a\in\mathbb Z\,\exists b\in\PositiveIntegers\,\\[-2pt]
    &\quad\bigl(q=\RatClass{a}{b}\land a\neq\IntZero\land u=\RatRecipRep{a}{b}\bigr)
  \end{aligned}}{%
    \FormulaRefAuto{UniqueExistenceUnderEquivalentConjuncts}[thm]{26,25}}
\end{tabproofwide}

\Needspace{15\baselineskip}
\FormulaDefDeltaK[Kehrwert einer von null verschiedenen rationalen Zahl]{%
  \begin{aligned}[t]
    &q\in\mathbb Q\dsep q\neq\RatZero\vdash{}\\
    &q^{-1}\coloneqq\iota u\Bigl(
      u\in\mathbb Q\land
      \exists a\in\mathbb Z\,\exists b\in\PositiveIntegers\,\bigl({}\\
    &\qquad q=\RatClass{a}{b}\land a\neq\IntZero\land
      u=\RatRecipRep{a}{b}\bigr)
      \Bigr)
  \end{aligned}%
}{RationalReciprocalDef}{%
  \DeltaRow{Mengen}{q\dsep u\dsep a\dsep b}%
  \DeltaRow{Neue Symbole}{q^{-1}}%
}

\FormulaThmDeltaK[Kehrwert auf Bruchklassen]{%
  a\in\mathbb Z\dsep b\in\PositiveIntegers\dsep
  a\neq\IntZero
  \vdash
  \bigl(\RatClass{a}{b}\bigr)^{-1}=\RatRecipRep{a}{b}%
}{RationalReciprocalClassValue}{%
  \DeltaRow{Mengen}{a\dsep b}%
  \DeltaRow{Funktionsoperatoren}{(\,\cdot\,)^{-1}}%
}
\begin{tabproofwide}
  \proofstepwidestarbreakkeep[1]{a\in\mathbb Z}{\rA}
  \proofstepwidestarbreakkeep[2]{b\in\PositiveIntegers}{\rA}
  \proofstepwidestarbreakkeep[3]{a\neq\IntZero}{\rA}
  \proofstepwidestarbreakkeep[1,2]{\RatClass{a}{b}\in\mathbb Q}{%
    \FormulaRefAuto{RationalClassInQ}[thm]{1,2}}
  \proofstepwidestarbreakkeep[1,2,3]{\RatRecipRep{a}{b}\in\mathbb Q}{%
    \FormulaRefAuto{RationalReciprocalRepresentativeInQ}[thm]{1,2,3}}
  \proofstepwidestar[1,2]{\RatClass{a}{b}=\RatZero\leftrightarrow a=\IntZero}{%
      \FormulaRefAuto{RationalClassZeroIffNumeratorZero}[thm]{1,2}}
  \proofstepwidestar[1,2,3]{\RatClass{a}{b}\neq\RatZero}{%
      \FormulaRefAuto{P\leftrightarrow Q,\neg Q\vdash\neg P}{6,3}}
  \proofstepwidestarbreakkeep[1,2,3]{\exists!u\in\mathbb Q\,
    \exists c\in\mathbb Z\,\exists d\in\PositiveIntegers\,
    (\RatClass{a}{b}=\RatClass{c}{d}\land c\neq\IntZero
     \land u=\RatRecipRep{c}{d})}{\FormulaRefAuto{RationalReciprocalQuotientDescent}[thm]{4,7}
}
  \proofstepwidestarbreakkeep[1,2,3]{(\RatClass{a}{b})^{-1}
    =\RatRecipRep{a}{b}}{%
    \FormulaRefAuto{RationalReciprocalDef}[def]{1,2,3,4,5,8}}
\end{tabproofwide}

\FormulaThmDeltaK[Abgeschlossenheit des rationalen Kehrwertes]{%
  q\in\mathbb Q\dsep q\neq\RatZero
  \vdash q^{-1}\in\mathbb Q%
}{RationalReciprocalClosure}{%
  \DeltaRow{Mengen}{q}%
}
\begin{tabproofwide}
  \proofstepwidestarbreakkeep[1]{q\in\mathbb Q}{\rA}
  \proofstepwidestarbreakkeep[2]{q\neq\RatZero}{\rA}
  \proofstepwidestarbreakkeep[1,2]{\exists!u\in\mathbb Q\,
    \exists a\in\mathbb Z\,\exists b\in\PositiveIntegers\,
    (q=\RatClass{a}{b}\land a\neq\IntZero\land
     u=\RatRecipRep{a}{b})}{%
    \FormulaRefAuto{RationalReciprocalQuotientDescent}[thm]{1,2}}
  \proofstepwidestarbreakkeep[1,2]{q^{-1}\in\mathbb Q}{%
    \FormulaRefAuto{RationalReciprocalDef}[def]{1,2,3}}
\end{tabproofwide}

\FormulaThmDeltaK[Repräsentantenrechnung für das multiplikative Inverse]{%
  a\in\mathbb Z\dsep b\in\PositiveIntegers\dsep a\neq\IntZero
  \vdash\RatClass{a}{b}\cdot\RatRecipRep{a}{b}=\RatOne%
}{RationalReciprocalRepresentativeProduct}{%
  \DeltaRow{Mengen}{a\dsep b}%
}
\begin{tabproofsplitwide}
  \proofpartwideindR[Positiver Zähler]{%
    a\in\mathbb Z\dsep b\in\PositiveIntegers\dsep
    \IntZero<a\vdash\RatClass{a}{b}\cdot\RatRecipRep{a}{b}=\RatOne
  }{%
    a\in\mathbb Z\dsep b\in\PositiveIntegers\dsep
    \IntZero<a\vdash\RatClass{a}{b}\cdot\RatRecipRep{a}{b}=\RatOne
  }[RationalReciprocalRepresentativeProductPositive]
    \proofstepwidestarbreakkeep[1]{a\in\mathbb Z}{\rA}
    \proofstepwidestarbreakkeep[2]{b\in\PositiveIntegers}{\rA}
    \proofstepwidestarbreakkeep[3]{\IntZero<a}{\rA}
    \proofstepwidestarbreakkeep[1,3]{a\in\PositiveIntegers}{%
      \FormulaRefAuto{PositiveIntegersDef}[def]{1,3}}
    \proofstepwidestar[2]{b\in\mathbb Z\land\IntZero<b}{%
      \FormulaRefAuto{PositiveIntegersDef}[def]{2}}
    \proofstepwidestarbreakkeep[2]{b\in\mathbb Z}{%
      \rAEa{5}}
    \proofstepwidestarbreakkeep[1,2,3]{\RatRecipRep{a}{b}=\RatClass{b}{a}}{%
      \FormulaRefAuto{RationalReciprocalRepresentativeDef}[def]{1,2,3}}
    \proofstepwidestar[1,2,3]{\RatClass a b \cdot \RatClass b a = \RatClass { a \cdot b } { b \cdot a }}{%
      \FormulaRefAuto{RationalMultiplicationClassValue}[thm]{%
        1,6,2,4}}
    \proofstepwidestarbreakkeep[1,2,3]{%
      \RatClass{a}{b}\cdot\RatRecipRep{a}{b}=\RatClass{a b}{b a}}{%
      \rIE{7,8}}
    \proofstepwidestarbreakkeep[1,2]{a b\in\mathbb Z}{%
      \FormulaRefAuto{IntMulClosure}[thm]{1,6}}
    \proofstepwidestarbreakkeep[1,2,3]{b a\in\PositiveIntegers}{%
      \FormulaRefAuto{PositiveIntegerMulClosure}[thm]{2,4}}
    \proofstepwidestarbreakkeep[1,2]{a b=b a}{%
      \FormulaRefAuto{IntegerMulCommutativeAxiom}[ax]{1,6}}
    \proofstepwidestar[1,2,3]{\IntOne \in \PositiveIntegers}{%
      \FormulaRefAuto{IntegerOnePositive}[thm]}
    \proofstepwidestar[1,2,3]{\IntOne \in \mathbb Z}{%
      \FormulaRefAuto{IntegerOneCarrierAxiom}[ax]}
    \proofstepwidestar[1,2]{(a b)\cdot\IntOne=a b\land\IntOne\cdot(a b)=a b}{%
      \FormulaRefAuto{IntegerMulOneAxiom}[ax]{10,6}}
    \proofstepwidestarbreakkeep[1,2,3]{\RatClass{a b}{b a}
      =\RatClass{\IntOne}{\IntOne}}{%
      \FormulaRefAuto{RationalClassEquality}[thm]{%
        10,14,11,%
        13,%
        \rIE{12,15}}}
    \proofstepwidestar[1,2,3]{\RatZero = \RatClass \IntZero \IntOne \land \RatOne = \RatClass \IntOne \IntOne}{%
      \FormulaRefAuto{RationalZeroOneClasses}[thm]}
    \proofstepwidestarbreakkeep[1,2,3]{%
      \RatClass{a}{b}\cdot\RatRecipRep{a}{b}=\RatOne}{%
      \rIE{9,16,\rAEb{17}}}
  \closeproofpartwideindR

  \proofpartwideindR[Negativer Zähler]{%
    a\in\mathbb Z\dsep b\in\PositiveIntegers\dsep
    a<\IntZero\vdash\RatClass{a}{b}\cdot\RatRecipRep{a}{b}=\RatOne
  }{%
    a\in\mathbb Z\dsep b\in\PositiveIntegers\dsep
    a<\IntZero\vdash\RatClass{a}{b}\cdot\RatRecipRep{a}{b}=\RatOne
  }[RationalReciprocalRepresentativeProductNegative]
    \proofstepwidestarbreakkeep[1]{a\in\mathbb Z}{\rA}
    \proofstepwidestarbreakkeep[2]{b\in\PositiveIntegers}{\rA}
    \proofstepwidestarbreakkeep[3]{a<\IntZero}{\rA}
    \proofstepwidestar[2]{b\in\mathbb Z\land\IntZero<b}{%
      \FormulaRefAuto{PositiveIntegersDef}[def]{2}}
    \proofstepwidestarbreakkeep[2]{b\in\mathbb Z}{%
      \rAEa{4}}
    \proofstepwidestarbreakkeep[1]{-b\in\mathbb Z}{%
      \FormulaRefAuto{IntegerNegClosureAxiom}[ax]{5}}
    \proofstepwidestarbreakkeep[1,3]{\IntZero<-a}{%
      \FormulaRefAuto{IntegerNegativeOrderReversal}[thm]{1,3}}
    \proofstepwidestar[1,3]{- a \in \mathbb Z}{%
      \FormulaRefAuto{IntegerNegClosureAxiom}[ax]{1}}
    \proofstepwidestarbreakkeep[1,3]{-a\in\PositiveIntegers}{%
      \FormulaRefAuto{PositiveIntegersDef}[def]{%
        8,7}}
    \proofstepwidestarbreakkeep[1,2,3]{\RatRecipRep{a}{b}=\RatClass{-b}{-a}}{%
      \FormulaRefAuto{RationalReciprocalRepresentativeDef}[def]{1,2,3}}
    \proofstepwidestar[1,2,3]{\RatClass a b \cdot \RatClass { ( - b ) } { ( - a ) } = \RatClass { a \cdot { ( - b ) } } { b \cdot { ( - a ) } }}{%
      \FormulaRefAuto{RationalMultiplicationClassValue}[thm]{%
        1,6,2,9}}
    \proofstepwidestarbreakkeep[1,2,3]{%
      \RatClass{a}{b}\cdot\RatRecipRep{a}{b}
      =\RatClass{a(-b)}{b(-a)}}{%
      \rIE{10,11}}
    \proofstepwide[1,2]{a(-b)}{=}{-(a b)}{%
      \FormulaRefAuto{IntegerRightNegativeProduct}[thm]{1,5}}
    \proofstepwide[1,2]{b(-a)}{=}{-(b a)}{%
      \FormulaRefAuto{IntegerRightNegativeProduct}[thm]{5,1}}
    \proofstepwidestarbreakkeep[1,2]{a b=b a}{%
      \FormulaRefAuto{IntegerMulCommutativeAxiom}[ax]{1,5}}
    \proofstepwidestarbreakkeep[1,2]{a(-b)=b(-a)}{\rIE{13,14,15}}
    \proofstepwidestarbreakkeep[1,2]{a(-b)\in\mathbb Z}{%
      \FormulaRefAuto{IntMulClosure}[thm]{1,6}}
    \proofstepwidestarbreakkeep[1,2,3]{b(-a)\in\PositiveIntegers}{%
      \FormulaRefAuto{PositiveIntegerMulClosure}[thm]{2,9}}
    \proofstepwidestar[1,2,3]{\IntOne \in \PositiveIntegers}{%
      \FormulaRefAuto{IntegerOnePositive}[thm]}
    \proofstepwidestar[1,2,3]{\IntOne \in \mathbb Z}{%
      \FormulaRefAuto{IntegerOneCarrierAxiom}[ax]}
    \proofstepwidestar[1,2]{a(-b)\cdot\IntOne=a(-b)\land\IntOne\cdot a(-b)=a(-b)}{%
      \FormulaRefAuto{IntegerMulOneAxiom}[ax]{17,5}}
    \proofstepwidestarbreakkeep[1,2,3]{\RatClass{a(-b)}{b(-a)}
      =\RatClass{\IntOne}{\IntOne}}{%
      \FormulaRefAuto{RationalClassEquality}[thm]{%
        17,20,18,%
        19,%
        \rIE{16,21}}}
    \proofstepwidestar[1,2,3]{\RatZero = \RatClass \IntZero \IntOne \land \RatOne = \RatClass \IntOne \IntOne}{%
      \FormulaRefAuto{RationalZeroOneClasses}[thm]}
    \proofstepwidestarbreakkeep[1,2,3]{%
      \RatClass{a}{b}\cdot\RatRecipRep{a}{b}=\RatOne}{%
      \rIE{12,22,\rAEb{23}}}
  \closeproofpartwideindR

  \proofpartwide{Fallschluss}
    \proofstepwidestarbreakkeep[1]{a\in\mathbb Z}{\rA}
    \proofstepwidestarbreakkeep[2]{b\in\PositiveIntegers}{\rA}
    \proofstepwidestarbreakkeep[3]{a\neq\IntZero}{\rA}
    \proofstepwidestar[]{\IntZero\in\mathbb Z}{%
      \FormulaRefAuto{IntegerZeroCarrierAxiom}[ax]}
    \proofstepwidestarbreakkeep[1,3]{\IntZero<a\lor a<\IntZero}{%
      \FormulaRefAuto{IntegerOrderComparableAxiom}[ax]{%
        4,1},%
      \FormulaRefAuto{IntegerStrictOrderDef}[def]{1,3}}
    \proofstepwidestarbreakkeep[6]{\IntZero<a}{\rA}
    \proofstepwidestarbreakkeep[1,2,6]{%
      \RatClass{a}{b}\cdot\RatRecipRep{a}{b}=\RatOne}{%
      \FormulaRefAuto{RationalReciprocalRepresentativeProductPositive}{1,2,6}}
    \proofstepwidestarbreakkeep[8]{a<\IntZero}{\rA}
    \proofstepwidestarbreakkeep[1,2,8]{%
      \RatClass{a}{b}\cdot\RatRecipRep{a}{b}=\RatOne}{%
      \FormulaRefAuto{RationalReciprocalRepresentativeProductNegative}{1,2,8}}
    \proofstepwidestarbreakkeep[1,2,3]{%
      \RatClass{a}{b}\cdot\RatRecipRep{a}{b}=\RatOne}{%
      \rOE{5,6,7,8,9}}
  \closeproofpartwide
\end{tabproofsplitwide}

\FormulaThmDeltaK[Multiplikatives Inverses rationaler Zahlen]{%
  q\in\mathbb Q\dsep q\neq\RatZero
  \vdash q\cdot q^{-1}=\RatOne%
}{RationalMulInverse}{%
  \DeltaRow{Mengen}{q}%
}
\begin{tabproofwide}
  \proofstepwidestarbreakkeep[1]{q\in\mathbb Q}{\rA}
  \proofstepwidestarbreakkeep[2]{q\neq\RatZero}{\rA}
  \proofstepwidestarbreakkeep[1]{\exists a\in\mathbb Z\,
    \exists b\in\PositiveIntegers\,q=\RatClass{a}{b}}{%
    \FormulaRefAuto{RationalPairRepresentative}[thm]{1}}
  \proofstepwidestarbreakkeep[4]{a\in\mathbb Z\land b\in\PositiveIntegers
    \land q=\RatClass{a}{b}}{\rA}
  \proofstepwidestar[4]{\RatClass{a}{b}=\RatZero\leftrightarrow a=\IntZero}{%
      \FormulaRefAuto{RationalClassZeroIffNumeratorZero}[thm]{\rAEa{4},\rAEa{\rAEb{4}}}}
  \proofstepwidestarbreakkeep[2,4]{a\neq\IntZero}{\FormulaRefAuto{P\leftrightarrow Q,\neg P\vdash\neg Q}{5,\rIE{\rAEb{\rAEb{4}},2}}
}
  \proofstepwidestarbreakkeep[4,6]{(\RatClass{a}{b})^{-1}
    =\RatRecipRep{a}{b}}{%
    \FormulaRefAuto{RationalReciprocalClassValue}[thm]{4,6}}
  \proofstepwidestarbreakkeep[4,6]{\RatClass{a}{b}\cdot\RatRecipRep{a}{b}
    =\RatOne}{%
    \FormulaRefAuto{RationalReciprocalRepresentativeProduct}[thm]{4,6}}
  \proofstepwidestarbreakkeep[4,6]{q\cdot q^{-1}=\RatOne}{\rIE{4,7,8}}
  \proofstepwidestarbreakkeep[1,2]{q\cdot q^{-1}=\RatOne}{\rEE{3,4,9}}
\end{tabproofwide}

\FormulaDefDeltaK[Division rationaler Zahlen]{%
  q,r\in\mathbb Q\dsep r\neq\RatZero\vdash
  \frac qr\coloneqq q\cdot r^{-1}%
}{RationalDivisionDef}{%
  \DeltaRow{Mengen}{q\dsep r}%
  \DeltaRow{Neue Symbole}{\frac qr}%
}

\FormulaThmDeltaK[Abgeschlossenheit der rationalen Division]{%
  q,r\in\mathbb Q\dsep r\neq\RatZero
  \vdash\frac qr\in\mathbb Q%
}{RationalDivisionClosure}{%
  \DeltaRow{Mengen}{q\dsep r}%
}
\begin{tabproof}
  \proofstep{1}{q,r\in\mathbb Q}{\rA}
  \proofstep{2}{r\neq\RatZero}{\rA}
  \proofstep{1,2}{r^{-1}\in\mathbb Q}{%
    \FormulaRefAuto{RationalReciprocalClosure}[thm]{\rAEb{1},2}}
  \proofstep{1,2}{q\cdot r^{-1}\in\mathbb Q}{%
    \FormulaRefAuto{RationalMulClosure}[thm]{\rAEa{1},3}}
  \proofstep{1,2}{\frac qr=q\cdot r^{-1}}{%
      \FormulaRefAuto{RationalDivisionDef}[def]{1,2}}
  \proofstep{1,2}{\frac qr\in\mathbb Q}{%
    \rIE{4,5}}
\end{tabproof}

\chapter{Ordnung der rationalen Zahlen}

\section{Quotientenordnung}

\FormulaThmDeltaK[Repräsentantenunabhängigkeit der rationalen Ordnung]{%
  \begin{aligned}[t]
    &a,a',c,c'\in\mathbb Z\dsep
      b,b',d,d'\in\PositiveIntegers\dsep
      \RatClass{a}{b}=\RatClass{a'}{b'}\dsep{}\\
    &\RatClass{c}{d}=\RatClass{c'}{d'}\vdash
      \bigl(a\cdot d\leq c\cdot b
      \leftrightarrow a'\cdot d'\leq c'\cdot b'\bigr)
  \end{aligned}%
}{RationalOrderCompatible}{%
  \DeltaRow{Mengen}{a\dsep a'\dsep b\dsep b'\dsep
    c\dsep c'\dsep d\dsep d'}%
}

\begin{tabproofwide}
  \proofstepwidestarbreakkeep[1]{a,a',c,c'\in\mathbb Z}{\rA}
  \proofstepwidestarbreakkeep[2]{b,b',d,d'\in\PositiveIntegers}{\rA}
  \proofstepwidestarbreakkeep[3]{\RatClass{a}{b}=\RatClass{a'}{b'}}{\rA}
  \proofstepwidestarbreakkeep[4]{\RatClass{c}{d}=\RatClass{c'}{d'}}{\rA}
  \proofstepwidestarbreakkeep[1,2,3]{a\cdot b'=a'\cdot b}{%
    \FormulaRefAuto{RationalClassEquality}[thm]{1,2,3}}
  \proofstepwidestarbreakkeep[1,2,4]{c\cdot d'=c'\cdot d}{%
    \FormulaRefAuto{RationalClassEquality}[thm]{1,2,4}}
  \proofstepwidestarbreakkeep[1,2,3,4]{%
    a\cdot d\leq c\cdot b
    \leftrightarrow a'\cdot d'\leq c'\cdot b'}{%
    \FormulaRefAuto{IntegerPositiveFactorOrderAxiom}[ax]{1,2},%
    \FormulaRefAuto{IntegerMulAssociativeAxiom}[ax]{1},%
    \FormulaRefAuto{IntegerMulCommutativeAxiom}[ax]{1},\rIE{5,6}}
\end{tabproofwide}

\FormulaDefDeltaK[Rationale Ordnung als Quotientenrelation]{%
  \begin{aligned}[t]
  \RatOrderGraph\coloneqq
  \bigl\{&(\RatClass{a}{b},\RatClass{c}{d})
    \in\mathbb Q\times\mathbb Q\bigm|{}\\[-2pt]
    &a,c\in\mathbb Z\land b,d\in\PositiveIntegers
      \land a\cdot d\leq c\cdot b\bigr\}.
  \end{aligned}%
}{RationalOrderRelationDef}{%
  \DeltaRow{Mengen}{a\dsep b\dsep c\dsep d\dsep\RatOrderGraph}%
  \DeltaRow{Neue Symbole}{\RatOrderGraph}%
}

\FormulaThmDeltaK[Quotientenabstieg der rationalen Ordnung]{%
  \begin{aligned}[t]
  &\RatOrderGraph\subseteq\mathbb Q\times\mathbb Q\land{}\\[-2pt]
  &\forall a,c\in\mathbb Z\,\forall b,d\in\PositiveIntegers\,
    \bigl((\RatClass{a}{b},\RatClass{c}{d})\in\RatOrderGraph
      \leftrightarrow a\cdot d\leq c\cdot b\bigr).
  \end{aligned}%
}{RationalOrderQuotientDescent}{%
  \DeltaRow{Mengen}{a\dsep b\dsep c\dsep d\dsep\RatOrderGraph}%
}
\begin{tabproofsplitwide}
\proofpartwide{Trägerrelation}
\proofstepwidestarbreakkeep[]{\RatOrderGraph\subseteq\mathbb Q\times\mathbb Q}{%
    \FormulaRefAuto{RationalOrderRelationDef}[def]}

\closeproofpartwide

\proofpartwide{Auswertung auf Repräsentanten}
  \setcounter{proofstepnr}{1}% Fortlaufende Schritte dieses Beweises.
\proofstepwidestarbreakkeep[2]{a,c\in\mathbb Z\land
    b,d\in\PositiveIntegers}{\rA}
  \proofstepwidestarbreakkeep[2]{(\RatClass{a}{b},\RatClass{c}{d})
    \in\RatOrderGraph\leftrightarrow a\cdot d\leq c\cdot b}{%
    \FormulaRefAuto{RationalOrderRelationDef}[def]{2},%
    \FormulaRefAuto{RationalOrderCompatible}[thm]}
  \proofstepwidestarbreakkeep[]{\forall a,c\in\mathbb Z\,
    \forall b,d\in\PositiveIntegers\,
    \bigl((\RatClass{a}{b},\RatClass{c}{d})\in\RatOrderGraph
      \leftrightarrow a\cdot d\leq c\cdot b\bigr)}{%
    \rUI{\rRI{2,3}}}

\closeproofpartwide

\proofpartwide{Zusammenführung}
  \setcounter{proofstepnr}{4}% Fortlaufende Schritte dieses Beweises.
\proofstepwidestarbreakkeep[]{\RatOrderGraph\subseteq\mathbb Q\times\mathbb Q
    \land\forall a,c\in\mathbb Z\,\forall b,d\in\PositiveIntegers\,
    \bigl((\RatClass{a}{b},\RatClass{c}{d})\in\RatOrderGraph
      \leftrightarrow a\cdot d\leq c\cdot b\bigr)}{\rAI{1,4}}

\closeproofpartwide
\end{tabproofsplitwide}

\FormulaDefDeltaK[Ordnungsoperator der rationalen Zahlen]{%
  q,r\in\mathbb Q\vdash
  q\leq_{\mathbb Q}r\coloneqq(q,r)\in\RatOrderGraph%
}{RationalOrderDef}{%
  \DeltaRow{Mengen}{q\dsep r\dsep\RatOrderGraph}%
  \DeltaRow{Relationsoperatoren}{\leq_{\mathbb Q}}%
  \DeltaRow{Neue Symbole}{\leq_{\mathbb Q}}%
}

\FormulaDefDeltaK[Notation der rationalen Ordnung]{%
  q,r\in\mathbb Q\vdash
  q\leq r\coloneqq q\leq_{\mathbb Q}r%
}{RationalOrderNotation}{%
  \DeltaRow{Mengen}{q\dsep r}%
  \DeltaRow{Relationsoperatoren}{\leq\dsep\leq_{\mathbb Q}}%
}

\FormulaThmDeltaK[Klassencharakterisierung der rationalen Ordnung]{%
  a,c\in\mathbb Z\dsep b,d\in\PositiveIntegers\vdash
  \bigl(\RatClass{a}{b}\leq\RatClass{c}{d}
  \leftrightarrow a\cdot d\leq c\cdot b\bigr)%
}{RationalOrderClassChar}{%
  \DeltaRow{Mengen}{a\dsep b\dsep c\dsep d}%
  \DeltaRow{Relationsoperatoren}{\leq\dsep\leq_{\mathbb Q}}%
}
\begin{tabproofwide}
  \proofstepwidestarbreakkeep[1]{a,c\in\mathbb Z}{\rA}
  \proofstepwidestarbreakkeep[2]{b,d\in\PositiveIntegers}{\rA}
  \proofstepwidestar[]{\RatOrderGraph\subseteq\mathbb Q\times\mathbb Q\land\forall a,c\in\mathbb Z\,\forall b,d\in\PositiveIntegers\,((\RatClass{a}{b},\RatClass{c}{d})\in\RatOrderGraph\leftrightarrow a\cdot d\leq c\cdot b)}{%
      \FormulaRefAuto{RationalOrderQuotientDescent}[thm]}
  \proofstepwidestarbreakkeep[1,2]{(\RatClass{a}{b},\RatClass{c}{d})
    \in\RatOrderGraph\leftrightarrow a\cdot d\leq c\cdot b}{\rRE{\rUE{\rAEb{3}},\rAI{1,2}}
}
  \proofstepwidestar[1,2]{\RatClass{a}{b}\leq_{\mathbb Q}\RatClass{c}{d}\leftrightarrow(\RatClass{a}{b},\RatClass{c}{d})\in\RatOrderGraph}{%
      \FormulaRefAuto{RationalOrderDef}[def]{1,2}}
  \proofstepwidestar[1,2]{\RatClass{a}{b}\leq\RatClass{c}{d}\leftrightarrow\RatClass{a}{b}\leq_{\mathbb Q}\RatClass{c}{d}}{%
      \FormulaRefAuto{RationalOrderNotation}[def]{1,2}}
  \proofstepwidestarbreakkeep[1,2]{\RatClass{a}{b}\leq\RatClass{c}{d}
    \leftrightarrow a\cdot d\leq c\cdot b}{\rChain{6,5,4}
}
\end{tabproofwide}

\FormulaThmDeltaK[Totale Ordnung der rationalen Zahlen]{%
  \TotOrd{\mathbb Q,\leq_{\mathbb Q}}%
}{RationalOrderTotal}{%
  \DeltaRow{Mengen}{q\dsep r\dsep s}%
  \DeltaRow{Relationsoperatoren}{\leq_{\mathbb Q}}%
}
\begin{tabproofsplitwide}
\proofpartwide{Reflexivität}
\proofstepwidestarbreakkeep[1]{q,r,s\in\mathbb Q}{\rA}
    \proofstepwidestarbreakkeep[1]{q\leq q}{%
      \FormulaRefAuto{RationalPairRepresentative}[thm]{\rAEa{1}},%
      \FormulaRefAuto{RationalOrderClassChar}[thm],%
      \FormulaRefAuto{IntegerOrderReflexiveAxiom}[ax]}

\closeproofpartwide

\proofpartwide{Antisymmetrie}
  \setcounter{proofstepnr}{2}% Fortlaufende Schritte dieses Beweises.
\proofstepwidestarbreakkeep[1]{(q\leq r\land r\leq q)\to q=r}{%
      \FormulaRefAuto{RationalPairRepresentative}[thm]{1},%
      \FormulaRefAuto{RationalOrderClassChar}[thm],%
      \FormulaRefAuto{RationalClassEquality}[thm],%
      \FormulaRefAuto{IntegerOrderAntisymmetricAxiom}[ax]}

\closeproofpartwide

\proofpartwide{Transitivität}
  \setcounter{proofstepnr}{3}% Fortlaufende Schritte dieses Beweises.
\proofstepwidestarbreakkeep[1]{(q\leq r\land r\leq s)\to q\leq s}{%
      \FormulaRefAuto{RationalPairRepresentative}[thm]{1},%
      \FormulaRefAuto{RationalOrderClassChar}[thm],%
      \FormulaRefAuto{IntegerOrderTransitiveAxiom}[ax],%
      \FormulaRefAuto{IntegerPositiveFactorOrderAxiom}[ax],%
      \FormulaRefAuto{IntegerPositiveFactorCancellationFromAxioms}[thm]}

\closeproofpartwide

\proofpartwide{Totalität}
  \setcounter{proofstepnr}{4}% Fortlaufende Schritte dieses Beweises.
\proofstepwidestarbreakkeep[1]{q\leq r\lor r\leq q}{%
      \FormulaRefAuto{RationalPairRepresentative}[thm]{1},%
      \FormulaRefAuto{RationalOrderClassChar}[thm],%
      \FormulaRefAuto{IntegerOrderComparableAxiom}[ax]}

\closeproofpartwide

\proofpartwideindR[Ordnungsregeln aus Kreuzprodukten]{%
    \begin{aligned}[t]
      &q,r,s\in\mathbb Q\vdash
      q\leq q\ \land\
      ((q\leq r\land r\leq q)\to q=r)\ \land{}\\
      &((q\leq r\land r\leq s)\to q\leq s)\ \land\
      (q\leq r\lor r\leq q)
    \end{aligned}
  }{%
    \begin{aligned}[t]
      &q,r,s\in\mathbb Q\vdash
      q\leq q\ \land\
      ((q\leq r\land r\leq q)\to q=r)\ \land{}\\
      &((q\leq r\land r\leq s)\to q\leq s)\ \land\
      (q\leq r\lor r\leq q)
    \end{aligned}
  }[RationalOrderLawsOnRepresentatives]
  \setcounter{proofstepnr}{5}% Fortlaufende Schritte dieses Beweises.
\proofstepwidestarbreakkeep[1]{q\leq q\land
      ((q\leq r\land r\leq q)\to q=r)\land
      ((q\leq r\land r\leq s)\to q\leq s)\land
      (q\leq r\lor r\leq q)}{\rAI{2,3,4,5}}

\closeproofpartwideindR


  \proofpartwide{Schluss}
    \proofstepwidestarbreakkeep[]{\forall q\in\mathbb Q\,(q\leq q)}{%
      \FormulaRefAuto{RationalOrderLawsOnRepresentatives}}
    \proofstepwidestarbreakkeep[]{\forall q,r,s\in\mathbb Q\,
      ((q\leq r\land r\leq s)\to q\leq s)}{%
      \FormulaRefAuto{RationalOrderLawsOnRepresentatives}}
    \proofstepwidestarbreakkeep[]{\forall q,r\in\mathbb Q\,
      ((q\leq r\land r\leq q)\to q=r)}{%
      \FormulaRefAuto{RationalOrderLawsOnRepresentatives}}
    \proofstepwidestarbreakkeep[]{\PartOrd{\mathbb Q,\leq_{\mathbb Q}}}{%
      \FormulaRefAuto{Partielle Ordnung}[def]{1,2,3}}
    \proofstepwidestarbreakkeep[]{\forall q,r\in\mathbb Q\,
      (q\leq r\lor r\leq q)}{%
      \FormulaRefAuto{RationalOrderLawsOnRepresentatives}}
    \proofstepwidestarbreakkeep[]{\TotOrd{\mathbb Q,\leq_{\mathbb Q}}}{%
      \FormulaRefAuto{Totale Ordnung}[def]{4,5}}
  \closeproofpartwide
\end{tabproofsplitwide}

\FormulaDefDeltaK[Strikte Ordnung der rationalen Zahlen]{%
  q,r\in\mathbb Q\vdash
  q<r\coloneqq q\leq r\land q\neq r%
}{RationalStrictOrderDef}{%
  \DeltaRow{Mengen}{q\dsep r}%
  \DeltaRow{Relationsoperatoren}{<}%
}

\FormulaThmDeltaK[Strikte Klassencharakterisierung der rationalen Ordnung]{%
  a,c\in\mathbb Z\dsep b,d\in\PositiveIntegers\vdash
  \bigl(\RatClass{a}{b}<\RatClass{c}{d}
  \leftrightarrow a\cdot d<c\cdot b\bigr)%
}{RationalStrictOrderClassChar}{%
  \DeltaRow{Mengen}{a\dsep b\dsep c\dsep d}%
}
\begin{tabproofwide}
  \proofstepwidestarbreakkeep[1]{a,c\in\mathbb Z}{\rA}
  \proofstepwidestarbreakkeep[2]{b,d\in\PositiveIntegers}{\rA}
  \proofstepwidestarbreakkeep[1,2]{\RatClass{a}{b}\leq\RatClass{c}{d}
    \leftrightarrow a\cdot d\leq c\cdot b}{%
    \FormulaRefAuto{RationalOrderClassChar}[thm]{1,2}}
  \proofstepwidestarbreakkeep[1,2]{\RatClass{a}{b}=\RatClass{c}{d}
    \leftrightarrow a\cdot d=c\cdot b}{%
    \FormulaRefAuto{RationalClassEquality}[thm]{1,2}}
  \proofstepwidestar[1,2]{a\cdot d,c\cdot b\in\mathbb Z}{%
      \FormulaRefAuto{IntMulClosure}[thm]{1,2}}
  \proofstepwidestarbreakkeep[1,2]{\RatClass{a}{b}<\RatClass{c}{d}
    \leftrightarrow a\cdot d<c\cdot b}{%
    \FormulaRefAuto{RationalStrictOrderDef}[def]{1,2,3,4},%
    \FormulaRefAuto{IntegerStrictOrderDef}[def]{%
      5}}
\end{tabproofwide}

\section{Verträglichkeit von Ordnung und Arithmetik}

\FormulaThmDeltaK[Translationsinvarianz der rationalen Ordnung]{%
  p,q,r\in\mathbb Q\vdash
  \bigl(q\leq r\leftrightarrow p+q\leq p+r\bigr)%
}{RationalOrderTranslation}{%
  \DeltaRow{Mengen}{p\dsep q\dsep r}%
}
\begin{tabproofwide}
  \proofstepwidestarbreakkeep[1]{p,q,r\in\mathbb Q}{\rA}
  \proofstepwidestarbreakkeep[1]{\exists a\in\mathbb Z\,
    \exists b\in\PositiveIntegers\,p=\RatClass{a}{b}}{%
    \FormulaRefAuto{RationalPairRepresentative}[thm]{\rAEa{1}}}
  \proofstepwidestarbreakkeep[1]{\exists c\in\mathbb Z\,
    \exists d\in\PositiveIntegers\,q=\RatClass{c}{d}}{%
    \FormulaRefAuto{RationalPairRepresentative}[thm]{\rAEa{\rAEb{1}}}}
  \proofstepwidestarbreakkeep[1]{\exists e\in\mathbb Z\,
    \exists f\in\PositiveIntegers\,r=\RatClass{e}{f}}{%
    \FormulaRefAuto{RationalPairRepresentative}[thm]{\rAEb{\rAEb{1}}}}
  \proofstepwidestarbreakkeep[5]{a,c,e\in\mathbb Z\land
    b,d,f\in\PositiveIntegers\land p=\RatClass{a}{b}\land
    q=\RatClass{c}{d}\land r=\RatClass{e}{f}}{\rA}
  \proofstepwidestarbreakkeep[5]{\RatClass{c}{d}\leq\RatClass{e}{f}
    \leftrightarrow c f\leq e d}{%
    \FormulaRefAuto{RationalOrderClassChar}[thm]{5}}
  \proofstepwidestarbreakkeep[5]{\RatClass{a}{b}+\RatClass{c}{d}
    =\RatClass{a d+c b}{b d}}{%
    \FormulaRefAuto{RationalAdditionClassValue}[thm]{5}}
  \proofstepwidestarbreakkeep[5]{\RatClass{a}{b}+\RatClass{e}{f}
    =\RatClass{a f+e b}{b f}}{%
    \FormulaRefAuto{RationalAdditionClassValue}[thm]{5}}
  \proofstepwidestar[5]{\RatClass{a d+c b}{b d}\leq\RatClass{a f+e b}{b f}\leftrightarrow(a d+c b)(b f)\leq(a f+e b)(b d)}{%
      \FormulaRefAuto{RationalOrderClassChar}[thm]{5,7,8}}
  \proofstepwidestarbreakkeep[5]{%
    \RatClass{a}{b}+\RatClass{c}{d}
      \leq\RatClass{a}{b}+\RatClass{e}{f}
    \leftrightarrow
    (a d+c b)(b f)\leq(a f+e b)(b d)}{%
    \rIE{7,8,9}}
  \proofstepwidestarbreakkeep[5]{%
    (a d+c b)(b f)=(a d)(b f)+(c f)(b b)}{%
    \FormulaRefAuto{IntMulLeftDistributive}[thm]{5},%
    \FormulaRefAuto{IntMulAssociative}[thm]{5},%
    \FormulaRefAuto{IntMulCommutative}[thm]{5}}
  \proofstepwidestarbreakkeep[5]{%
    (a f+e b)(b d)=(a d)(b f)+(e d)(b b)}{%
    \FormulaRefAuto{IntMulLeftDistributive}[thm]{5},%
    \FormulaRefAuto{IntMulAssociative}[thm]{5},%
    \FormulaRefAuto{IntMulCommutative}[thm]{5}}
  \proofstepwidestar[5]{(c f)(b b)+(a d)(b f)=(a d)(b f)+(c f)(b b)}{%
      \FormulaRefAuto{IntAddCommutative}[thm]{5}}
  \proofstepwidestar[5]{(e d)(b b)+(a d)(b f)=(a d)(b f)+(e d)(b b)}{%
      \FormulaRefAuto{IntAddCommutative}[thm]{5}}
  \proofstepwidestar[5]{(c f)(b b)\leq(e d)(b b)\leftrightarrow(c f)(b b)+(a d)(b f)\leq(e d)(b b)+(a d)(b f)}{%
      \FormulaRefAuto{IntegerOrderTranslationEquivalenceAxiom}[ax]{5}}
  \proofstepwidestar[5]{(c f)(b b)\leq(e d)(b b)\leftrightarrow(a d)(b f)+(c f)(b b)\leq(a d)(b f)+(e d)(b b)}{\rIE{15,13,14}}
  \proofstepwidestarbreakkeep[5]{%
    (a d+c b)(b f)\leq(a f+e b)(b d)
    \leftrightarrow(c f)(b b)\leq(e d)(b b)}{%
    \rIE{11,12,16}}
  \proofstepwidestarbreakkeep[5]{b b\in\PositiveIntegers}{%
    \FormulaRefAuto{PositiveIntegerMulClosure}[thm]{5}}
  \proofstepwidestarbreakkeep[5]{c f\leq e d
    \leftrightarrow(c f)(b b)\leq(e d)(b b)}{%
    \FormulaRefAuto{IntegerPositiveFactorOrderAxiom}[ax]{5,18}}
  \proofstepwidestarbreakkeep[5]{%
    \RatClass{c}{d}\leq\RatClass{e}{f}
    \leftrightarrow
    \RatClass{a}{b}+\RatClass{c}{d}
      \leq\RatClass{a}{b}+\RatClass{e}{f}}{\rIE{6,10,17,19}}
  \proofstepwidestarbreakkeep[5]{q\leq r\leftrightarrow p+q\leq p+r}{%
    \rIE{5,20}}
  \proofstepwidestarbreakkeep[1]{q\leq r\leftrightarrow p+q\leq p+r}{%
    \rEE{2,3,4,5,21}}
\end{tabproofwide}

\FormulaThmDeltaK[Produkt nichtnegativer rationaler Zahlen ist nichtnegativ]{%
  q,r\in\mathbb Q\dsep\RatZero\leq q\dsep\RatZero\leq r
  \vdash\RatZero\leq q\cdot r%
}{RationalNonnegativeProduct}{%
  \DeltaRow{Mengen}{q\dsep r}%
}
\begin{tabproofwide}
  \proofstepwidestarbreakkeep[1]{q,r\in\mathbb Q}{\rA}
  \proofstepwidestarbreakkeep[2]{\RatZero\leq q}{\rA}
  \proofstepwidestarbreakkeep[3]{\RatZero\leq r}{\rA}
  \proofstepwidestarbreakkeep[1]{\exists a\in\mathbb Z\,
    \exists b\in\PositiveIntegers\,q=\RatClass{a}{b}}{%
    \FormulaRefAuto{RationalPairRepresentative}[thm]{\rAEa{1}}}
  \proofstepwidestarbreakkeep[1]{\exists c\in\mathbb Z\,
    \exists d\in\PositiveIntegers\,r=\RatClass{c}{d}}{%
    \FormulaRefAuto{RationalPairRepresentative}[thm]{\rAEb{1}}}
  \proofstepwidestarbreakkeep[6]{a,c\in\mathbb Z\land b,d\in\PositiveIntegers
    \land q=\RatClass{a}{b}\land r=\RatClass{c}{d}}{\rA}
  \proofstepwidestarbreakkeep[2,6]{\RatZero\leq\RatClass{a}{b}}{%
    \rIE{2,\rAEa{\rAEb{\rAEb{6}}}}}
  \proofstepwidestarbreakkeep[2,6]{\IntZero\leq a}{%
    \FormulaRefAuto{RationalOrderClassChar}[thm]{6,7},%
    \FormulaRefAuto{RationalZeroOneClasses}[thm],%
    \FormulaRefAuto{IntegerMulOneAxiom}[ax]{6},%
    \FormulaRefAuto{IntegerMulZeroFromAxioms}[thm]{6}}
  \proofstepwidestarbreakkeep[3,6]{\RatZero\leq\RatClass{c}{d}}{%
    \rIE{3,\rAEb{\rAEb{\rAEb{6}}}}}
  \proofstepwidestarbreakkeep[3,6]{\IntZero\leq c}{%
    \FormulaRefAuto{RationalOrderClassChar}[thm]{6,9},%
    \FormulaRefAuto{RationalZeroOneClasses}[thm],%
    \FormulaRefAuto{IntegerMulOneAxiom}[ax]{6},%
    \FormulaRefAuto{IntegerMulZeroFromAxioms}[thm]{6}}
  \proofstepwidestar[]{\IntZero\in\mathbb Z}{%
      \FormulaRefAuto{IntegerZeroCarrierAxiom}[ax]}
  \proofstepwidestarbreakkeep[2,3,6]{\IntZero\leq a\cdot c}{%
    \FormulaRefAuto{IntegerPositiveProductAxiom}[ax]{6},%
    \FormulaRefAuto{IntegerStrictOrderDef}[def]{6},%
    \FormulaRefAuto{IntegerMulZeroFromAxioms}[thm]{6},%
    \FormulaRefAuto{IntegerOrderReflexiveAxiom}[ax]{%
      11}}
  \proofstepwidestarbreakkeep[6]{b d\in\PositiveIntegers}{%
    \FormulaRefAuto{PositiveIntegerMulClosure}[thm]{6}}
  \proofstepwidestarbreakkeep[2,3,6]{\RatZero
    \leq\RatClass{a c}{b d}}{%
    \FormulaRefAuto{RationalOrderClassChar}[thm]{6,12,13},%
    \FormulaRefAuto{RationalZeroOneClasses}[thm],%
    \FormulaRefAuto{IntegerMulOneAxiom}[ax]{6},%
    \FormulaRefAuto{IntegerMulZeroFromAxioms}[thm]{6}}
  \proofstepwidestar[6]{\RatClass{a}{b}\cdot\RatClass{c}{d}=\RatClass{a c}{b d}}{%
      \FormulaRefAuto{RationalMultiplicationClassValue}[thm]{6}}
  \proofstepwidestarbreakkeep[6]{q\cdot r=\RatClass{a c}{b d}}{%
    \rIE{\rAEa{\rAEb{\rAEb{6}}},\rAEb{\rAEb{\rAEb{6}}},%
      15}}
  \proofstepwidestarbreakkeep[2,3,6]{\RatZero\leq q\cdot r}{\rIE{14,16}}
  \proofstepwidestarbreakkeep[1,2,3]{\RatZero\leq q\cdot r}{%
    \rEE{4,5,6,17}}
\end{tabproofwide}

\FormulaThmDeltaK[Positive rationale Faktoren bewahren die strikte Ordnung]{%
  p,q,r\in\mathbb Q\dsep q<r\dsep\RatZero<p
  \vdash p\cdot q<p\cdot r%
}{RationalPositiveMulStrictOrder}{%
  \DeltaRow{Mengen}{p\dsep q\dsep r}%
}
\begin{tabproofwide}
  \proofstepwidestarbreakkeep[1]{p,q,r\in\mathbb Q}{\rA}
  \proofstepwidestarbreakkeep[2]{q<r}{\rA}
  \proofstepwidestarbreakkeep[3]{\RatZero<p}{\rA}
  \proofstepwidestarbreakkeep[1]{\exists a\in\mathbb Z\,
    \exists b\in\PositiveIntegers\,p=\RatClass{a}{b}}{%
    \FormulaRefAuto{RationalPairRepresentative}[thm]{\rAEa{1}}}
  \proofstepwidestarbreakkeep[1]{\exists c\in\mathbb Z\,
    \exists d\in\PositiveIntegers\,q=\RatClass{c}{d}}{%
    \FormulaRefAuto{RationalPairRepresentative}[thm]{\rAEa{\rAEb{1}}}}
  \proofstepwidestarbreakkeep[1]{\exists e\in\mathbb Z\,
    \exists f\in\PositiveIntegers\,r=\RatClass{e}{f}}{%
    \FormulaRefAuto{RationalPairRepresentative}[thm]{\rAEb{\rAEb{1}}}}
  \proofstepwidestarbreakkeep[7]{a,c,e\in\mathbb Z\land
    b,d,f\in\PositiveIntegers\land p=\RatClass{a}{b}\land
    q=\RatClass{c}{d}\land r=\RatClass{e}{f}}{\rA}
  \proofstepwidestarbreakkeep[2,3,7]{\IntZero<a
    \land c\cdot f<e\cdot d}{%
    \FormulaRefAuto{RationalStrictOrderClassChar}[thm]{7},%
    \FormulaRefAuto{RationalZeroOneClasses}[thm],%
    \FormulaRefAuto{IntegerMulOneAxiom}[ax]{7},%
    \FormulaRefAuto{IntegerMulZeroFromAxioms}[thm]{7}}
  \proofstepwidestarbreakkeep[2,3,7]{(a\cdot c)\cdot(b\cdot f)
    <(a\cdot e)\cdot(b\cdot d)}{%
    \FormulaRefAuto{IntegerPositiveFactorStrictOrderFromAxioms}[thm]{7,8},%
    \FormulaRefAuto{IntegerPositiveProductAxiom}[ax]{7},%
    \FormulaRefAuto{IntegerMulAssociativeAxiom}[ax]{7},%
    \FormulaRefAuto{IntegerMulCommutativeAxiom}[ax]{7}}
  \proofstepwidestarbreakkeep[2,3,7]{%
    \RatClass{a c}{b d}<\RatClass{a e}{b f}}{%
    \FormulaRefAuto{RationalStrictOrderClassChar}[thm]{7,9}}
  \proofstepwidestarbreakkeep[7]{p\cdot q=\RatClass{a c}{b d}
    \land p\cdot r=\RatClass{a e}{b f}}{%
    \rAIRepeat{2}{%
      \FormulaRefAuto{RationalMultiplicationClassValue}[thm]}{7}}
  \proofstepwidestarbreakkeep[2,3,7]{p\cdot q<p\cdot r}{\rIE{10,11}}
  \proofstepwidestarbreakkeep[1,2,3]{p\cdot q<p\cdot r}{%
    \rEE{4,5,6,7,12}}
\end{tabproofwide}

\FormulaThmDeltaK[Positive rationale Zahlen haben positive Kehrwerte]{%
  q\in\mathbb Q\dsep\RatZero<q
  \vdash\RatZero<q^{-1}%
}{RationalPositiveReciprocal}{%
  \DeltaRow{Mengen}{q\dsep a\dsep b}%
  \DeltaRow{Relationsoperatoren}{<}%
  \DeltaRow{Funktionsoperatoren}{(\,\cdot\,)^{-1}}%
}
\begin{tabproofwide}
  \proofstepwidestarbreakkeep[1]{q\in\mathbb Q}{\rA}
  \proofstepwidestarbreakkeep[2]{\RatZero<q}{\rA}
  \proofstepwidestarbreakkeep[1]{\exists a\in\mathbb Z\,
    \exists b\in\PositiveIntegers\,q=\RatClass{a}{b}}{%
    \FormulaRefAuto{RationalPairRepresentative}[thm]{1}}
  \proofstepwidestarbreakkeep[4]{a\in\mathbb Z\land b\in\PositiveIntegers
    \land q=\RatClass{a}{b}}{\rA}
  \proofstepwidestarbreakkeep[2,4]{\RatZero<\RatClass{a}{b}}{%
    \rIE{2,\rAEb{\rAEb{4}}}}
  \proofstepwidestarbreakkeep[2,4]{\IntZero\cdot b<a\cdot\IntOne}{%
    \FormulaRefAuto{RationalStrictOrderClassChar}[thm]{4,5},%
    \FormulaRefAuto{RationalZeroOneClasses}[thm]}
  \proofstepwidestar[4]{a\cdot\IntOne=a\land\IntOne\cdot a=a}{%
      \FormulaRefAuto{IntegerMulOneAxiom}[ax]{4}}
  \proofstepwidestar[4]{b\cdot\IntZero=\IntZero\land\IntZero\cdot b=\IntZero}{%
      \FormulaRefAuto{IntegerMulZeroFromAxioms}[thm]{4}}
  \proofstepwidestarbreakkeep[2,4]{\IntZero<a}{%
    \rIE{6,7,%
      8}}
  \proofstepwidestar[2,4]{\IntZero\leq a\land\IntZero\neq a}{%
      \FormulaRefAuto{IntegerStrictOrderDef}[def]{4,9}}
  \proofstepwidestarbreakkeep[2,4]{a\neq\IntZero}{\FormulaRefAuto{a\neq b\vdash b\neq a}[thm]{\rAEb{10}}}
  \proofstepwidestar[2,4]{(\RatClass{a}{b})^{-1}=\RatRecipRep{a}{b}}{%
      \FormulaRefAuto{RationalReciprocalClassValue}[thm]{4,11}}
  \proofstepwidestar[2,4]{\RatRecipRep{a}{b}=\RatClass{b}{a}}{%
      \FormulaRefAuto{RationalReciprocalRepresentativeDef}[def]{4,9}}
  \proofstepwidestarbreakkeep[2,4]{q^{-1}=\RatClass{b}{a}}{%
    \rIE{\rAEb{\rAEb{4}},%
      12,%
      13}}
  \proofstepwidestar[4]{b\in\mathbb Z\land\IntZero<b}{%
      \FormulaRefAuto{PositiveIntegersDef}[def]{\rAEa{\rAEb{4}}}}
  \proofstepwidestarbreakkeep[4]{\IntZero<b}{%
    \rAEb{15}}
  \proofstepwidestar[4]{b\cdot\IntOne=b\land\IntOne\cdot b=b}{%
      \FormulaRefAuto{IntegerMulOneAxiom}[ax]{4}}
  \proofstepwidestar[4]{a\cdot\IntZero=\IntZero\land\IntZero\cdot a=\IntZero}{%
      \FormulaRefAuto{IntegerMulZeroFromAxioms}[thm]{4}}
  \proofstepwidestarbreakkeep[2,4]{\IntZero\cdot a<b\cdot\IntOne}{%
    \rIE{16,17,%
      18}}
  \proofstepwidestarbreakkeep[2,4]{\RatZero<\RatClass{b}{a}}{%
    \FormulaRefAuto{RationalStrictOrderClassChar}[thm]{4,9,19},%
    \FormulaRefAuto{RationalZeroOneClasses}[thm]}
  \proofstepwidestarbreakkeep[2,4]{\RatZero<q^{-1}}{\rIE{14,20}}
  \proofstepwidestarbreakkeep[1,2]{\RatZero<q^{-1}}{\rEE{3,4,21}}
\end{tabproofwide}

\FormulaThmDeltaK[Ordnungstreue der ganzzahligen Einbettung]{%
  z,w\in\mathbb Z\vdash
  \bigl(z\leq w\leftrightarrow\RatEmbed(z)\leq\RatEmbed(w)\bigr)%
}{IntegerRationalEmbeddingOrder}{%
  \DeltaRow{Mengen}{z\dsep w}%
  \DeltaRow{Funktion}{\RatEmbed}%
}
\begin{tabproofwide}
  \proofstepwidestarbreakkeep[1]{z,w\in\mathbb Z}{\rA}
  \proofstepwidestarbreakkeep[1]{\RatEmbed(z)=\RatClass{z}{\IntOne}
    \land\RatEmbed(w)=\RatClass{w}{\IntOne}}{%
    \rAIRepeat{2}{\FormulaRefAuto{IntegerRationalEmbedding}[def]}{%
      \rAEa{1}\mid\rAEb{1}}}
  \proofstepwidestar[1]{\IntZero < \IntOne}{%
      \FormulaRefAuto{IntegerOnePositiveAxiom}[ax]}
  \proofstepwidestarbreakkeep[1]{\RatClass{z}{\IntOne}
    \leq\RatClass{w}{\IntOne}
    \leftrightarrow z\cdot\IntOne\leq w\cdot\IntOne}{%
    \FormulaRefAuto{RationalOrderClassChar}[thm]{%
      1,3}}
  \proofstepwidestarbreakkeep[1]{z\cdot\IntOne\leq w\cdot\IntOne
    \leftrightarrow z\leq w}{%
    \FormulaRefAuto{IntegerMulOneAxiom}[ax]{1}}
  \proofstepwidestarbreakkeep[1]{z\leq w
    \leftrightarrow\RatEmbed(z)\leq\RatEmbed(w)}{\rIE{2,4,5}}
\end{tabproofwide}

\chapter{Axiomatische Weiterarbeit}

Der Quotientenkern ist damit vollständig verifiziert.  Ab jetzt werden keine
Eigenschaften von Bruchpaaren mehr vorausgesetzt.  Die folgenden drei
Modellsätze halten noch die für \(\mathbb Q\) charakteristische Erzeugung aus
ganzen Zählern und positiven ganzen Nennern sowie die Orientierung der Eins
fest; anschließend wird ausschließlich über die registrierten Axiome
weitergearbeitet.

\FormulaThmDeltaK[Null und Eins sind im rationalen Modell verschieden]{%
  \RatZero\neq\RatOne%
}{RationalZeroNeOneModel}{%
  \DeltaRow{Mengen}{\mathbb Q}%
}
\begin{tabproof}
  \proofstep{}{\IntZero\neq\IntOne}{%
    \FormulaRefAuto{IntegerZeroNeOneAxiom}[ax]}
  \proofstep{}{\IntOne \in \mathbb Z}{%
      \FormulaRefAuto{IntegerOneCarrierAxiom}[ax]}
  \proofstep{}{\IntZero \in \mathbb Z}{%
      \FormulaRefAuto{IntegerZeroCarrierAxiom}[ax]}
  \proofstep{}{\RatEmbed(\IntZero)\neq\RatEmbed(\IntOne)}{%
    \FormulaRefAuto{IntegerRationalEmbeddingInjective}[thm]{%
      3,%
      2,1}}
  \proofstep{}{\RatOne = \RatOfInt \IntOne}{%
      \FormulaRefAuto{RationalOneDef}[def]}
  \proofstep{}{\RatZero = \RatOfInt \IntZero}{%
      \FormulaRefAuto{RationalZeroDef}[def]}
  \proofstep{}{\RatZero\neq\RatOne}{%
    \rIE{4,6,%
      5}}
\end{tabproof}

\FormulaThmDeltaK[Die rationale Eins ist im Modell positiv]{%
  \RatZero<\RatOne%
}{RationalOnePositiveModel}{%
  \DeltaRow{Mengen}{\mathbb Q}%
}
\begin{tabproof}
  \proofstep{}{\IntZero<\IntOne}{%
    \FormulaRefAuto{IntegerOnePositiveAxiom}[ax]}
  \proofstep{}{\IntOne\in\mathbb Z}{%
      \FormulaRefAuto{IntegerOneCarrierAxiom}[ax]}
  \proofstep{}{\IntZero\in\mathbb Z}{%
      \FormulaRefAuto{IntegerZeroCarrierAxiom}[ax]}
  \proofstep{}{\RatEmbed(\IntZero)<\RatEmbed(\IntOne)}{%
    \FormulaRefAuto{IntegerRationalEmbeddingOrder}[thm]{%
      3,%
      2,1},%
    \FormulaRefAuto{RationalStrictOrderDef}[def]}
  \proofstep{}{\RatOne = \RatOfInt \IntOne}{%
      \FormulaRefAuto{RationalOneDef}[def]}
  \proofstep{}{\RatZero = \RatOfInt \IntZero}{%
      \FormulaRefAuto{RationalZeroDef}[def]}
  \proofstep{}{\RatZero<\RatOne}{%
    \rIE{4,6,%
      5}}
\end{tabproof}

\FormulaThmDeltaK[Brucherzeugung des rationalen Modells]{%
  q\in\mathbb Q\vdash
  \exists a\in\mathbb Z\,\exists b\in\PositiveIntegers\,
  q=\RatEmbed(a)\cdot\bigl(\RatEmbed(b)\bigr)^{-1}%
}{RationalFractionGenerationModel}{%
  \DeltaRow{Mengen}{q\dsep a\dsep b}%
  \DeltaRow{Funktion}{\RatEmbed}%
}
\begin{tabproofwide}
  \proofstepwidestarbreakkeep[1]{q\in\mathbb Q}{\rA}
  \proofstepwidestarbreakkeep[1]{\exists a\in\mathbb Z\,
    \exists b\in\PositiveIntegers\,q=\RatClass{a}{b}}{%
    \FormulaRefAuto{RationalPairRepresentative}[thm]{1}}
  \proofstepwidestarbreakkeep[3]{a\in\mathbb Z\land b\in\PositiveIntegers
    \land q=\RatClass{a}{b}}{\rA}
  \proofstepwidestar[3]{\RatEmbed(a)=\RatClass{a}{\IntOne}}{%
      \FormulaRefAuto{IntegerRationalEmbedding}[def]{3}}
  \proofstepwidestar[3]{b\in\mathbb Z\land\IntZero<b}{%
      \FormulaRefAuto{PositiveIntegersDef}[def]{\rAEa{\rAEb{3}}}}
  \proofstepwidestar[3]{\RatEmbed(b)=\RatClass{b}{\IntOne}}{%
      \FormulaRefAuto{IntegerRationalEmbedding}[def]{%
        \rAEa{5}}}
  \proofstepwidestarbreakkeep[3]{\RatEmbed(a)=\RatClass{a}{\IntOne}
    \land\RatEmbed(b)=\RatClass{b}{\IntOne}}{%
    \rAI{4,%
      6}}
  \proofstepwidestarbreakkeep[3]{b\neq\IntZero}{%
    \FormulaRefAuto{PositiveIntegerNonzero}[thm]{\rAEa{\rAEb{3}}}}
  \proofstepwidestarbreakkeep[3]{\bigl(\RatEmbed(b)\bigr)^{-1}
    =\RatClass{\IntOne}{b}}{%
    \FormulaRefAuto{RationalReciprocalClassValue}[thm]{3,8},%
    \FormulaRefAuto{RationalReciprocalRepresentativeDef}[def]{3}}
  \proofstepwidestar[3]{a\cdot\IntOne=a\land\IntOne\cdot b=b}{%
      \FormulaRefAuto{IntegerMulOneAxiom}[ax]{3}}
  \proofstepwidestar[3]{\RatClass{a}{\IntOne}\cdot\RatClass{\IntOne}{b}=\RatClass{a\cdot\IntOne}{\IntOne\cdot b}}{%
      \FormulaRefAuto{RationalMultiplicationClassValue}[thm]{3}}
  \proofstepwidestarbreakkeep[3]{\RatEmbed(a)\cdot
    \bigl(\RatEmbed(b)\bigr)^{-1}=\RatClass{a}{b}}{%
    \rIE{7,9,11,%
      10}}
  \proofstepwidestarbreakkeep[3]{q=\RatEmbed(a)\cdot
    \bigl(\RatEmbed(b)\bigr)^{-1}}{\rIE{\rAEb{\rAEb{3}},12}}
  \proofstepwidestarbreakkeep[3]{\exists a\in\mathbb Z\,
    \exists b\in\PositiveIntegers\,
    q=\RatEmbed(a)\cdot\bigl(\RatEmbed(b)\bigr)^{-1}}{%
    \rEI{\rAI{\rAEa{3},\rEI{\rAI{\rAEa{\rAEb{3}},13}}}}}
  \proofstepwidestarbreakkeep[1]{\exists a\in\mathbb Z\,
    \exists b\in\PositiveIntegers\,
    q=\RatEmbed(a)\cdot\bigl(\RatEmbed(b)\bigr)^{-1}}{\rEE{2,3,14}}
\end{tabproofwide}

Die folgende Tafel ist die Schnittstelle für alle weiteren Sätze dieses und
der folgenden Bände.  Anders als die Axiome eines beliebigen geordneten
Körpers enthält sie ausdrücklich die Brucherzeugung; erst diese Zeile hält
fest, dass der Träger gerade der von \(\mathbb Z\) erzeugte Primkörper ist.
Dichte und Archimedizität werden danach als Folgesätze bewiesen. Die beiden
abschließenden Ganzzahlgrundlagen werden nur als importierte, stabil
registrierte Hilfsregeln mitgeführt; sie sind keine zusätzlichen
Körperaxiome der rationalen Zahlen.

\FormulaDefDeltaK[Rationale Zahlen als erzeugter geordneter Körper]{%
  \RationalOrderedField{\mathbb Q}%
}{RationalOrderedFieldAxioms}{%
  \DeltaRow{\textbf{Signatur}}{}%
  \DeltaRow{Träger}{\mathbb Q}%
  \DeltaRow{Konstanten}{\RatZero\dsep\RatOne}%
  \DeltaRow{Totale Operationen}{+\dsep-\dsep\cdot}%
  \DeltaRow{Partieller Kehrwert}{%
    q\longmapsto q^{-1}\quad
    (q\in\mathbb Q\setminus\{\RatZero\})}%
  \DeltaRow{Ordnung}{\leq_{\mathbb Q}}%
  \DeltaRow{Einbettung}{\RatEmbed}%
  \DeltaRow{Null}{\RatZero\in\mathbb Q}
    [\FormulaRefAuto{RationalZeroCarrierAxiom}[ax]]%
  \DeltaRow{Eins}{\RatOne\in\mathbb Q}
    [\FormulaRefAuto{RationalOneCarrierAxiom}[ax]]%
  \DeltaRow{Verschieden}{\RatZero\neq\RatOne}
    [\FormulaRefAuto{RationalZeroNeOneAxiom}[ax]]%
  \end{DeltaContext}
  \begin{DeltaContext}{ }%
  \DeltaRow{\makecell[l]{\textbf{Additive abelsche}\\\textbf{Gruppe}}}{}%
  \DeltaRow{Abgeschlossen}{q,r\in\mathbb Q\vdash q+r\in\mathbb Q}
    [\FormulaRefAuto{RationalAddClosureAxiom}[ax]]%
  \DeltaRow{Assoziativ}{\begin{aligned}[t]
    &p,q,r\in\mathbb Q\vdash{}\\[-2pt]
    &(p+q)+r=p+(q+r)
    \end{aligned}}
    [\FormulaRefAuto{RationalAddAssociativeAxiom}[ax]]%
  \DeltaRow{Kommutativ}{q,r\in\mathbb Q\vdash q+r=r+q}
    [\FormulaRefAuto{RationalAddCommutativeAxiom}[ax]]%
  \DeltaRow{Negation}{q\in\mathbb Q\vdash-q\in\mathbb Q}
    [\FormulaRefAuto{RationalNegationClosureAxiom}[ax]]%
  \DeltaRow{Nullgesetz}{q\in\mathbb Q\vdash q+\RatZero=q}
    [\FormulaRefAuto{RationalAddZeroAxiom}[ax]]%
  \DeltaRow{Inverse}{q\in\mathbb Q\vdash q+(-q)=\RatZero}
    [\FormulaRefAuto{RationalAddInverseAxiom}[ax]]%
  \end{DeltaContext}
  \begin{DeltaContext}{ }%
  \DeltaRow{\makecell[l]{\textbf{Kommutatives Monoid}\\
    \textbf{und Kehrwertaxiome}}}{}%
  \DeltaRow{Abgeschlossen}{q,r\in\mathbb Q\vdash q\cdot r\in\mathbb Q}
    [\FormulaRefAuto{RationalMulClosureAxiom}[ax]]%
  \DeltaRow{Assoziativ}{\begin{aligned}[t]
    &p,q,r\in\mathbb Q\vdash{}\\[-2pt]
    &(p\cdot q)\cdot r=p\cdot(q\cdot r)
    \end{aligned}}
    [\FormulaRefAuto{RationalMulAssociativeAxiom}[ax]]%
  \DeltaRow{Kommutativ}{q,r\in\mathbb Q\vdash q\cdot r=r\cdot q}
    [\FormulaRefAuto{RationalMulCommutativeAxiom}[ax]]%
  \DeltaRow{Einsgesetz}{q\in\mathbb Q\vdash q\cdot\RatOne=q}
    [\FormulaRefAuto{RationalMulOneAxiom}[ax]]%
  \DeltaRow{Kehrwert}{q\in\mathbb Q\dsep q\neq\RatZero
    \vdash q^{-1}\in\mathbb Q}
    [\FormulaRefAuto{RationalReciprocalClosureAxiom}[ax]]%
  \DeltaRow{Inverse}{q\in\mathbb Q\dsep q\neq\RatZero
    \vdash q\cdot q^{-1}=\RatOne}
    [\FormulaRefAuto{RationalMulInverseAxiom}[ax]]%
  \end{DeltaContext}
  \begin{DeltaContext}{ }%
  \DeltaRow{\textbf{Distributivität}}{}%
  \DeltaRow{Distributiv}{\begin{aligned}[t]
    &p,q,r\in\mathbb Q\vdash{}\\[-2pt]
    &p\cdot(q+r)=p\cdot q+p\cdot r
    \end{aligned}}
    [\FormulaRefAuto{RationalDistributiveAxiom}[ax]]%
  \end{DeltaContext}
  \begin{DeltaContext}{ }%
  \DeltaRow{\makecell[l]{\textbf{Geordneter}\\\textbf{Körper}}}{}%
  \DeltaRow{Orientierung}{\RatZero<\RatOne}
    [\FormulaRefAuto{RationalOnePositiveAxiom}[ax]]%
  \DeltaRow{Total}{\TotOrd{\mathbb Q,\leq_{\mathbb Q}}}
    [\FormulaRefAuto{RationalOrderTotalAxiom}[ax]]%
  \DeltaRow{Translation}{\begin{aligned}[t]
    &p,q,r\in\mathbb Q\vdash{}\\[-2pt]
    &q\leq r\leftrightarrow p+q\leq p+r
    \end{aligned}}
    [\FormulaRefAuto{RationalOrderTranslationAxiom}[ax]]%
  \DeltaRow{\makecell[l]{Nichtnegative\\Produkte}}{\begin{aligned}[t]
    &q,r\in\mathbb Q\dsep\RatZero\leq q\dsep{}\\[-2pt]
    &\RatZero\leq r\vdash\RatZero\leq q\cdot r
    \end{aligned}}
    [\FormulaRefAuto{RationalNonnegativeProductAxiom}[ax]]%
  \end{DeltaContext}
  \begin{DeltaContext}{ }%
  \DeltaRow{\makecell[l]{\textbf{Abgeleitete}\\\textbf{Ordnungsregeln}}}{}%
  \DeltaRow{\makecell[l]{Positive\\Faktoren}}{\begin{aligned}[t]
    &p,q,r\in\mathbb Q\dsep q<r\dsep\RatZero<p\vdash{}\\[-2pt]
    &p\cdot q<p\cdot r
    \end{aligned}}
    [\FormulaRefAuto{RationalPositiveMulStrictOrderAxiom}[ax]]%
  \DeltaRow{\makecell[l]{Positive\\Kehrwerte}}{\begin{aligned}[t]
    &q\in\mathbb Q\dsep\RatZero<q\vdash{}\\[-2pt]
    &\RatZero<q^{-1}
    \end{aligned}}
    [\FormulaRefAuto{RationalPositiveReciprocalAxiom}[ax]]%
  \end{DeltaContext}
  \begin{DeltaContext}{ }%
  \DeltaRow{\makecell[l]{\textbf{Ganzzahleinbettung}\\
    \textbf{und Brucherzeugung}}}{}%
  \DeltaRow{Funktion}{\RatEmbed\colon\mathbb Z\to\mathbb Q}
    [\FormulaRefAuto{RationalIntegerEmbeddingFunctionAxiom}[ax]]%
  \DeltaRow{Injektiv}{\RatEmbed\colon\mathbb Z\inj\mathbb Q}
    [\FormulaRefAuto{RationalIntegerEmbeddingInjectiveAxiom}[ax]]%
  \DeltaRow{Additiv}{\begin{aligned}[t]
    &z,w\in\mathbb Z\vdash{}\\[-2pt]
    &\RatEmbed(z+w)=\RatEmbed(z)+\RatEmbed(w)
    \end{aligned}}
    [\FormulaRefAuto{RationalIntegerEmbeddingAdditiveAxiom}[ax]]%
  \DeltaRow{Multiplikativ}{\begin{aligned}[t]
    &z,w\in\mathbb Z\vdash{}\\[-2pt]
    &\RatEmbed(z\cdot w)=\RatEmbed(z)\cdot\RatEmbed(w)
    \end{aligned}}
    [\FormulaRefAuto{RationalIntegerEmbeddingMultiplicativeAxiom}[ax]]%
  \DeltaRow{Ordnungstreu}{\begin{aligned}[t]
    &z,w\in\mathbb Z\vdash{}\\[-2pt]
    &z\leq w\leftrightarrow\RatEmbed(z)\leq\RatEmbed(w)
    \end{aligned}}
    [\FormulaRefAuto{RationalIntegerEmbeddingOrderAxiom}[ax]]%
  \DeltaRow{Brucherzeugung}{\begin{aligned}[t]
    &q\in\mathbb Q\vdash
      \exists a\in\mathbb Z\,\exists b\in\PositiveIntegers\,{}\\[-2pt]
    &q=\RatEmbed(a)\cdot(\RatEmbed(b))^{-1}
  \end{aligned}}
    [\FormulaRefAuto{RationalFractionGenerationAxiom}[ax]]%
  \DeltaRow{\makecell[l]{\textbf{Importierte}\\
    \textbf{Ganzzahlgrundlagen}}}{}%
  \DeltaRow{Positive Nenner}{\begin{aligned}[t]
    &b\in\PositiveIntegers\vdash{}\\[-2pt]
    &b\in\mathbb Z\land\IntZero<b
    \end{aligned}}
    [\FormulaRefAuto{RationalPositiveDenominatorAxiom}[ax]]%
  \DeltaRow{Nenner-Schranke}{\begin{aligned}[t]
    &a\in\mathbb Z\dsep b\in\PositiveIntegers\vdash{}\\[-2pt]
    &\exists n\in\mathbb N\,a<\IntEmbed(n)\cdot b
    \end{aligned}}
    [\FormulaRefAuto{RationalIntegerFractionBoundAxiom}[ax]]%
}

\FormulaThmDeltaK[Das Quotientenmodell erfüllt die rationalen Körperaxiome]{%
  \begin{aligned}[t]
    &\mathbb Q=\QuotientSet{\RatPairCarrier}{\RatPairRel}\ \land{}\\
    &\RationalOrderedField{\mathbb Q}
  \end{aligned}%
}{RationalQuotientModelSatisfiesOrderedField}{%
  \DeltaRow{Mengen}{\mathbb Q}%
  \DeltaRow{Funktionensymbol}{\RatEmbed}%
  \DeltaRow{Relationsoperatoren}{\leq_{\mathbb Q}}%
}
\begin{tabproofwide}
  \proofstepwidestarbreakkeep[]{\mathbb Q
    =\QuotientSet{\RatPairCarrier}{\RatPairRel}}{%
    \FormulaRefAuto{RationalsAsQuotient}[def]}
  \proofstepwidestarbreakkeep[]{\RationalOrderedField{\mathbb Q}}{\begin{gathered}[t]
      \FormulaRefAuto{RationalOrderedFieldAxioms}[def]
      \\[-2pt]\FormulaRefAuto{RationalZeroOneCarrier}[thm]
      \\[-2pt]\FormulaRefAuto{RationalZeroNeOneModel}[thm]
      \\[-2pt]\FormulaRefAuto{RationalOnePositiveModel}[thm]
      \\[-2pt]\FormulaRefAuto{RationalNegationLaws}[thm]
      \\[-2pt]\FormulaRefAuto{RationalAddClosure}[thm]
      \\[-2pt]\FormulaRefAuto{RationalAddAssociative}[thm]
      \\[-2pt]\FormulaRefAuto{RationalAddCommutative}[thm]
      \\[-2pt]\FormulaRefAuto{RationalAddZero}[thm]
      \\[-2pt]\FormulaRefAuto{RationalAddInverse}[thm]
      \\[-2pt]\FormulaRefAuto{RationalMulClosure}[thm]
      \\[-2pt]\FormulaRefAuto{RationalMulAssociative}[thm]
      \\[-2pt]\FormulaRefAuto{RationalMulCommutative}[thm]
      \\[-2pt]\FormulaRefAuto{RationalMulOne}[thm]
      \\[-2pt]\FormulaRefAuto{RationalReciprocalClosure}[thm]
      \\[-2pt]\FormulaRefAuto{RationalMulInverse}[thm]
      \\[-2pt]\FormulaRefAuto{RationalDistributive}[thm]
      \\[-2pt]\FormulaRefAuto{RationalOrderTotal}[thm]
      \\[-2pt]\FormulaRefAuto{RationalOrderTranslation}[thm]
      \\[-2pt]\FormulaRefAuto{RationalNonnegativeProduct}[thm]
      \\[-2pt]\FormulaRefAuto{RationalPositiveMulStrictOrder}[thm]
      \\[-2pt]\FormulaRefAuto{RationalPositiveReciprocal}[thm]
      \\[-2pt]\FormulaRefAuto{IntegerRationalEmbeddingFunction}[thm]
      \\[-2pt]\FormulaRefAuto{IntegerRationalEmbeddingInjective}[thm]
      \\[-2pt]\FormulaRefAuto{IntegerRationalEmbeddingAdditive}[thm]
      \\[-2pt]\FormulaRefAuto{IntegerRationalEmbeddingMultiplicative}[thm]
      \\[-2pt]\FormulaRefAuto{IntegerRationalEmbeddingOrder}[thm]
      \\[-2pt]\FormulaRefAuto{RationalFractionGenerationModel}[thm]
      \\[-2pt]\FormulaRefAuto{PositiveIntegersDef}[def]
      \\[-2pt]\FormulaRefAuto{RationalIntegerFractionBoundModel}[thm]
      \end{gathered}}
  \proofstepwidestarbreakkeep[]{\mathbb Q
    =\QuotientSet{\RatPairCarrier}{\RatPairRel}
    \land\RationalOrderedField{\mathbb Q}}{\rAI{1,2}}
\end{tabproofwide}

\medskip\noindent\textbf{Signatur und Grundkonstanten.}\par

\FormulaAxiomAutoR[Rationale Null im Träger]{%
  \RatZero\in\mathbb Q}{RationalZeroCarrierAxiom}
\par\noindent\emph{Modellursprung:}
\FormulaRefAuto{RationalZeroOneCarrier}[thm].\par

\FormulaAxiomAutoR[Rationale Eins im Träger]{%
  \RatOne\in\mathbb Q}{RationalOneCarrierAxiom}
\par\noindent\emph{Modellursprung:}
\FormulaRefAuto{RationalZeroOneCarrier}[thm].\par

\FormulaAxiomAutoR[Rationale Null und Eins sind verschieden]{%
  \RatZero\neq\RatOne}{RationalZeroNeOneAxiom}
\par\noindent\emph{Modellursprung:}
\FormulaRefAuto{RationalZeroNeOneModel}[thm].\par

\medskip\noindent\textbf{Additive abelsche Gruppenstruktur.}\par

\FormulaAxiomAutoR[Abgeschlossenheit der rationalen Addition]{%
  q,r\in\mathbb Q\vdash q+r\in\mathbb Q}{RationalAddClosureAxiom}
\FormulaAxiomAutoR[Assoziativität der rationalen Addition]{%
  p,q,r\in\mathbb Q\vdash(p+q)+r=p+(q+r)}{RationalAddAssociativeAxiom}
\FormulaAxiomAutoR[Kommutativität der rationalen Addition]{%
  q,r\in\mathbb Q\vdash q+r=r+q}{RationalAddCommutativeAxiom}
\FormulaAxiomAutoR[Abgeschlossenheit der rationalen Negation]{%
  q\in\mathbb Q\vdash-q\in\mathbb Q}{RationalNegationClosureAxiom}
\FormulaAxiomAutoR[Nullgesetz der rationalen Addition]{%
  q\in\mathbb Q\vdash q+\RatZero=q}{RationalAddZeroAxiom}
\FormulaAxiomAutoR[Additive Inversen in \(\mathbb Q\)]{%
  q\in\mathbb Q\vdash q+(-q)=\RatZero}{RationalAddInverseAxiom}

\medskip\noindent\textbf{Kommutatives multiplikatives Monoid und
Kehrwertaxiome.}\par

\FormulaAxiomAutoR[Abgeschlossenheit der rationalen Multiplikation]{%
  q,r\in\mathbb Q\vdash q\cdot r\in\mathbb Q}{RationalMulClosureAxiom}
\FormulaAxiomAutoR[Assoziativität der rationalen Multiplikation]{%
  p,q,r\in\mathbb Q\vdash(p\cdot q)\cdot r=p\cdot(q\cdot r)%
}{RationalMulAssociativeAxiom}
\FormulaAxiomAutoR[Kommutativität der rationalen Multiplikation]{%
  q,r\in\mathbb Q\vdash q\cdot r=r\cdot q}{RationalMulCommutativeAxiom}
\FormulaAxiomAutoR[Einheitsgesetz der rationalen Multiplikation]{%
  q\in\mathbb Q\vdash q\cdot\RatOne=q}{RationalMulOneAxiom}
\FormulaAxiomAutoR[Abgeschlossenheit des rationalen Kehrwertes]{%
  q\in\mathbb Q\dsep q\neq\RatZero\vdash q^{-1}\in\mathbb Q%
}{RationalReciprocalClosureAxiom}
\FormulaAxiomAutoR[Multiplikative Inversen in \(\mathbb Q\)]{%
  q\in\mathbb Q\dsep q\neq\RatZero
  \vdash q\cdot q^{-1}=\RatOne}{RationalMulInverseAxiom}

\medskip\noindent\textbf{Distributivität.}\par

\FormulaAxiomAutoR[Distributivität der rationalen Multiplikation]{%
  p,q,r\in\mathbb Q\vdash p\cdot(q+r)=p\cdot q+p\cdot r%
}{RationalDistributiveAxiom}

\medskip\noindent\textbf{Geordnete Körperstruktur.}\par

\FormulaAxiomAutoR[Die rationale Eins ist positiv]{%
  \RatZero<\RatOne}{RationalOnePositiveAxiom}
\par\noindent\emph{Modellursprung:}
\FormulaRefAuto{RationalOnePositiveModel}[thm].\par

\FormulaAxiomAutoR[Totale Ordnung der rationalen Zahlen]{%
  \TotOrd{\mathbb Q,\leq_{\mathbb Q}}}{RationalOrderTotalAxiom}
\FormulaAxiomAutoR[Translationsäquivalenz der rationalen Ordnung]{%
  p,q,r\in\mathbb Q\vdash(q\leq r\leftrightarrow p+q\leq p+r)%
}{RationalOrderTranslationAxiom}
\FormulaAxiomAutoR[Nichtnegative rationale Produkte]{%
  q,r\in\mathbb Q\dsep\RatZero\leq q\dsep\RatZero\leq r
  \vdash\RatZero\leq q\cdot r}{RationalNonnegativeProductAxiom}

\medskip\noindent\textbf{Abgeleitete Ordnungsregeln.}\par

\FormulaAxiomAutoR[Positive rationale Faktoren bewahren strikte Ordnung]{%
  p,q,r\in\mathbb Q\dsep q<r\dsep\RatZero<p
  \vdash p\cdot q<p\cdot r}{RationalPositiveMulStrictOrderAxiom}
\FormulaAxiomAutoR[Positive rationale Zahlen haben positive Kehrwerte]{%
  q\in\mathbb Q\dsep\RatZero<q\vdash\RatZero<q^{-1}%
}{RationalPositiveReciprocalAxiom}

\medskip\noindent\textbf{Ganzzahleinbettung und Brucherzeugung.}\par

\FormulaAxiomAutoR[Typisierung der Ganzzahleinbettung]{%
  \RatEmbed\colon\mathbb Z\to\mathbb Q%
}{RationalIntegerEmbeddingFunctionAxiom}
\FormulaAxiomAutoR[Injektivität der Ganzzahleinbettung]{%
  \RatEmbed\colon\mathbb Z\inj\mathbb Q%
}{RationalIntegerEmbeddingInjectiveAxiom}
\FormulaAxiomAutoR[Additivität der Ganzzahleinbettung]{%
  z,w\in\mathbb Z\vdash\RatEmbed(z+w)=\RatEmbed(z)+\RatEmbed(w)%
}{RationalIntegerEmbeddingAdditiveAxiom}
\FormulaAxiomAutoR[Multiplikativität der Ganzzahleinbettung]{%
  z,w\in\mathbb Z\vdash
  \RatEmbed(z\cdot w)=\RatEmbed(z)\cdot\RatEmbed(w)%
}{RationalIntegerEmbeddingMultiplicativeAxiom}
\FormulaAxiomAutoR[Ordnungstreue der Ganzzahleinbettung]{%
  z,w\in\mathbb Z\vdash
  (z\leq w\leftrightarrow\RatEmbed(z)\leq\RatEmbed(w))%
}{RationalIntegerEmbeddingOrderAxiom}
\FormulaAxiomAutoR[Brucherzeugung der rationalen Zahlen]{%
  q\in\mathbb Q\vdash
  \exists a\in\mathbb Z\,\exists b\in\PositiveIntegers\,
  q=\RatEmbed(a)\cdot(\RatEmbed(b))^{-1}%
}{RationalFractionGenerationAxiom}
\FormulaAxiomAutoR[Charakterisierung positiver Nenner]{%
  b\in\PositiveIntegers\vdash
  b\in\mathbb Z\land\IntZero<b%
}{RationalPositiveDenominatorAxiom}
\FormulaAxiomAutoR[Ganzzahlschranke für positive Nenner]{%
  a\in\mathbb Z\dsep b\in\PositiveIntegers\vdash
  \exists n\in\mathbb N\,a<\IntEmbed(n)\cdot b%
}{RationalIntegerFractionBoundAxiom}

\begin{sloppypar}
\noindent\emph{Modellursprünge der arithmetischen, geordneten und
Einbettungsaxiome:}
\FormulaRefAuto{RationalAddClosure}[thm],
\FormulaRefAuto{RationalAddAssociative}[thm],
\FormulaRefAuto{RationalAddCommutative}[thm],
\FormulaRefAuto{RationalNegationLaws}[thm],
\FormulaRefAuto{RationalAddZero}[thm],
\FormulaRefAuto{RationalAddInverse}[thm],
\FormulaRefAuto{RationalMulClosure}[thm],
\FormulaRefAuto{RationalMulAssociative}[thm],
\FormulaRefAuto{RationalMulCommutative}[thm],
\FormulaRefAuto{RationalMulOne}[thm],
\FormulaRefAuto{RationalReciprocalClosure}[thm],
\FormulaRefAuto{RationalMulInverse}[thm],
\FormulaRefAuto{RationalDistributive}[thm],
\FormulaRefAuto{RationalOrderTotal}[thm],
\FormulaRefAuto{RationalOrderTranslation}[thm],
\FormulaRefAuto{RationalNonnegativeProduct}[thm],
\FormulaRefAuto{RationalPositiveMulStrictOrder}[thm],
\FormulaRefAuto{RationalPositiveReciprocal}[thm],
\FormulaRefAuto{IntegerRationalEmbeddingFunction}[thm],
\FormulaRefAuto{IntegerRationalEmbeddingInjective}[thm],
\FormulaRefAuto{IntegerRationalEmbeddingAdditive}[thm],
\FormulaRefAuto{IntegerRationalEmbeddingMultiplicative}[thm],
\FormulaRefAuto{IntegerRationalEmbeddingOrder}[thm] und
\FormulaRefAuto{RationalFractionGenerationModel}[thm] und
\FormulaRefAuto{RationalIntegerFractionBoundModel}[thm]; die
Nennercharakterisierung ist die Definition
\FormulaRefAuto{PositiveIntegersDef}[def].\par
\end{sloppypar}

\chapter{Folgesätze aus den rationalen Axiomen}

\section{Dichte und archimedische Eigenschaft}

\FormulaThmDeltaK[Ordnungsumkehr positiver Kehrwerte]{%
  q,r\in\mathbb Q\dsep\RatZero<q\dsep\RatZero<r
  \vdash\bigl(q^{-1}<r^{-1}\leftrightarrow r<q\bigr)%
}{RationalPositiveReciprocalOrderReversal}{%
  \DeltaRow{Mengen}{q\dsep r}%
  \DeltaRow{Relationsoperatoren}{<}%
  \DeltaRow{Funktionsoperatoren}{(\,\cdot\,)^{-1}}%
}
\begin{tabproofsplitwide}
  \proofpartwideindR[Hinrichtung]{%
    q,r\in\mathbb Q\dsep\RatZero<q\dsep\RatZero<r\dsep
    q^{-1}<r^{-1}\vdash r<q
  }{%
    q,r\in\mathbb Q\dsep\RatZero<q\dsep\RatZero<r\dsep
    q^{-1}<r^{-1}\vdash r<q
  }[RationalPositiveReciprocalOrderForward]
    \proofstepwidestarbreakkeep[1]{q,r\in\mathbb Q}{\rA}
    \proofstepwidestarbreakkeep[2]{\RatZero<q}{\rA}
    \proofstepwidestarbreakkeep[3]{\RatZero<r}{\rA}
    \proofstepwidestarbreakkeep[4]{q^{-1}<r^{-1}}{\rA}
    \proofstepwidestarbreakkeep[2,3]{q\neq\RatZero\land r\neq\RatZero}{%
      \rAIRepeat{2}{%
        \FormulaRefAuto{a\neq b\vdash b\neq a}[thm]\circ
        \FormulaRefAuto{RationalStrictOrderDef}[def]}{2\mid3}}
    \proofstepwidestarbreakkeep[1,2,3]{q^{-1},r^{-1}\in\mathbb Q}{%
      \rAIRepeat{2}{%
        \FormulaRefAuto{RationalReciprocalClosureAxiom}[ax]}{%
        \rAEa{1},\rAEa{5}\mid\rAEb{1},\rAEb{5}}}
    \proofstepwidestarbreakkeep[1,2,3,4]{%
      q\cdot q^{-1}<q\cdot r^{-1}}{%
      \FormulaRefAuto{RationalPositiveMulStrictOrderAxiom}[ax]{%
        1,2,4,6}}
    \proofstepwidestar[1,2,3,4]{q \cdot q ^ { - 1 } = \RatOne}{%
      \FormulaRefAuto{RationalMulInverseAxiom}[ax]{%
        \rAEa{1},\rAEa{5}}}
    \proofstepwidestarbreakkeep[1,2,3,4]{%
      \RatOne<q\cdot r^{-1}}{%
      \rIE{8,7}}
    \proofstepwidestar[]{\RatOne\in\mathbb Q}{%
      \FormulaRefAuto{RationalOneCarrierAxiom}[ax]}
    \proofstepwidestar[1,2,3]{q\cdot r^{-1}\in\mathbb Q}{%
      \FormulaRefAuto{RationalMulClosureAxiom}[ax]{1,6}}
    \proofstepwidestarbreakkeep[1,2,3,4]{%
      r\cdot\RatOne<r\cdot(q\cdot r^{-1})}{%
      \FormulaRefAuto{RationalPositiveMulStrictOrderAxiom}[ax]{%
        1,3,9,10,%
        11}}
    \proofstepwidestarbreakkeep[1]{r\cdot\RatOne=r}{%
      \FormulaRefAuto{RationalMulOneAxiom}[ax]{\rAEb{1}}}
    \proofstepwidestarbreakkeep[1,2,3]{%
      r\cdot(q\cdot r^{-1})=(r\cdot r^{-1})\cdot q}{%
      \FormulaRefAuto{RationalMulAssociativeAxiom}[ax]{1,6},%
      \FormulaRefAuto{RationalMulCommutativeAxiom}[ax]{1,6}}
    \proofstepwidestarbreakkeep[1,2,3]{r\cdot r^{-1}=\RatOne}{%
      \FormulaRefAuto{RationalMulInverseAxiom}[ax]{%
        \rAEb{1},\rAEb{5}}}
    \proofstepwidestar[1]{q\cdot\RatOne=q}{%
      \FormulaRefAuto{RationalMulOneAxiom}[ax]{\rAEa{1}}}
    \proofstepwidestar[1]{\RatOne\cdot q=q\cdot\RatOne}{%
      \FormulaRefAuto{RationalMulCommutativeAxiom}[ax]{1}}
    \proofstepwidestarbreakkeep[1,2,3]{(r\cdot r^{-1})\cdot q=q}{%
      \rIE{15,16,%
        17}}
    \proofstepwidestarbreakkeep[1,2,3,4]{r<q}{\rIE{12,13,14,18}}
    \proofstepwidestarbreakkeep[1,2,3]{q^{-1}<r^{-1}\to r<q}{%
      \rRI{4,19}}
  \closeproofpartwideindR

  \proofpartwideindR[R\"uckrichtung]{%
    q,r\in\mathbb Q\dsep\RatZero<q\dsep\RatZero<r\dsep
    r<q\vdash q^{-1}<r^{-1}
  }{%
    q,r\in\mathbb Q\dsep\RatZero<q\dsep\RatZero<r\dsep
    r<q\vdash q^{-1}<r^{-1}
  }[RationalPositiveReciprocalOrderBackward]
    \proofstepwidestarbreakkeep[1]{q,r\in\mathbb Q}{\rA}
    \proofstepwidestarbreakkeep[2]{\RatZero<q}{\rA}
    \proofstepwidestarbreakkeep[3]{\RatZero<r}{\rA}
    \proofstepwidestarbreakkeep[4]{r<q}{\rA}
    \proofstepwidestarbreakkeep[2,3]{q\neq\RatZero\land r\neq\RatZero}{%
      \rAIRepeat{2}{%
        \FormulaRefAuto{a\neq b\vdash b\neq a}[thm]\circ
        \FormulaRefAuto{RationalStrictOrderDef}[def]}{2\mid3}}
    \proofstepwidestarbreakkeep[1,2,3]{q^{-1},r^{-1}\in\mathbb Q}{%
      \rAIRepeat{2}{%
        \FormulaRefAuto{RationalReciprocalClosureAxiom}[ax]}{%
        \rAEa{1},\rAEa{5}\mid\rAEb{1},\rAEb{5}}}
    \proofstepwidestarbreakkeep[1,2,3]{%
      \RatZero<q^{-1}\land\RatZero<r^{-1}}{%
      \rAIRepeat{2}{%
        \FormulaRefAuto{RationalPositiveReciprocalAxiom}[ax]}{%
        \rAEa{1},2\mid\rAEb{1},3}}
    \proofstepwidestarbreakkeep[1,2,3,4]{%
      r^{-1}\cdot r<r^{-1}\cdot q}{%
      \FormulaRefAuto{RationalPositiveMulStrictOrderAxiom}[ax]{%
        1,4,\rAEb{7},6}}
    \proofstepwidestar[1,2,3,4]{r \cdot r ^ { - 1 } = \RatOne}{%
      \FormulaRefAuto{RationalMulInverseAxiom}[ax]{%
          \rAEb{1},\rAEb{5}}}
    \proofstepwidestar[1,2,3]{r^{-1}\cdot r=r\cdot r^{-1}}{%
      \FormulaRefAuto{RationalMulCommutativeAxiom}[ax]{1,6}}
    \proofstepwidestarbreakkeep[1,2,3,4]{%
      \RatOne<r^{-1}\cdot q}{%
      \rIE{%
        10,%
        9,8}}
    \proofstepwidestarbreakkeep[1,2,3,4]{%
      q^{-1}\cdot\RatOne<q^{-1}\cdot(r^{-1}\cdot q)}{%
      \FormulaRefAuto{RationalPositiveMulStrictOrderAxiom}[ax]{%
        1,\rAEa{7},11,6}}
    \proofstepwidestarbreakkeep[1,2,3]{q^{-1}\cdot\RatOne=q^{-1}}{%
      \FormulaRefAuto{RationalMulOneAxiom}[ax]{\rAEa{6}}}
    \proofstepwidestarbreakkeep[1,2,3]{%
      q^{-1}\cdot(r^{-1}\cdot q)=(q^{-1}\cdot q)\cdot r^{-1}}{%
      \FormulaRefAuto{RationalMulAssociativeAxiom}[ax]{1,6},%
      \FormulaRefAuto{RationalMulCommutativeAxiom}[ax]{1,6}}
    \proofstepwidestarbreakkeep[1,2]{q^{-1}\cdot q=\RatOne}{%
      \FormulaRefAuto{RationalMulCommutativeAxiom}[ax]{1,6},%
      \FormulaRefAuto{RationalMulInverseAxiom}[ax]{\rAEa{1},\rAEa{5}}}
    \proofstepwidestar[1,2,3]{r^{-1}\cdot\RatOne=r^{-1}}{%
      \FormulaRefAuto{RationalMulOneAxiom}[ax]{\rAEb{6}}}
    \proofstepwidestar[1]{\RatOne\cdot r^{-1}=r^{-1}\cdot\RatOne}{%
      \FormulaRefAuto{RationalMulCommutativeAxiom}[ax]{1}}
    \proofstepwidestarbreakkeep[1,2,3]{%
      (q^{-1}\cdot q)\cdot r^{-1}=r^{-1}}{%
      \rIE{15,16,%
        17}}
    \proofstepwidestarbreakkeep[1,2,3,4]{q^{-1}<r^{-1}}{%
      \rIE{12,13,14,18}}
    \proofstepwidestarbreakkeep[1,2,3]{r<q\to q^{-1}<r^{-1}}{%
      \rRI{4,19}}
  \closeproofpartwideindR

  \proofpartwide{Zusammenfassung}
    \proofstepwidestarbreakkeep[1]{q,r\in\mathbb Q}{\rA}
    \proofstepwidestarbreakkeep[2]{\RatZero<q}{\rA}
    \proofstepwidestarbreakkeep[3]{\RatZero<r}{\rA}
    \proofstepwidestarbreakkeep[1,2,3]{q^{-1}<r^{-1}\to r<q}{%
      \FormulaRefAuto{RationalPositiveReciprocalOrderForward}[thm]{1,2,3}}
    \proofstepwidestarbreakkeep[1,2,3]{r<q\to q^{-1}<r^{-1}}{%
      \FormulaRefAuto{RationalPositiveReciprocalOrderBackward}[thm]{1,2,3}}
    \proofstepwidestarbreakkeep[1,2,3]{%
      q^{-1}<r^{-1}\leftrightarrow r<q}{\rLRI{4,5}}
  \closeproofpartwide
\end{tabproofsplitwide}

\FormulaThmDeltaK[Dichte der rationalen Ordnung]{%
  q,r\in\mathbb Q\dsep q<r
  \vdash\exists s\in\mathbb Q\,(q<s\land s<r)%
}{RationalOrderDense}{%
  \DeltaRow{Mengen}{q\dsep r\dsep s}%
}
\begin{tabproofsplitwide}
\proofpartwide{Rationalität des Zwischenwerts}
\proofstepwidestarbreakkeep[1]{q,r\in\mathbb Q}{\rA}
  \proofstepwidestarbreakkeep[2]{q<r}{\rA}
  \proofstepwidestarbreakkeep[]{\RatZero<\RatOne}{%
    \FormulaRefAuto{RationalOnePositiveAxiom}[ax]}
  \proofstepwidestar[]{\RatOne\in\mathbb Q}{%
      \FormulaRefAuto{RationalOneCarrierAxiom}[ax]}
  \proofstepwidestar[]{\RatOne\in\mathbb Q}{%
      \FormulaRefAuto{RationalOneCarrierAxiom}[ax]}
  \proofstepwidestar[]{\RatZero\in\mathbb Q}{%
      \FormulaRefAuto{RationalZeroCarrierAxiom}[ax]}
  \proofstepwidestarbreakkeep[]{\RatOne<\RatOne+\RatOne}{%
    \FormulaRefAuto{RationalOrderTranslationAxiom}[ax]{%
      6,%
      4,3},%
    \FormulaRefAuto{RationalAddZeroAxiom}[ax]{%
      5},%
    \FormulaRefAuto{RationalStrictOrderDef}[def]}
  \proofstepwidestarbreakkeep[]{\RatZero<\RatOne+\RatOne}{%
    \FormulaRefAuto{RationalOrderTotalAxiom}[ax]{3,7},%
    \FormulaRefAuto{RationalStrictOrderDef}[def]}
  \proofstepwidestar[]{\RatOne \in \mathbb Q}{%
      \FormulaRefAuto{RationalOneCarrierAxiom}[ax]}
  \proofstepwidestarbreakkeep[]{\RatOne+\RatOne\in\mathbb Q}{%
    \FormulaRefAuto{RationalAddClosureAxiom}[ax]{%
      9}}
  \proofstepwidestar[]{\RatZero\leq\RatOne+\RatOne\land\RatZero\neq\RatOne+\RatOne}{%
      \FormulaRefAuto{RationalStrictOrderDef}[def]{8}}
  \proofstepwidestarbreakkeep[]{\RatZero\neq\RatOne+\RatOne}{%
    \rAEb{11}}
  \proofstepwidestarbreakkeep[]{\RatOne+\RatOne\neq\RatZero}{%
    \FormulaRefAuto{a\neq b\vdash b\neq a}[thm]{12}}
  \proofstepwidestarbreakkeep[]{(\RatOne+\RatOne)^{-1}\in\mathbb Q}{%
    \FormulaRefAuto{RationalReciprocalClosureAxiom}[ax]{%
      10,13}}
  \proofstepwidestarbreakkeep[]{\RatZero
    <(\RatOne+\RatOne)^{-1}}{%
    \FormulaRefAuto{RationalPositiveReciprocalAxiom}[ax]{10,8}}
  \proofstepwidestar[1]{q+r\in\mathbb Q}{%
      \FormulaRefAuto{RationalAddClosureAxiom}[ax]{1}}
  \proofstepwidestarbreakkeep[1]{(q+r)\cdot
    (\RatOne+\RatOne)^{-1}\in\mathbb Q}{%
    \FormulaRefAuto{RationalMulClosureAxiom}[ax]{%
      16,14}}

\closeproofpartwide

\proofpartwide{Untere strikte Schranke}
  \setcounter{proofstepnr}{17}% Fortlaufende Schritte dieses Beweises.
\proofstepwidestarbreakkeep[1,2]{q+q<q+r}{%
    \FormulaRefAuto{RationalOrderTranslationAxiom}[ax]{1,2},%
    \FormulaRefAuto{RationalStrictOrderDef}[def]}
  \proofstepwidestarbreakkeep[1,2]{%
    (\RatOne+\RatOne)^{-1}\cdot(q+q)
    <(\RatOne+\RatOne)^{-1}\cdot(q+r)}{%
    \FormulaRefAuto{RationalPositiveMulStrictOrderAxiom}[ax]{%
      1,15,18,14}}
  \proofstepwidestarbreakkeep[1]{%
    (\RatOne+\RatOne)^{-1}\cdot(q+q)
    =\bigl((\RatOne+\RatOne)^{-1}\cdot
      (\RatOne+\RatOne)\bigr)\cdot q}{%
    \FormulaRefAuto{RationalDistributiveAxiom}[ax]{1,10,14},%
    \FormulaRefAuto{RationalMulAssociativeAxiom}[ax]{1,10,14},%
    \FormulaRefAuto{RationalMulCommutativeAxiom}[ax]{1,10,14}}
  \proofstepwidestarbreakkeep[]{%
    (\RatOne+\RatOne)^{-1}\cdot(\RatOne+\RatOne)=\RatOne}{%
    \FormulaRefAuto{RationalMulCommutativeAxiom}[ax]{10,14},%
    \FormulaRefAuto{RationalMulInverseAxiom}[ax]{10,13}}
  \proofstepwidestar[1]{q\cdot\RatOne=q}{%
      \FormulaRefAuto{RationalMulOneAxiom}[ax]{\rAEa{1}}}
  \proofstepwidestar[1]{\RatOne\cdot q=q\cdot\RatOne}{%
      \FormulaRefAuto{RationalMulCommutativeAxiom}[ax]{1}}
  \proofstepwidestarbreakkeep[1]{%
    \bigl((\RatOne+\RatOne)^{-1}\cdot
      (\RatOne+\RatOne)\bigr)\cdot q=q}{%
    \rIE{21,22,%
      23}}
  \proofstepwidestarbreakkeep[1]{%
    (\RatOne+\RatOne)^{-1}\cdot(q+r)
    =(q+r)\cdot(\RatOne+\RatOne)^{-1}}{%
    \FormulaRefAuto{RationalMulCommutativeAxiom}[ax]{1,14}}
  \proofstepwidestarbreakkeep[1,2]{q<(q+r)\cdot
    (\RatOne+\RatOne)^{-1}}{\rIE{19,20,24,25}}

\closeproofpartwide

\proofpartwide{Obere strikte Schranke}
  \setcounter{proofstepnr}{26}% Fortlaufende Schritte dieses Beweises.
\proofstepwidestarbreakkeep[1,2]{q+r<r+r}{%
    \FormulaRefAuto{RationalOrderTranslationAxiom}[ax]{1,2},%
    \FormulaRefAuto{RationalAddCommutativeAxiom}[ax]{1},%
    \FormulaRefAuto{RationalStrictOrderDef}[def]}
  \proofstepwidestarbreakkeep[1,2]{%
    (\RatOne+\RatOne)^{-1}\cdot(q+r)
    <(\RatOne+\RatOne)^{-1}\cdot(r+r)}{%
    \FormulaRefAuto{RationalPositiveMulStrictOrderAxiom}[ax]{%
      1,15,27,14}}
  \proofstepwidestarbreakkeep[1]{%
    (\RatOne+\RatOne)^{-1}\cdot(r+r)
    =\bigl((\RatOne+\RatOne)^{-1}\cdot
      (\RatOne+\RatOne)\bigr)\cdot r}{%
    \FormulaRefAuto{RationalDistributiveAxiom}[ax]{1,10,14},%
    \FormulaRefAuto{RationalMulAssociativeAxiom}[ax]{1,10,14},%
    \FormulaRefAuto{RationalMulCommutativeAxiom}[ax]{1,10,14}}
  \proofstepwidestar[1]{r\cdot\RatOne=r}{%
      \FormulaRefAuto{RationalMulOneAxiom}[ax]{\rAEb{1}}}
  \proofstepwidestar[1]{\RatOne\cdot r=r\cdot\RatOne}{%
      \FormulaRefAuto{RationalMulCommutativeAxiom}[ax]{1}}
  \proofstepwidestarbreakkeep[1]{%
    \bigl((\RatOne+\RatOne)^{-1}\cdot
      (\RatOne+\RatOne)\bigr)\cdot r=r}{%
    \rIE{21,30,%
      31}}
  \proofstepwidestarbreakkeep[1,2]{(q+r)\cdot
    (\RatOne+\RatOne)^{-1}<r}{\rIE{25,28,29,32}}

\closeproofpartwide

\proofpartwide{Existenz des Zwischenwerts}
  \setcounter{proofstepnr}{33}% Fortlaufende Schritte dieses Beweises.
\proofstepwidestarbreakkeep[1,2]{\exists s\in\mathbb Q\,(q<s\land s<r)}{%
    \rEI{\rAI{17,\rAI{26,33}}}}

\closeproofpartwide
\end{tabproofsplitwide}

\FormulaThmDeltaK[Archimedische Eigenschaft der rationalen Zahlen]{%
  q\in\mathbb Q\vdash
  \exists n\in\mathbb N\,
  q<\RatEmbed(\IntEmbed(n))%
}{RationalArchimedean}{%
  \DeltaRow{Mengen}{q\dsep n}%
  \DeltaRow{Funktionensymbole}{\IntEmbed\dsep\RatEmbed}%
}
\begin{tabproofwide}
  \proofstepwidestar[1]{q\in\mathbb Q}{\rA}
  \proofstepwidestar[1]{\exists a\in\mathbb Z\,\exists b\in\PositiveIntegers\,q=\RatEmbed(a)\cdot(\RatEmbed(b))^{-1}}{\FormulaRefAuto{RationalFractionGenerationAxiom}[ax]{1}}
  \proofstepwidestar[3]{a\in\mathbb Z\land(b\in\PositiveIntegers\land q=\RatEmbed(a)\cdot(\RatEmbed(b))^{-1})}{\rA}
  \proofstepwidestar[3]{a\in\mathbb Z}{\rAEa{3}}
  \proofstepwidestar[3]{b\in\PositiveIntegers}{\rAEa{\rAEb{3}}}
  \proofstepwidestar[3]{b\in\mathbb Z\land\IntZero<b}{\FormulaRefAuto{RationalPositiveDenominatorAxiom}[ax]{5}}
  \proofstepwidestar[3]{b\in\mathbb Z}{\rAEa{6}}
  \proofstepwidestar[]{\RatEmbed\colon\mathbb Z\to\mathbb Q}{\FormulaRefAuto{RationalIntegerEmbeddingFunctionAxiom}[ax]}
  \proofstepwidestar[3]{\RatEmbed(a)\in\mathbb Q}{\FormulaRefAuto{F\colon A\to B\dsep x\in A\vdash F(x)\in B}[thm]{8,4}}
  \proofstepwidestar[3]{\RatEmbed(b)\in\mathbb Q}{\FormulaRefAuto{F\colon A\to B\dsep x\in A\vdash F(x)\in B}[thm]{8,7}}
  \proofstepwidestar[3]{\RatZero<\RatEmbed(b)}{\FormulaRefAuto{RationalIntegerEmbeddingOrderAxiom}[ax]{7,\rAEb{6}},\FormulaRefAuto{RationalIntegerEmbeddingInjectiveAxiom}[ax],\FormulaRefAuto{RationalZeroDef}[def],\FormulaRefAuto{RationalStrictOrderDef}[def]}
  \proofstepwidestar[3]{\RatZero\leq\RatEmbed(b)\land\RatZero\neq\RatEmbed(b)}{\FormulaRefAuto{RationalStrictOrderDef}[def]{11}}
  \proofstepwidestar[3]{\RatZero\neq\RatEmbed(b)}{\rAEb{12}}
  \proofstepwidestar[3]{\RatEmbed(b)\neq\RatZero}{\FormulaRefAuto{a\neq b\vdash b\neq a}[thm]{13}}
  \proofstepwidestar[3]{(\RatEmbed(b))^{-1}\in\mathbb Q}{\FormulaRefAuto{RationalReciprocalClosureAxiom}[ax]{10,14}}
  \proofstepwidestar[3]{\RatZero<(\RatEmbed(b))^{-1}}{\FormulaRefAuto{RationalPositiveReciprocalAxiom}[ax]{10,11}}
  \proofstepwidestar[3]{\exists n\in\mathbb N\,(a<\IntEmbed(n)\cdot b)}{\FormulaRefAuto{RationalIntegerFractionBoundAxiom}[ax]{4,5}}
  \proofstepwidestar[18]{n\in\mathbb N\land a<\IntEmbed(n)\cdot b}{\rA}
  \proofstepwidestar[18]{n\in\mathbb N}{\rAEa{18}}
  \proofstepwidestar[]{\IntEmbed\colon\mathbb N\to\mathbb Z}{\FormulaRefAuto{NaturalIntegerEmbeddingFunction}[thm]}
  \proofstepwidestar[18]{\IntEmbed(n)\in\mathbb Z}{\FormulaRefAuto{F\colon A\to B\dsep x\in A\vdash F(x)\in B}[thm]{20,19}}
  \proofstepwidestar[3,18]{\IntEmbed(n)\cdot b\in\mathbb Z}{\FormulaRefAuto{IntegerMulClosureAxiom}[ax]{21,7}}
  \proofstepwidestar[18]{\RatEmbed(\IntEmbed(n))\in\mathbb Q}{\FormulaRefAuto{F\colon A\to B\dsep x\in A\vdash F(x)\in B}[thm]{8,21}}
  \proofstepwidestar[3,18]{\RatEmbed(\IntEmbed(n)\cdot b)\in\mathbb Q}{\FormulaRefAuto{F\colon A\to B\dsep x\in A\vdash F(x)\in B}[thm]{8,22}}
  \proofstepwidestar[3,18]{\RatEmbed(a)<\RatEmbed(\IntEmbed(n)\cdot b)}{\FormulaRefAuto{RationalIntegerEmbeddingOrderAxiom}[ax]{4,22,\rAEb{18}},\FormulaRefAuto{RationalIntegerEmbeddingInjectiveAxiom}[ax],\FormulaRefAuto{RationalStrictOrderDef}[def]}
  \proofstepwidestar[3,18]{\RatEmbed(\IntEmbed(n)\cdot b)=\RatEmbed(\IntEmbed(n))\cdot\RatEmbed(b)}{\FormulaRefAuto{RationalIntegerEmbeddingMultiplicativeAxiom}[ax]{21,7}}
  \proofstepwidestar[3,18]{(\RatEmbed(b))^{-1}\cdot\RatEmbed(a)<(\RatEmbed(b))^{-1}\cdot\RatEmbed(\IntEmbed(n)\cdot b)}{\FormulaRefAuto{RationalPositiveMulStrictOrderAxiom}[ax]{\rAI{15,\rAI{9,24}},25,16}}
  \proofstepwidestar[3]{(\RatEmbed(b))^{-1}\cdot\RatEmbed(a)=\RatEmbed(a)\cdot(\RatEmbed(b))^{-1}}{\FormulaRefAuto{RationalMulCommutativeAxiom}[ax]{15,9}}
  \proofstepwidestar[3,18]{(\RatEmbed(b))^{-1}\cdot\RatEmbed(\IntEmbed(n)\cdot b)=\RatEmbed(\IntEmbed(n)\cdot b)\cdot(\RatEmbed(b))^{-1}}{\FormulaRefAuto{RationalMulCommutativeAxiom}[ax]{15,24}}
  \proofstepwidestar[3,18]{\RatEmbed(a)\cdot(\RatEmbed(b))^{-1}<\RatEmbed(\IntEmbed(n)\cdot b)\cdot(\RatEmbed(b))^{-1}}{\rIE{28,29,27}}
  \proofstepwidestar[3]{\RatEmbed(b)\cdot(\RatEmbed(b))^{-1}=\RatOne}{\FormulaRefAuto{RationalMulInverseAxiom}[ax]{10,14}}
  \proofstepwidestar[3,18]{(\RatEmbed(\IntEmbed(n))\cdot\RatEmbed(b))\cdot(\RatEmbed(b))^{-1}=\RatEmbed(\IntEmbed(n))\cdot(\RatEmbed(b)\cdot(\RatEmbed(b))^{-1})}{\FormulaRefAuto{RationalMulAssociativeAxiom}[ax]{23,10,15}}
  \proofstepwidestar[18]{\RatEmbed(\IntEmbed(n))\cdot\RatOne=\RatEmbed(\IntEmbed(n))}{\FormulaRefAuto{RationalMulOneAxiom}[ax]{23}}
  \proofstepwidestar[3,18]{\RatEmbed(\IntEmbed(n)\cdot b)\cdot(\RatEmbed(b))^{-1}=\RatEmbed(\IntEmbed(n))}{\rIE{26,31,32,33}}
  \proofstepwidestar[3,18]{q<\RatEmbed(\IntEmbed(n))}{\rIE{\rAEb{\rAEb{3}},30,34}}
  \proofstepwidestar[3,18]{\exists n\in\mathbb N\,(q<\RatEmbed(\IntEmbed(n)))}{\rEI{\rAI{19,35}}}
  \proofstepwidestar[3]{\exists n\in\mathbb N\,(q<\RatEmbed(\IntEmbed(n)))}{\rEE{17,18,36}}
  \proofstepwidestar[1]{\exists n\in\mathbb N\,(q<\RatEmbed(\IntEmbed(n)))}{\rEE{2,3,37}}
\end{tabproofwide}

\end{document}
